%% --- BEGIN page-0106.tex ---
The fundamental group $\pi_1(X)$ is especially useful when studying spaces of low
dimension, as one would expect from its definition which involves only maps from
low-dimensional spaces into $X$, namely loops $I \to X$ and homotopies of loops, maps
$I \times I \to X$. The definition in terms of objects that are at most 2-dimensional manifests
itself for example in the fact that when $X$ is a CW complex, $\pi_1(X)$ depends only on
the 2-skeleton of $X$. In view of the low-dimensional nature of the fundamental group,
we should not expect it to be a very refined tool for dealing with high-dimensional
spaces. Thus it cannot distinguish between spheres $S^n$ with $n \geq 2$. This limitation
to low dimensions can be removed by considering the natural higher-dimensional
analogs of $\pi_1(X)$, the homotopy groups $\pi_n(X)$, which are defined in terms of maps
of the $n$-dimensional cube $I^n$ into $X$ and homotopies $I^n \times I \to X$ of such maps. Not
surprisingly, when $X$ is a CW complex, $\pi_n(X)$ depends only on the $(n+1)$-skeleton
of $X$. And as one might hope, homotopy groups do indeed distinguish spheres of all
dimensions since $\pi_i(S^n)$ is 0 for $i < n$ and $\mathbb{Z}$ for $i = n$.

However, the higher-dimensional homotopy groups have the serious drawback
that they are extremely difficult to compute in general. Even for simple spaces like
spheres, the calculation of $\pi_i(S^n)$ for $i > n$ turns out to be a huge problem. For-
tunately there is a more computable alternative to homotopy groups: the homology
groups $H_n(X)$. Like $\pi_n(X)$, the homology group $H_n(X)$ for a CW complex $X$ de-
pends only on the $(n+1)$-skeleton. For spheres, the homology groups $H_i(S^n)$ are
isomorphic to the homotopy groups $\pi_i(S^n)$ in the range $1 \leq i \leq n$, but homology
groups have the advantage that $H_i(S^n) = 0$ for $i > n$.

The computability of homology groups does not come for free, unfortunately.
The definition of homology groups is decidedly less transparent than the definition
of homotopy groups, and once one gets beyond the definition there is a certain amount
of technical machinery to be set up before any real calculations and applications can
be given. In the exposition below we approach the definition of $H_n(X)$ by two prelim-
inary stages, first giving a few motivating examples nonrigorously, then constructing
%% --- END page-0106.tex ---
%% --- BEGIN page-0108.tex ---
Let us look at some examples to see what the idea is. Consider the graph $X_1$ shown in the figure, consisting of two vertices joined by four edges.



When studying the fundamental group of $X_1$, we consider loops formed by sequences of edges, starting and ending at a fixed basepoint. For example, at the basepoint $x$, the loop $ab^{-1}$ travels forward along the edge $a$, then backward along $b$, as indicated by the exponent $-1$. A more complicated loop would be $ac^{-1}bd^{-1}ca^{-1}$. A salient feature of the fundamental group is that it is generally nonabelian, which both enriches and complicates the theory. Suppose we simplify matters by abelianizing. Thus for example the two loops $ab^{-1}$ and $b^{-1}a$ are to be regarded as equal if we make $a$ commute with $b^{-1}$. These two loops $ab^{-1}$ and $b^{-1}a$ are really the same circle, just with a different choice of starting and ending point: $x$ for $ab^{-1}$ and $y$ for $b^{-1}a$. The same thing happens for all loops: Rechoosing the basepoint in a loop just permutes its letters cyclically, so a byproduct of abelianizing is that we no longer have to pin all our loops down to a fixed basepoint. Thus loops become \emph{cycles}, without a chosen basepoint.

Having abelianized, let us switch to additive notation, so cycles become linear combinations of edges with integer coefficients, such as $a - b + c - d$. Let us call these linear combinations \emph{chains} of edges. Some chains can be decomposed into cycles in several different ways, for example $(a - c) + (b - d) = (a - d) + (b - c)$, and if we adopt an algebraic viewpoint then we do not want to distinguish between these different decompositions. Thus we broaden the meaning of the term `cycle' to be simply any linear combination of edges for which at least one decomposition into cycles in the previous more geometric sense exists.

What is the condition for a chain to be a cycle in this more algebraic sense? A geometric cycle, thought of as a path traversed in time, is distinguished by the property that it enters each vertex the same number of times that it leaves the vertex. For an arbitrary chain $ka + \ell b + mc + nd$, the net number of times this chain enters $y$ is $k + \ell + m + n$ since each of $a, b, c,$ and $d$ enters $y$ once. Similarly, each of the four edges leaves $x$ once, so the net number of times the chain $ka + \ell b + mc + nd$ enters $x$ is $-k - \ell - m - n$. Thus the condition for $ka + \ell b + mc + nd$ to be a cycle is simply $k + \ell + m + n = 0$.

To describe this result in a way that would generalize to all graphs, let $C_1$ be the free abelian group with basis the edges $a, b, c, d$ and let $C_0$ be the free abelian group with basis the vertices $x, y$. Elements of $C_1$ are chains of edges, or 1-dimensional chains, and elements of $C_0$ are linear combinations of vertices, or 0-dimensional chains. Define a homomorphism $\partial: C_1 \to C_0$ by sending each basis element $a, b, c, d$ to $y - x$, the vertex at the head of the edge minus the vertex at the tail. Thus we have $\partial(ka + \ell b + mc + nd) = (k + \ell + m + n)y - (k + \ell + m + n)x$, and the cycles are precisely the kernel of $\partial$. It is a simple calculation to verify that $a - b, b - c,$ and $c - d$

%% --- END page-0108.tex ---
%% --- BEGIN page-0109.tex ---
form a basis for this kernel. Thus every cycle in $X_1$ is a unique linear combination of
these three most obvious cycles. By means of these three basic cycles we convey the
geometric information that the graph $X_1$ has three visible `holes', the empty spaces
between the four edges.

Let us now enlarge the preceding graph $X_1$ by attaching a 2-cell $A$ along the
cycle $a - b$, producing a 2-dimensional cell complex $X_2$. If

we think of the 2-cell $A$ as being oriented clockwise, then
we can regard its boundary as the cycle $a - b$. This cycle is
now homotopically trivial since we can contract it to a point
by sliding over $A$. In other words, it no longer encloses a
hole in $X_2$. This suggests that we form a quotient of the
group of cycles in the preceding example by factoring out
the subgroup generated by $a - b$. In this quotient the cycles $a - c$ and $b - c$, for
example, become equivalent, consistent with the fact that they are homotopic in $X_2$.

Algebraically, we can define now a pair of homomorphisms $C_2 \xrightarrow{\partial_2} C_1 \xrightarrow{\partial_1} C_0$
where $C_2$ is the infinite cyclic group generated by $A$ and $\partial_2(A) = a - b$. The map
$\partial_1$ is the boundary homomorphism in the previous example. The quotient group we
are interested in is $\ker \partial_1 / \text{im} \partial_2$, the kernel of $\partial_1$ modulo the image of $\partial_2$, or in other
words, the 1-dimensional cycles modulo those that are boundaries, the multiples of
$a - b$. This quotient group is the \textit{homology group} $H_1(X_2)$. The previous example can
be fit into this scheme too by taking $C_2$ to be zero since there are no 2-cells in $X_1$,
so in this case $H_1(X_1) = \ker \partial_1 / \text{im} \partial_2 = \ker \partial_1$, which as we saw was free abelian on
three generators. In the present example, $H_1(X_2)$ is free abelian on two generators,
$b - c$ and $c - d$, expressing the geometric fact that by filling in the 2-cell $A$ we have
reduced the number of `holes' in our space from three to two.

Suppose we enlarge $X_2$ to a space $X_3$ by attaching a second 2-cell $B$ along the
same cycle $a - b$. This gives a 2-dimensional chain group $C_2$

consisting of linear combinations of $A$ and $B$, and the bound-
ary homomorphism $\partial_2: C_2 \to C_1$ sends both $A$ and $B$ to $a-b$.
The homology group $H_1(X_3) = \ker \partial_1 / \text{im} \partial_2$ is the same as
for $X_2$, but now $\partial_2$ has a nontrivial kernel, the infinite cyclic
group generated by $A - B$. We view $A - B$ as a 2-dimensional
cycle, generating the homology group $H_2(X_3) = \ker \partial_2 \cong \mathbb{Z}$.
Topologically, the cycle $A - B$ is the sphere formed by the cells $A$ and $B$ together
with their common boundary circle. This spherical cycle detects the presence of a
`hole' in $X_3$, the missing interior of the sphere. However, since this hole is enclosed
by a sphere rather than a circle, it is of a different sort from the holes detected by
$H_1(X_3) \cong \mathbb{Z} \times \mathbb{Z}$, which are detected by the cycles $b - c$ and $c - d$.

Let us continue one more step and construct a complex $X_4$ from $X_3$ by attaching
a 3-cell $C$ along the 2-sphere formed by $A$ and $B$. This creates a chain group $C_3$

%% --- END page-0109.tex ---
%% --- BEGIN page-0111.tex ---
\section{Simplicial and Singular Homology}

The most important homology theory in algebraic topology, and the one we shall be studying almost exclusively, is called singular homology. Since the technical apparatus of singular homology is somewhat complicated, we will first introduce a more primitive version called simplicial homology in order to see how some of the apparatus works in a simpler setting before beginning the general theory.

The natural domain of definition for simplicial homology is a class of spaces we call $\Delta$-complexes, which are a mild generalization of the more classical notion of a simplicial complex. Historically, the modern definition of singular homology was first given in [Eilenberg 1944], and $\Delta$-complexes were introduced soon thereafter in [Eilenberg-Zilber 1950] where they were called semisimplicial complexes. Within a few years this term came to be applied to what Eilenberg and Zilber called complete semisimplicial complexes, and later there was yet another shift in terminology as the latter objects came to be called simplicial sets. In theory this frees up the term semisimplicial complex to have its original meaning, but to avoid potential confusion it seems best to introduce a new name, and the term $\Delta$-complex has at least the virtue of brevity.

\textbf{$\Delta$-Complexes}

The torus, the projective plane, and the Klein bottle can each be obtained from a square by identifying opposite edges in the way indicated by the arrows in the following figures:



Cutting a square along a diagonal produces two triangles, so each of these surfaces can also be built from two triangles by identifying their edges in pairs. In similar fashion a polygon with any number of sides can be cut along diagonals into triangles, so in fact all closed surfaces can be constructed from triangles by identifying edges. Thus we have a single building block, the triangle, from which all surfaces can be constructed. Using only triangles we could also construct a large class of 2-dimensional spaces that are not surfaces in the strict sense, by allowing more than two edges to be identified together at a time.



%% --- END page-0111.tex ---
%% --- BEGIN page-0112.tex ---
The idea of a $\Delta$-complex is to generalize constructions like these to any number of dimensions. The $n$-dimensional analog of the triangle is the \textbf{$n$-simplex}. This is the smallest convex set in a Euclidean space $\mathbb{R}^m$ containing $n + 1$ points $v_0, \cdots, v_n$ that do not lie in a hyperplane of dimension less than $n$, where by a hyperplane we mean the set of solutions of a system of linear equations. An equivalent condition would be that the difference vectors $v_1 - v_0, \cdots, v_n - v_0$ are linearly independent. The points $v_i$ are the \textbf{vertices} of the simplex, and the simplex itself is denoted $[v_0, \cdots, v_n]$. For example, there is the standard $n$-simplex

\[
\Delta^n = \{(t_0, \cdots, t_n) \in \mathbb{R}^{n+1} \mid \sum_i t_i = 1 \text{ and } t_i \geq 0 \text{ for all } i\}
\]

whose vertices are the unit vectors along the coordinate axes.



For purposes of homology it will be important to keep track of the order of the vertices of a simplex, so `$n$-simplex' will really mean `$n$-simplex with an ordering of its vertices'. A by-product of ordering the vertices of a simplex $[v_0, \cdots, v_n]$ is that this determines orientations of the edges $[v_i, v_j]$ according to increasing subscripts, as shown in the two preceding figures. Specifying the ordering of the vertices also determines a canonical linear homeomorphism from the standard $n$-simplex $\Delta^n$ onto any other $n$-simplex $[v_0, \cdots, v_n]$, preserving the order of vertices, namely, $(t_0, \cdots, t_n) \mapsto \sum_i t_i v_i$. The coefficients $t_i$ are the \textbf{barycentric coordinates} of the point $\sum_i t_i v_i$ in $[v_0, \cdots, v_n]$.

If we delete one of the $n + 1$ vertices of an $n$-simplex $[v_0, \cdots, v_n]$, then the remaining $n$ vertices span an $(n - 1)$-simplex, called a \textbf{face} of $[v_0, \cdots, v_n]$. We adopt the following convention:

\textit{The vertices of a face, or of any subsimplex spanned by a subset of the vertices, will always be ordered according to their order in the larger simplex.}

The union of all the faces of $\Delta^n$ is the \textbf{boundary} of $\Delta^n$, written $\partial \Delta^n$. The \textbf{open simplex} $\dot{\Delta}^n$ is $\Delta^n - \partial \Delta^n$, the interior of $\Delta^n$.

A \textbf{$\Delta$-complex structure} on a space $X$ is a collection of maps $\sigma_\alpha: \Delta^n \to X$, with $n$ depending on the index $\alpha$, such that:

(i) The restriction $\sigma_\alpha|_{\dot{\Delta}^n}$ is injective, and each point of $X$ is in the image of exactly one such restriction $\sigma_\alpha|_{\dot{\Delta}^n}$.

(ii) Each restriction of $\sigma_\alpha$ to a face of $\Delta^n$ is one of the maps $\sigma_\beta: \Delta^{n-1} \to X$. Here we are identifying the face of $\Delta^n$ with $\Delta^{n-1}$ by the canonical linear homeomorphism between them that preserves the ordering of the vertices.

(iii) A set $A \subset X$ is open iff $\sigma_\alpha^{-1}(A)$ is open in $\Delta^n$ for each $\sigma_\alpha$.

%% --- END page-0112.tex ---
%% --- BEGIN page-0113.tex ---
Among other things, this last condition rules out trivialities like regarding all the
points of $X$ as individual vertices. The earlier decompositions of the torus, projective
plane, and Klein bottle into two triangles, three edges, and one or two vertices define
$\Delta$-complex structures with a total of six $\sigma_\alpha$'s for the torus and Klein bottle and seven
for the projective plane. The orientations on the edges in the pictures are compatible
with a unique ordering of the vertices of each simplex, and these orderings determine
the maps $\sigma_\alpha$.

A consequence of (iii) is that $X$ can be built as a quotient space of a collection
of disjoint simplices $\Delta^\alpha$, one for each $\sigma_\alpha:\Delta^n \to X$, the quotient space obtained by
identifying each face of a $\Delta^\alpha$ with the $\Delta^\beta$ corresponding to the restriction $\sigma_\beta$ of
$\sigma_\alpha$ to the face in question, as in condition (ii). One can think of building the quotient
space inductively, starting with a discrete set of vertices, then attaching edges to
these to produce a graph, then attaching 2-simplices to the graph, and so on. From
this viewpoint we see that the data specifying a $\Delta$-complex can be described purely
combinatorially as collections of $n$-simplices $\Delta^\alpha$ for each $n$ together with functions
associating to each face of each $n$-simplex $\Delta^\alpha$ an $(n-1)$-simplex $\Delta^\beta$.

More generally, $\Delta$-complexes can be built from collections of disjoint simplices by
identifying various subsimplices spanned by subsets of the vertices, where the identifications are performed using the canonical linear homeomorphisms that preserve
the orderings of the vertices. The earlier $\Delta$-complex structures on a torus, projective
plane, or Klein bottle can be obtained in this way, by identifying pairs of edges of
two 2-simplices. If one starts with a single 2-simplex and identifies all three edges
to a single edge, preserving the orientations given by the ordering of the vertices,
this produces a $\Delta$-complex known as the `dunce hat'. By contrast, if the three edges
of a 2-simplex are identified preserving a cyclic orientation of the three edges, as in
the first figure at the right, this does not produce a
$\Delta$-complex structure, although if the 2-simplex is
subdivided into three smaller 2-simplices about a
central vertex, then one does obtain a $\Delta$-complex
structure on the quotient space.

Thinking of a $\Delta$-complex $X$ as a quotient space of a collection of disjoint simplices, it is not hard to see that $X$ must be a Hausdorff space. Condition (iii) then
implies that each restriction $\sigma_\alpha|_{\Delta^n}$ is a homeomorphism onto its image, which is
thus an open simplex in $X$. It follows from Proposition A.2 in the Appendix that
these open simplices $\sigma_\alpha(\hat{\Delta}^n)$ are the cells $e^\alpha_n$ of a CW complex structure on $X$ with
the $\sigma_\alpha$'s as characteristic maps. We will not need this fact at present, however.

\section*{Simplicial Homology}

Our goal now is to define the simplicial homology groups of a $\Delta$-complex $X$. Let
$\Delta_n(X)$ be the free abelian group with basis the open $n$-simplices $e^\alpha_n$ of $X$. Elements
%% --- END page-0113.tex ---
%% --- BEGIN page-0114.tex ---
of $\Delta_n(X)$, called \textbf{$n$-chains}, can be written as finite formal sums $\sum_\alpha n_\alpha e^\alpha_n$ with coefficients $n_\alpha \in \mathbb{Z}$. Equivalently, we could write $\sum_\alpha n_\alpha \sigma_\alpha$ where $\sigma_\alpha:\Delta^n \to X$ is the
characteristic map of $e^\alpha_n$, with image the closure of $e^\alpha_n$ as described above. Such a
sum $\sum_\alpha n_\alpha\sigma_\alpha$ can be thought of as a finite collection, or `chain', of $n$-simplices in $X$,
with integer multiplicities, the coefficients $n_\alpha$.

As one can see in the next figure, the boundary of the $n$-simplex $[v_0, \cdots, v_n]$ consists of the various $(n-1)$-dimensional simplices $[v_0, \cdots, \hat{v}_i, \cdots, v_n]$, where the `hat'
symbol $\hat{}$ over $v_i$ indicates that this vertex is deleted from the sequence $v_0, \cdots, v_n$.
In terms of chains, we might then wish to say that the boundary of $[v_0, \cdots, v_n]$ is the
$(n-1)$-chain formed by the sum of the faces $[v_0, \cdots, \hat{v}_i, \cdots, v_n]$. However, it turns
out to be better to insert certain signs and instead let the boundary of $[v_0, \cdots, v_n]$ be
$\sum_i (-1)^i [v_0, \cdots, \hat{v}_i, \cdots, v_n]$. Heuristically, the signs are inserted to take orientations
into account, so that all the faces of a simplex are coherently oriented, as indicated in
the following figure:



In the last case, the orientations of the two hidden faces are also counterclockwise
when viewed from outside the 3-simplex.

With this geometry in mind we define for a general $\Delta$-complex $X$ a \textbf{boundary
homomorphism} $\partial_n : \Delta_n(X) \to \Delta_{n-1}(X)$ by specifying its values on basis elements:

$$\partial_n(\sigma_\alpha) = \sum_i (-1)^i \sigma_\alpha |[v_0, \cdots, \hat{v}_i, \cdots, v_n]$$

Note that the right side of this equation does indeed lie in $\Delta_{n-1}(X)$ since each restriction $\sigma_\alpha |[v_0, \cdots, \hat{v}_i, \cdots, v_n]$ is the characteristic map of an $(n-1)$-simplex of $X$.

\begin{lemma}[Composition of boundary maps is zero]
\label{lem:boundary-composition-zero}
\dcref{ch2:2.1}
\group{Singular Homology}
\uses{}
The composition $\Delta_n(X) \xrightarrow{\partial_n} \Delta_{n-1}(X) \xrightarrow{\partial_{n-1}} \Delta_{n-2}(X)$ is zero.
\end{lemma}
\begin{proof}
 We have $\partial_n(\sigma) = \sum_j (-1)^j \sigma|[v_0, \cdots, \hat{v}_j, \cdots, v_n]$, and hence
$$\partial_{n-1} \partial_n(\sigma) = \sum (-1)^i (-1)^j \sigma|[v_0, \cdots, \hat{v}_i, \cdots, \hat{v}_j, \cdots, v_n]$$
$$+ \sum_{j > i} (-1)^i(-1)^{j-1}\sigma|[v_0, \cdots, \hat{v}_i, \cdots, \hat{v}_j, \cdots, v_n]$$

%% --- END page-0114.tex ---
%% --- BEGIN page-0115.tex ---
The latter two summations cancel since after switching $i$ and $j$ in the second sum, it
becomes the negative of the first.
\end{proof}
The algebraic situation we have now is a sequence of homomorphisms of abelian
groups

$$\cdots \to C_{n+1} \xrightarrow{\partial_{n+1}} C_n \xrightarrow{\partial_n} C_{n-1} \to \cdots \to C_1 \xrightarrow{\partial_1} C_0 \xrightarrow{\partial_0} 0$$

with $\partial_n\partial_{n+1} = 0$ for each $n$. Such a sequence is called a \textbf{chain complex}. Note that
we have extended the sequence by a 0 at the right end, with $\partial_0 = 0$. The equation
$\partial_n\partial_{n+1} = 0$ is equivalent to the inclusion $\text{Im}\partial_{n+1} \subset \ker\partial_n$, where Im and Ker denote
image and kernel. So we can define the $n^{\text{th}}$ \textbf{homology group} of the chain complex to
be the quotient group $H_n = \ker\partial_n / \text{Im}\partial_{n+1}$. Elements of $\ker\partial_n$ are called \textbf{cycles} and
elements of $\text{Im}\partial_{n+1}$ are called \textbf{boundaries}. Elements of $H_n$ are cosets of $\text{Im}\partial_{n+1}$,
called \textbf{homology classes}. Two cycles representing the same homology class are said
to be \textbf{homologous}. This means their difference is a boundary.

Returning to the case that $C_n = \Delta_n(X)$, the homology group $\ker\partial_n / \text{Im}\partial_{n+1}$ will
be denoted $H^\Delta_n(X)$ and called the $n^{\text{th}}$ \textbf{simplicial homology group} of $X$.

\begin{example}[Simplicial homology of $S^1$]
\label{ex:simplicial-homology-circle}
\dcref{ch2:2.2}
\group{Singular Homology}
\uses{lem:boundary-composition-zero}
$X = S^1$, with one vertex $v$ and one edge $e$. Then $\Delta_0(S^1)$
and $\Delta_1(S^1)$ are both $\mathbb{Z}$ and the boundary map $\partial_1$ is zero since $\partial e = v - v$.
The groups $\Delta_n(S^1)$ are 0 for $n \geq 2$ since there are no simplices in these
dimensions. Hence
\end{example}

$$H^\Delta_n(S^1) \cong \begin{cases} \mathbb{Z} & \text{for } n = 0, 1 \\ 0 & \text{for } n \geq 2 \end{cases}$$

This is an illustration of the general fact that if the boundary maps in a chain complex
are all zero, then the homology groups of the complex are isomorphic to the chain
groups themselves.

\begin{example}[Simplicial homology of the torus]
\label{ex:simplicial-homology-torus}
\dcref{ch2:2.3}
\group{Singular Homology}
\uses{lem:boundary-composition-zero}
$X = T$, the torus with the $\Delta$-complex structure pictured earlier, having
one vertex, three edges $a, b,$ and $c$, and two 2-simplices $U$ and $L$. As in the previous
example, $\partial_1 = 0$ so $H^\Delta_0(T) \cong \mathbb{Z}$. Since $\partial_2 U = a + b - c = \partial_2 L$ and $\{a, b, a+b-c\}$ is
a basis for $\Delta_1(T)$, it follows that $H^\Delta_1(T) \cong \mathbb{Z} \oplus \mathbb{Z}$ with basis the homology classes $[a]$
and $[b]$. Since there are no 3-simplices, $H^\Delta_2(T)$ is equal to $\ker\partial_2$, which is infinite
cyclic generated by $U - L$ since $\partial(pU + qL) = (p+q)(a+b-c) = 0$ only if $p = -q$.
Thus
\end{example}

$$H^\Delta_n(T) \cong \begin{cases} \mathbb{Z} \oplus \mathbb{Z} & \text{for } n = 1 \\ \mathbb{Z} & \text{for } n = 0, 2 \\ 0 & \text{for } n \geq 3 \end{cases}$$

\begin{example}[Simplicial homology of $\mathbb{RP}^2$]
\label{ex:simplicial-homology-rp2}
\dcref{ch2:2.4}
\group{Singular Homology}
\uses{lem:boundary-composition-zero}
$X = \mathbb{RP}^2$, as pictured earlier, with two vertices $v$ and $w$, three edges
$a, b,$ and $c$, and two 2-simplices $U$ and $L$. Then $\text{Im}\partial_1$ is generated by $w - v$, so
$H^\Delta_0(X) \cong \mathbb{Z}$ with either vertex as a generator. Since $\partial_2 U = -a + b + c$ and $\partial_2 L = a - b + c$,
we see that $\partial_2$ is injective, so $H^\Delta_2(X) = 0$. Further, $\ker\partial_1 \cong \mathbb{Z} \oplus \mathbb{Z}$ with basis $a - b$ and
$c$, and $\text{Im}\partial_2$ is an index-two subgroup of $\ker\partial_1$ since we can choose $c$ and $a - b + c$ as
\end{example}

%% --- END page-0115.tex ---
%% --- BEGIN page-0116.tex ---
a basis for $\ker\partial_1$ and $a - b + c$ and $2c = (a - b + c) + (-a + b + c)$ as a basis for
$\text{Im}\partial_2$. Thus $H^\Delta_1(X) \cong \mathbb{Z}_2$.

\begin{example}[$\Delta$-complex structure on $S^n$]
\label{ex:delta-complex-sn}
\dcref{ch2:2.5}
\group{Singular Homology}
\uses{lem:boundary-composition-zero}
We can obtain a $\Delta$-complex structure on $S^n$ by taking two copies of $\Delta^n$
and identifying their boundaries via the identity map. Labeling these two $n$-simplices
$U$ and $L$, then it is obvious that $\ker\partial_n$ is infinite cyclic generated by $U - L$. Thus
$H^\Delta_n(S^n) \cong \mathbb{Z}$ for this $\Delta$-complex structure on $S^n$. Computing the other homology
groups would be more difficult.
\end{example}

Many similar examples could be worked out without much trouble, such as the
other closed orientable and nonorientable surfaces. However, the calculations do tend
to increase in complexity before long, particularly for higher-dimensional complexes.

Some obvious general questions arise: Are the groups $H^\Delta_n(X)$ independent of
the choice of $\Delta$-complex structure on $X$? In other words, if two $\Delta$-complexes are
homeomorphic, do they have isomorphic homology groups? More generally, do they
have isomorphic homology groups if they are merely homotopy equivalent? To answer
such questions and to develop a general theory it is best to leave the rather rigid
simplicial realm and introduce the singular homology groups. These have the added
advantage that they are defined for all spaces, not just $\Delta$-complexes. At the end of
this section, after some theory has been developed, we will show that simplicial and
singular homology groups coincide for $\Delta$-complexes.

Traditionally, simplicial homology is defined for \textbf{simplicial complexes}, which are
the $\Delta$-complexes whose simplices are uniquely determined by their vertices. This
amounts to saying that each $n$-simplex has $n + 1$ distinct vertices, and that no other
$n$-simplex has this same set of vertices. Thus a simplicial complex can be described
combinatorially as a set $X_0$ of vertices together with sets $X_n$ of $n$-simplices, which
are $(n+1)$-element subsets of $X_0$. The only requirement is that each $(k+1)$-element
subset of the vertices of an $n$-simplex in $X_n$ is a $k$-simplex, in $X_k$. From this combinatorial data a $\Delta$-complex $X$ can be constructed, once we choose a partial ordering
of the vertices $X_0$ that restricts to a linear ordering on the vertices of each simplex
in $X_n$. For example, we could just choose a linear ordering of all the vertices. This
might perhaps involve invoking the Axiom of Choice for large vertex sets.

An exercise at the end of this section is to show that every $\Delta$-complex can be
subdivided to be a simplicial complex. In particular, every $\Delta$-complex is then homeomorphic to a simplicial complex.

Compared with simplicial complexes, $\Delta$-complexes have the advantage of simpler
computations since fewer simplices are required. For example, to put a simplicial
complex structure on the torus one needs at least 14 triangles, 21 edges, and 7 vertices,
and for $\mathbb{RP}^2$ one needs at least 10 triangles, 15 edges, and 6 vertices. This would slow
down calculations considerably!

%% --- END page-0116.tex ---
%% --- BEGIN page-0117.tex ---
\section*{Singular Homology}

A \textbf{singular $n$-simplex} in a space $X$ is by definition just a map $\sigma:\Delta^n \to X$. The
word `singular' is used here to express the idea that $\sigma$ need not be a nice embedding
but can have `singularities' where its image does not look at all like a simplex. All that
is required is that $\sigma$ be continuous. Let $C_n(X)$ be the free abelian group with basis
the set of singular $n$-simplices in $X$. Elements of $C_n(X)$, called \textbf{$n$-chains}, or more
precisely singular $n$-chains, are finite formal sums $\sum_i n_i \sigma_i$ for $n_i \in \mathbb{Z}$ and $\sigma_i:\Delta^n \to X$.
A boundary map $\partial_n : C_n(X) \to C_{n-1}(X)$ is defined by the same formula as before:

$$\partial_n(\sigma) = \sum_{i=0}^{n} (-1)^i \sigma|[v_0, \cdots, \hat{v}_i, \cdots, v_n]$$

Implicit in this formula is the canonical identification of $[v_0, \cdots, \hat{v}_i, \cdots, v_n]$ with
$\Delta^{n-1}$, preserving the ordering of vertices, so that $\sigma|[v_0, \cdots, \hat{v}_i, \cdots, v_n]$ is regarded
as a map $\Delta^{n-1} \to X$, that is, a singular $(n-1)$-simplex.

Often we write the boundary map $\partial_n$ from $C_n(X)$ to $C_{n-1}(X)$ simply as $\partial$ when
this does not lead to serious ambiguities. The proof of Lemma 2.1 applies equally well
to singular simplices, showing that $\partial_n \partial_{n+1} = 0$ or more concisely $\partial^2 = 0$, so we can
define the \textbf{singular homology group} $H_n(X) = \ker\partial_n / \text{Im}\partial_{n+1}$.

It is evident from the definition that homeomorphic spaces have isomorphic singular
homology groups $H_n$, in contrast with the situation for $H^\Delta_n$. On the other hand,
since the groups $C_n(X)$ are so large, the number of singular $n$-simplices in $X$ usually
being uncountable, it is not at all clear that for a $\Delta$-complex $X$ with finitely many simplices, $H_n(X)$ should be finitely generated for all $n$, or that $H_n(X)$ should be zero
for $n$ larger than the dimension of $X$ -- two properties that are trivial for $H^\Delta_n(X)$.

Though singular homology looks so much more general than simplicial homology,
it can actually be regarded as a special case of simplicial homology by means of the
following construction. For an arbitrary space $X$, define the \textbf{singular complex} $S(X)$
to be the $\Delta$-complex with one $n$-simplex $\Delta^n_\sigma$ for each singular $n$-simplex $\sigma:\Delta^n \to X$,
with $\Delta^n_\sigma$ attached in the obvious way to the $(n-1)$-simplices of $S(X)$ that are the
restrictions of $\sigma$ to the various $(n-1)$-simplices in $\partial\Delta^n$. It is clear from the definitions that $H^\Delta_n(S(X))$ is identical with $H_n(X)$ for all $n$, and in this sense the singular
homology group $H_n(X)$ is a special case of a simplicial homology group. One can
regard $S(X)$ as a $\Delta$-complex model for $X$, although it is usually an extremely large
object compared to $X$.

Cycles in singular homology are defined algebraically, but they can be given a
somewhat more geometric interpretation in terms of maps from finite $\Delta$-complexes.
To see this, note first that a singular $n$-chain $\xi$ can always be written in the form
$\sum_i \epsilon_i \sigma_i$ with $\epsilon_i = \pm 1$, allowing repetitions of the singular $n$-simplices $\sigma_i$. Given such
an $n$-chain $\xi = \sum_i \epsilon_i \sigma_i$, when we compute $\partial \xi$ as a sum of singular $(n-1)$-simplices
with signs $\pm 1$, there may be some \textit{canceling pairs} consisting of two identical singular
$(n-1)$-simplices with opposite signs. Choosing a maximal collection of such

%% --- END page-0117.tex ---
%% --- BEGIN page-0118.tex ---
canceling pairs, construct an $n$-dimensional $\Delta$-complex $K_\xi$ from a disjoint union of
$n$-simplices $\Delta^n_i$, one for each $\sigma_i$, by identifying the pairs of $(n-1)$-dimensional faces
corresponding to the chosen canceling pairs. The $\sigma_i$'s then induce a map $K_\xi \to X$. If
$\xi$ is a cycle, all the $(n-1)$-dimensional faces of the $\Delta^n_i$'s are identified in pairs. Thus
$K_\xi$ is a manifold, locally homeomorphic to $\mathbb{R}^n$, near all points in the complement
of the $(n-2)$-skeleton $K^{n-2}_\xi$ of $K_\xi$. All the $n$-simplices of $K_\xi$ can be coherently
oriented by taking the signs of the $\sigma_i$'s into account, so $K_\xi - K^{n-2}_\xi$ is actually an
oriented manifold. A closer inspection shows that $K_\xi$ is also a manifold near points
in the interiors of $(n-2)$-simplices, so the nonmanifold points of $K_\xi$ in fact lie in
the $(n-3)$-skeleton. However, near points in the interiors of $(n-3)$-simplices it can
very well happen that $K_\xi$ is not a manifold.

In particular, elements of $H_1(X)$ are represented by collections of oriented loops
in $X$, and elements of $H_2(X)$ are represented by maps of closed oriented surfaces
into $X$. With a bit more work it can be shown that an oriented 1-cycle $\bigsqcup_\alpha S^1_\alpha \to X$ is
zero in $H_1(X)$ iff it extends to a map of a compact oriented surface with boundary
$\bigsqcup_\alpha S^1_\alpha$ into $X$. The analogous statement for 2-cycles is also true. In the early days of
homology theory it may have been believed, or at least hoped, that this close connection
with manifolds continued in all higher dimensions, but this has turned out not to
be the case. There is a sort of homology theory built from manifolds, called \textit{bordism},
but it is quite a bit more complicated than the homology theory we are studying here.

After these preliminary remarks let us begin to see what can be proved about
singular homology.

\begin{proposition}[Homology of disjoint union]
\label{prop:homology-disjoint-union}
\dcref{ch2:2.6}
\group{Singular Homology}
\uses{lem:boundary-composition-zero}
Corresponding to the decomposition of a space $X$ into its path-
components $X_\alpha$ there is an isomorphism of $H_n(X)$ with the direct sum $\bigoplus_\alpha H_n(X_\alpha)$.
\end{proposition}
\begin{proof}
 Since a singular simplex always has path-connected image, $C_n(X)$ splits as the
direct sum of its subgroups $C_n(X_\alpha)$. The boundary maps $\partial_n$ preserve this direct sum
decomposition, taking $C_n(X_\alpha)$ to $C_{n-1}(X_\alpha)$, so $\ker\partial_n$ and $\text{Im}\partial_{n+1}$ split similarly as
direct sums, hence the homology groups also split, $H_n(X) \cong \bigoplus_\alpha H_n(X_\alpha)$.
\end{proof}
\begin{proposition}[Zero-dimensional homology of path-connected space]
\label{prop:h0-path-connected}
\dcref{ch2:2.7}
\group{Singular Homology}
\uses{prop:homology-disjoint-union}
If $X$ is nonempty and path-connected, then $H_0(X) \cong \mathbb{Z}$. Hence for
any space $X$, $H_0(X)$ is a direct sum of $\mathbb{Z}$'s, one for each path-component of $X$.
\end{proposition}
\begin{proof}
 By definition, $H_0(X) = C_0(X) / \text{Im} \partial_1$ since $\partial_0 = 0$. Define a homomorphism
$\varepsilon : C_0(X) \to \mathbb{Z}$ by $\varepsilon(\sum_i n_i \sigma_i) = \sum_i n_i$. This is obviously surjective if $X$ is nonempty.
The claim is that $\ker \varepsilon = \text{Im} \partial_1$ if $X$ is path-connected, and hence $\varepsilon$ induces an isomorphism $H_0(X) \cong \mathbb{Z}$.

To verify the claim, observe first that $\text{Im} \partial_1 \subset \ker\varepsilon$ since for a singular 1-simplex
$\sigma:\Delta^1 \to X$ we have $\varepsilon \partial_1(\sigma) = \varepsilon(\sigma|[v_1] - \sigma|[v_0]) = 1 - 1 = 0$. For the reverse
inclusion $\ker\varepsilon \subset \text{Im}\partial_1$, suppose $\varepsilon(\sum_i n_i \sigma_i) = 0$, so $\sum_i n_i = 0$. The $\sigma_i$'s are singular
0-simplices, which are simply points of $X$. Choose a path $\tau_i : I \to X$ from a basepoint

%% --- END page-0118.tex ---
%% --- BEGIN page-0119.tex ---
$x_0$ to $\sigma_i(v_0)$ and let $\sigma_0$ be the singular 0-simplex with image $x_0$. We can view $\tau_i$
as a singular 1-simplex, a map $\tau_i:[v_0, v_1]\to X$, and then we have $\partial\tau_i = \sigma_i - \sigma_0$.
Hence $\partial(\sum_i n_i\tau_i) = \sum_i n_i\sigma_i - \sum_i n_i\sigma_0 = \sum_i n_i\sigma_i$ since $\sum_i n_i = 0$. Thus $\sum_i n_i\sigma_i$ is a
boundary, which shows that $\text{Ker}\partial \subset \text{Im}\,\partial_1$.
\end{proof}
\begin{proposition}[Homology of a point]
\label{prop:homology-point}
\dcref{ch2:2.8}
\group{Singular Homology}
\uses{}
If $X$ is a point, then $H_n(X) = 0$ for $n > 0$ and $H_0(X) \cong \mathbb{Z}$.
\end{proposition}
\begin{proof}
 In this case there is a unique singular $n$-simplex $\sigma_n$ for each $n$, and $\partial(\sigma_n) =
\sum_i(-1)^i\sigma_{n-1}$, a sum of $n+1$ terms, which is therefore 0 for $n$ odd and $\sigma_{n-1}$ for $n$
even, $n \ne 0$. Thus we have the chain complex
$$\cdots \to \mathbb{Z} \xrightarrow{\cong} \mathbb{Z} \xrightarrow{0} \mathbb{Z} \xrightarrow{\cong} \mathbb{Z} \xrightarrow{0} \mathbb{Z} \to 0$$
with boundary maps alternately isomorphisms and trivial maps, except at the last $\mathbb{Z}$.
The homology groups of this complex are trivial except for $H_0 \cong \mathbb{Z}$.
\end{proof}
It is often very convenient to have a slightly modified version of homology for
which a point has trivial homology groups in all dimensions, including zero. This is
done by defining the \textbf{reduced homology groups} $\tilde H_n(X)$ to be the homology groups
of the augmented chain complex
$$\cdots \to C_2(X) \xrightarrow{\partial_2} C_1(X) \xrightarrow{\partial_1} C_0(X) \xrightarrow{\varepsilon} \mathbb{Z} \to 0$$
where $\varepsilon(\sum_i n_i\sigma_i) = \sum_i n_i$ as in the proof of Proposition 2.7. Here we had better
require $X$ to be nonempty, to avoid having a nontrivial homology group in dimension
$-1$. Since $\varepsilon\partial_1 = 0$, $\varepsilon$ vanishes on $\text{Im}\,\partial_1$ and hence induces a map $H_0(X) \to \mathbb{Z}$ with
kernel $\tilde H_0(X)$, so $H_0(X) \cong \tilde H_0(X) \oplus \mathbb{Z}$. Obviously $H_n(X) \cong \tilde H_n(X)$ for $n > 0$.

Formally, one can think of the extra $\mathbb{Z}$ in the augmented chain complex as generated by the unique map $[\varnothing]\to X$ where $[\varnothing]$ is the empty simplex, with no vertices.
The augmentation map $\varepsilon$ is then the usual boundary map since $\partial[v_0] = [\bar v_0] = [\varnothing]$.

Readers who know about the fundamental group $\pi_1(X)$ may wish to make a
detour here to look at $\S2.A$ where it is shown that $H_1(X)$ is the abelianization of
$\pi_1(X)$ whenever $X$ is path-connected. This result will not be needed elsewhere in the
chapter, however.

\section*{Homotopy Invariance}

The first substantial result we will prove about singular homology is that homotopy equivalent spaces have isomorphic homology groups. This will be done by
showing that a map $f:X \to Y$ induces a homomorphism $f_*:H_n(X) \to H_n(Y)$ for each
$n$, and that $f_*$ is an isomorphism if $f$ is a homotopy equivalence.

For a map $f:X\to Y$, an induced homomorphism $f_*:C_n(X)\to C_n(Y)$ is defined
by composing each singular $n$-simplex $\sigma:\Delta^n \to X$ with $f$ to get a singular $n$-simplex
%% --- END page-0119.tex ---
%% --- BEGIN page-0120.tex ---
$f_*(\sigma) = f\circ\sigma:\Delta^n \to Y$, then extending $f_*$ linearly via $f_*(\sum_i n_i\sigma_i) = \sum_i n_if_*(\sigma_i) =
\sum_i n_if\circ\sigma_i$. The maps $f_*:C_n(X)\to C_n(Y)$ satisfy $f_*\partial = \partial f_*$ since
$$f_*(\partial(\sigma)) = f_*(\sum_i(-1)^i\sigma|[v_0,\cdots,\hat v_i,\cdots,v_n])
= \sum_i(-1)^if\circ\sigma|[v_0,\cdots,\hat v_i,\cdots,v_n] = \partial f_*(\sigma)$$

Thus we have a diagram
$$\cdots \to C_{n+1}(X) \xrightarrow{\partial} C_n(X) \xrightarrow{\partial} C_{n-1}(X) \to \cdots$$
$$\phantom{\cdots} \downarrow f_* \phantom{\cdots} \downarrow f_* \phantom{\cdots} \downarrow f_*$$
$$\cdots \to C_{n+1}(Y) \xrightarrow{\partial} C_n(Y) \xrightarrow{\partial} C_{n-1}(Y) \to \cdots$$

such that in each square the composition $f_*\circ\partial$ equals the composition $\partial\circ f_*$. A diagram
of maps with the property that any two compositions of maps starting at one point in
the diagram and ending at another are equal is called a \textbf{commutative diagram}. In the
present case commutativity of the diagram is equivalent to the commutativity relation
$f_*\circ\partial = \partial\circ f_*$, but commutative diagrams can contain commutative triangles, pentagons,
etc., as well as commutative squares.

The fact that the maps $f_*:C_n(X)\to C_n(Y)$ satisfy $f_*\partial = \partial f_*$ is also expressed
by saying that the $f_*$'s define a \textbf{chain map} from the singular chain complex of $X$
to that of $Y$. The relation $f_*\partial = \partial f_*$ implies that $f_*$ takes cycles to cycles since
$\partial\alpha = 0$ implies $\partial(f_*(\alpha)) = f_*(\partial\alpha) = 0$. Also, $f_*$ takes boundaries to boundaries
since $f_*(\partial\beta) = \partial(f_*(\beta))$. Hence $f_*$ induces a homomorphism $f_*:H_n(X)\to H_n(Y)$. An
algebraic statement of what we have just proved is:

\begin{proposition}[Chain map induces homology homomorphism]
\label{prop:chain-map-induces-homology-homomorphism}
\dcref{ch2:2.9}
\group{Singular Homology}
\uses{lem:boundary-composition-zero}
A chain map between chain complexes induces homomorphisms
between the homology groups of the two complexes. $\square$
\end{proposition}

Two basic properties of induced homomorphisms which are important in spite
of being rather trivial are:
\begin{enumerate}
\item[{\rm (i)}] $(fg)_* = f_*g_*$ for a composed mapping $X \xrightarrow{g} Y \xrightarrow{f} Z$. This follows from
associativity of compositions $\Delta^n \xrightarrow{\sigma} X \xrightarrow{g} Y \xrightarrow{f} Z$.
\item[{\rm (ii)}] $\mathbb{1}_* = \mathbb{1}$ where $\mathbb{1}$ denotes the identity map of a space or a group.
\end{enumerate}

Less trivially, we have:

\begin{theorem}[Homotopic maps induce the same homology map]
\label{thm:homotopic-maps-induce-same-homology-map}
\dcref{ch2:2.10}
\group{Singular Homology}
\uses{def:homotopy,prop:chain-map-induces-homology-homomorphism}
If two maps $f,g:X \to Y$ are homotopic, then they induce the same
homomorphism $f_* = g_*:H_n(X) \to H_n(Y)$.
\end{theorem}

In view of the formal properties $(fg)_* = f_*g_*$ and $\mathbb{1}_* = \mathbb{1}$, this immediately
implies:

\begin{corollary}[Homotopy equivalence induces homology isomorphism]
\label{cor:homotopy-equivalence-induces-homology-isomorphism}
\dcref{ch2:2.11}
\group{Singular Homology}
\uses{def:homotopy-equivalence,thm:homotopic-maps-induce-same-homology-map}
The maps $f_*:H_n(X)\to H_n(Y)$ induced by a homotopy equivalence
$f:X\to Y$ are isomorphisms for all $n$. $\square$
\end{corollary}

For example, if $X$ is contractible then $\tilde H_n(X) = 0$ for all $n$.

%% --- END page-0120.tex ---
%% --- BEGIN page-0121.tex ---
\begin{proof}[2.10]
 The essential ingredient is a procedure for
subdividing $\Delta^n \times I$ into simplices. The figure shows the
cases $n = 1, 2$. In $\Delta^n \times I$, let $\Delta^n \times \{0\} = [v_0,\cdots,v_n]$
and $\Delta^n \times \{1\} = [w_0,\cdots,w_n]$, where $v_i$ and $w_i$ have the
same image under the projection $\Delta^n \times I \to \Delta^n$. We can pass
from $[v_0,\cdots,v_n]$ to $[w_0,\cdots,w_n]$ by interpolating a sequence of $n$-simplices, each obtained from the preceding
one by moving one vertex $v_i$ up to $w_i$, starting with $v_n$
and working backwards to $v_0$. Thus the first step is to
move $[v_0,\cdots,v_n]$ up to $[v_0,\cdots,v_{n-1},w_n]$, then the second step is to move this up to $[v_0,\cdots,v_{n-2},w_{n-1},w_n]$,
and so on. In the typical step $[v_0,\cdots,v_i,w_{i+1},\cdots,w_n]$
moves up to $[v_0,\cdots,v_{i-1},w_i,\cdots,w_n]$. The region between these two $n$-simplices is exactly the $(n+1)$-simplex
$[v_0,\cdots,v_i,w_i,\cdots,w_n]$ which has $[v_0,\cdots,v_i,w_{i+1},\cdots,w_n]$ as its lower face and
$[v_0,\cdots,v_{i-1},w_i,\cdots,w_n]$ as its upper face. Altogether, $\Delta^n \times I$ is the union of the
$(n+1)$-simplices $[v_0,\cdots,v_i,w_i,\cdots,w_n]$, each intersecting the next in an $n$-simplex
face.

Given a homotopy $F:X \times I \to Y$ from $f$ to $g$ and a singular simplex $\sigma:\Delta^n \to X$,
we can form the composition $F \circ (\sigma \times \mathbb{1}):\Delta^n \times I \to X \times I \to Y$. Using this, we can define
\textbf{prism operators} $P:C_n(X)\to C_{n+1}(Y)$ by the following formula:
$$P(\sigma) = \sum_i(-1)^i F \circ (\sigma \times \mathbb{1})|[v_0,\cdots,v_i,w_i,\cdots,w_n]$$

We will show that these prism operators satisfy the basic relation
$$\partial P = g_* - f_* - P\partial$$

Geometrically, the left side of this equation represents the boundary of the prism, and
the three terms on the right side represent the top $\Delta^n \times \{1\}$, the bottom $\Delta^n \times \{0\}$, and
the sides $\partial\Delta^n \times I$ of the prism. To prove the relation we calculate
$$\partial P(\sigma) = \sum_j\sum_i(-1)^i(-1)^jF \circ (\sigma \times \mathbb{1})|[v_0,\cdots,\hat v_j,\cdots,v_i,w_i,\cdots,w_n]$$
$$+ \sum_{j\geq i}(-1)^i(-1)^{j+1}F \circ (\sigma \times \mathbb{1})|[v_0,\cdots,v_i,w_i,\cdots,\hat w_j,\cdots,w_n]$$

The terms with $i = j$ in the two sums cancel except for $F \circ (\sigma \times \mathbb{1})|[\hat v_n,w_n,\cdots,w_n]$,
which is $g_*\circ\sigma = g_*(\sigma)$, and $-F \circ (\sigma \times \mathbb{1})|[v_0,\cdots,v_n,\hat w_n]$, which is $-f_*\circ\sigma = -f_*(\sigma)$.
The terms with $i \ne j$ are exactly $-P\partial(\sigma)$ since
$$P\partial(\sigma) = \sum_i(-1)^i(-1)^jF \circ (\sigma \times \mathbb{1})|[v_0,\cdots,v_i,w_i,\cdots,\hat w_j,\cdots,w_n]$$
$$+ \sum_{i>j}(-1)^{i-1}(-1)^jF \circ (\sigma \times \mathbb{1})|[v_0,\cdots,\hat v_j,\cdots,v_i,w_i,\cdots,w_n]$$

%% --- END page-0121.tex ---
%% --- BEGIN page-0122.tex ---
Now we can finish the proof of the theorem. If $\alpha \in C_n(X)$ is a cycle, then we
have $g_*(\alpha) - f_*(\alpha) = \partial P(\alpha) + P\partial(\alpha) = \partial P(\alpha)$ since $\partial\alpha = 0$. Thus $g_*(\alpha) - f_*(\alpha)$ is
a boundary, so $g_*(\alpha)$ and $f_*(\alpha)$ determine the same homology class, which means
that $g_*$ equals $f_*$ on the homology class of $\alpha$.
\end{proof}
The relationship $\partial P + P\partial = g_* - f_*$ is expressed by saying that $P$ is a \textbf{chain homotopy}
between the chain maps $f_*$ and $g_*$. We have just shown:

\begin{proposition}[Chain-homotopic maps induce the same homology map]
\label{prop:chain-homotopic-maps-induce-same-homology-map}
\dcref{ch2:2.12}
\group{Singular Homology}
\uses{prop:chain-map-induces-homology-homomorphism}
Chain-homotopic chain maps induce the same homomorphism on
homology. $\square$
\end{proposition}

There are also induced homomorphisms $f_*:\tilde H_n(X) \to \tilde H_n(Y)$ for reduced homology groups since $f_*\varepsilon = \varepsilon f_*$ where $f_*$ is the identity map on the added groups $\mathbb{Z}$ in the
augmented chain complexes. The properties of induced homomorphisms we proved
above hold equally well in the setting of reduced homology, with the same proofs.

\section*{Exact Sequences and Excision}

If there was always a simple relationship between the homology groups of a space
$X$, a subspace $A$, and the quotient space $X/A$, then this could be a very useful tool
in understanding the homology groups of spaces such as CW complexes that can be
built inductively from successively more complicated subspaces. Perhaps the simplest
possible relationship would be if $H_n(X)$ contained $H_n(A)$ as a subgroup and the
quotient group $H_n(X)/H_n(A)$ was isomorphic to $H_n(X/A)$. While this does hold
in some cases, if it held in general then homology theory would collapse totally since
every space $X$ can be embedded as a subspace of a space with trivial homology groups,
namely the cone $CX = (X \times I)/(X \times \{0\})$, which is contractible.

It turns out that this overly simple model does not have to be modified too much
to get a relationship that is valid in fair generality. The novel feature of the actual
relationship is that it involves the groups $H_n(X)$, $H_n(A)$, and $H_n(X/A)$ for all values
of $n$ simultaneously. In practice this is not as bad as it might sound, and in addition
it has the pleasant side effect of sometimes allowing higher-dimensional homology
groups to be computed in terms of lower-dimensional groups which may already be
known, for example by induction.

In order to formulate the relationship we are looking for, we need an algebraic
definition which is central to algebraic topology. A sequence of homomorphisms
$$\cdots \to A_{n+1} \xrightarrow{\alpha_{n+1}} A_n \xrightarrow{\alpha_n} A_{n-1} \to \cdots$$
is said to be exact if $\text{Ker}\,\alpha_n = \text{Im}\,\alpha_{n+1}$ for each $n$. The inclusions $\text{Im}\,\alpha_{n+1} \subset \text{Ker}\,\alpha_n$
are equivalent to $\alpha_n\alpha_{n+1} = 0$, so the sequence is a chain complex, and the opposite
inclusions $\text{Ker}\,\alpha_n \subset \text{Im}\,\alpha_{n+1}$ say that the homology groups of this chain complex are
trivial.

%% --- END page-0122.tex ---
%% --- BEGIN page-0123.tex ---
A number of basic algebraic concepts can be expressed in terms of exact sequences, for example:

\begin{enumerate}
\item[{\rm (i)}] $0 \to A \xrightarrow{\alpha} B$ is exact iff $\text{Ker}\,\alpha = 0$, i.e., $\alpha$ is injective.
\item[{\rm (ii)}] $A \xrightarrow{\alpha} B \to 0$ is exact iff $\text{Im}\,\alpha = B$, i.e., $\alpha$ is surjective.
\item[{\rm (iii)}] $0 \to A \xrightarrow{\alpha} B \to 0$ is exact iff $\alpha$ is an isomorphism, by (i) and (ii).
\item[{\rm (iv)}] $0 \to A \xrightarrow{\alpha} B \xrightarrow{\beta} C \to 0$ is exact iff $\alpha$ is injective, $\beta$ is surjective, and $\text{ker}\,\beta =
\text{Im}\,\alpha$, so $\beta$ induces an isomorphism $C \cong B/\text{Im}\,\alpha$. This can be written $C \cong B/A$ if
we think of $\alpha$ as an inclusion of $A$ as a subgroup of $B$.
\end{enumerate}

An exact sequence $0 \to A \to B \to C \to 0$ as in (iv) is called a \textbf{short exact sequence}.

Exact sequences provide the right tool to relate the homology groups of a space,
a subspace, and the associated quotient space:

\begin{theorem}[Excision Theorem]
\label{thm:excision}
\dcref{ch2:2.13}
\group{Singular Homology}
\uses{def:deformation-retraction}
If $X$ is a space and $A$ is a nonempty closed subspace that is a deformation retract of some neighborhood in $X$, then there is an exact sequence
$$\cdots \to \tilde H_n(A) \xrightarrow{i_*} \tilde H_n(X) \xrightarrow{j_*} \tilde H_n(X/A) \xrightarrow{\partial} \tilde H_{n-1}(A) \xrightarrow{i_*} \tilde H_{n-1}(X) \to \cdots$$
$$\cdots \to \tilde H_0(X/A) \to 0$$
where $i$ is the inclusion $A \hookrightarrow X$ and $j$ is the quotient map $X \to X/A$.
\end{theorem}

The map $\partial$ will be constructed in the course of the proof. The idea is that an
element $x \in \tilde H_n(X/A)$ can be represented by a chain $\alpha$ in $X$ with $\partial\alpha$ a cycle in $A$
whose homology class is $\partial x \in \tilde H_{n-1}(A)$.

Pairs of spaces $(X,A)$ satisfying the hypothesis of the theorem will be called
\textbf{good pairs}. For example, if $X$ is a CW complex and $A$ is a nonempty subcomplex,
then $(X,A)$ is a good pair by Proposition A.5 in the Appendix.

\begin{corollary}[Homology of spheres]
\label{cor:homology-spheres}
\dcref{ch2:2.14}
\group{Singular Homology}
\uses{thm:excision}
$\tilde H_n(S^n) \cong \mathbb{Z}$ and $\tilde H_i(S^n) = 0$ for $i \ne n$.
\end{corollary}
\begin{proof}
 For $n > 0$ take $(X,A) = (D^n,S^{n-1})$ so $X/A = S^n$. The terms $\tilde H_i(D^n)$ in the
long exact sequence for this pair are zero since $D^n$ is contractible. Exactness of the
sequence then implies that the maps $\tilde H_i(S^n) \xrightarrow{\partial} \tilde H_{i-1}(S^{n-1})$ are isomorphisms for
$i > 0$ and that $\tilde H_0(S^n) = 0$. The result now follows by induction on $n$, starting with
the case of $S^0$ where the result holds by Propositions 2.6 and 2.8.
\end{proof}
As an application of this calculation we have the following classical theorem of
Brouwer, the 2-dimensional case of which was proved in $\S1.1$.

\begin{corollary}[Non-retraction of disk onto boundary]
\label{cor:non-retraction-disk-boundary}
\dcref{ch2:2.15}
\group{Singular Homology}
\uses{cor:homology-spheres}
$\partial D^n$ is not a retract of $D^n$. Hence every map $f:D^n \to D^n$ has a
fixed point.
\end{corollary}
\begin{proof}
 If $r:D^n \to \partial D^n$ is a retraction, then $r \circ i = \mathbb{1}$ for $i:\partial D^n \to D^n$ the inclusion map.
The composition $\tilde H_{n-1}(\partial D^n) \xrightarrow{i_*} \tilde H_{n-1}(D^n) \xrightarrow{r_*} \tilde H_{n-1}(\partial D^n)$ is then the identity map

%% --- END page-0123.tex ---
%% --- BEGIN page-0124.tex ---
on $\tilde H_{n-1}(\partial D^n) \cong \mathbb{Z}$. But $i_*$ and $r_*$ are both 0 since $\tilde H_{n-1}(D^n) = 0$, and we have a
contradiction. The statement about fixed points follows as in Theorem 1.9.
\end{proof}
The derivation of the exact sequence of homology groups for a good pair $(X,A)$
will be rather a long story. We will in fact derive a more general exact sequence which
holds for arbitrary pairs $(X,A)$, but with the homology groups of the quotient space
$X/A$ replaced by \textit{relative homology groups}, denoted $H_n(X,A)$. These turn out to be
quite useful for many other purposes as well.

\section*{Relative Homology Groups}

It sometimes happens that by ignoring a certain amount of data or structure one
obtains a simpler, more flexible theory which, almost paradoxically, can give results
not readily obtainable in the original setting. A familiar instance of this is arithmetic
mod $n$, where one ignores multiples of $n$. Relative homology is another example. In
this case what one ignores is all singular chains in a subspace of the given space.

Relative homology groups are defined in the following way. Given a space $X$ and
a subspace $A \subset X$, let $C_n(X,A)$ be the quotient group $C_n(X)/C_n(A)$. Thus chains in
$A$ are trivial in $C_n(X,A)$. Since the boundary map $\partial:C_n(X) \to C_{n-1}(X)$ takes $C_n(A)$
to $C_{n-1}(A)$, it induces a quotient boundary map $\partial:C_n(X,A) \to C_{n-1}(X,A)$. Letting $n$
vary, we have a sequence of boundary maps
$$\cdots \to C_n(X,A) \xrightarrow{\partial} C_{n-1}(X,A) \to \cdots$$

The relation $\partial^2 = 0$ holds for these boundary maps since it holds before passing to
quotient groups. So we have a chain complex, and the homology groups Ker$\partial$/Im$\partial$
of this chain complex are by definition the \textbf{relative homology groups} $H_n(X,A)$. By
considering the definition of the relative boundary map we see:

\begin{itemize}
\item Elements of $H_n(X,A)$ are represented by \textbf{relative cycles}: $n$-chains $\alpha \in C_n(X)$
such that $\partial\alpha \in C_{n-1}(A)$.
\item A relative cycle $\alpha$ is trivial in $H_n(X,A)$ iff it is a \textbf{relative boundary}: $\alpha = \partial\beta + \gamma$
for some $\beta \in C_{n+1}(X)$ and $\gamma \in C_n(A)$.
\end{itemize}

These properties make precise the intuitive idea that $H_n(X,A)$ is `homology of $X$
modulo $A$'.

The quotient $C_n(X)/C_n(A)$ could also be viewed as a subgroup of $C_n(X)$, the
subgroup with basis the singular $n$-simplices $\sigma:\Delta^n \to X$ whose image is not contained in $A$. However, the boundary map does not take this subgroup of $C_n(X)$ to
the corresponding subgroup of $C_{n-1}(X)$, so it is usually better to regard $C_n(X,A)$ as
a quotient rather than a subgroup of $C_n(X)$.

Our goal now is to show that the relative homology groups $H_n(X,A)$ for any pair
$(X,A)$ fit into a long exact sequence
$$\cdots \to H_n(A) \to H_n(X) \to H_n(X,A) \to H_{n-1}(A) \to H_{n-1}(X) \to \cdots$$
$$\cdots \to H_0(X,A) \to 0$$

%% --- END page-0124.tex ---
%% --- BEGIN page-0125.tex ---
This will be entirely a matter of algebra. To start the process, consider the diagram
\[
0 \longrightarrow C_n(A) \xrightarrow{i} C_n(X) \xrightarrow{j} C_n(X,A) \longrightarrow 0
\]
\[
\begin{array}{ccccccc}
& & \partial \downarrow & & \partial \downarrow & & \partial \downarrow \\
0 \longrightarrow C_{n-1}(A) &\xrightarrow{i}& C_{n-1}(X) &\xrightarrow{j}& C_{n-1}(X,A) &\longrightarrow& 0
\end{array}
\]
where $i$ is inclusion and $j$ is the quotient map. The diagram is commutative by the definition of the boundary maps. Letting $n$ vary, and drawing these short exact sequences vertically rather than horizontally, we have a large commutative diagram of the form shown at the right, where the columns are exact and the rows are chain complexes which we denote $A$, $B$, and $C$. Such a diagram is called a \textbf{short exact sequence of chain complexes}. We will show that when we pass to homology groups, this short exact sequence of chain complexes stretches out into a long exact sequence of homology groups

\[
\cdots \to H_n(A) \xrightarrow{i_*} H_n(B) \xrightarrow{j_*} H_n(C) \xrightarrow{\partial} H_{n-1}(A) \xrightarrow{i_*} H_{n-1}(B) \to \cdots
\]

where $H_n(A)$ denotes the homology group $\ker\partial/\operatorname{Im}\partial$ at $A_n$ in the chain complex $A$, and $H_n(B)$ and $H_n(C)$ are defined similarly.

The commutativity of the squares in the short exact sequence of chain complexes means that $i$ and $j$ are chain maps. These therefore induce maps $i_*$ and $j_*$ on homology. To define the boundary map $\partial:H_n(C)\to H_{n-1}(A)$, let $c\in C_n$ be a cycle. Since $j$ is onto, $c=j(b)$ for some $b\in B_n$. The element $\partial b \in B_{n-1}$ is in $\ker j$ since $j(\partial b) = \partial j(b) = \partial c = 0$. So $\partial b = i(a)$ for some $a \in A_{n-1}$ since $\ker j = \operatorname{Im}i$. Note that $\partial a = 0$ since $i(\partial a) = \partial i(a) = \partial\partial b = 0$ and $i$ is injective. We define $\partial:H_n(C)\to H_{n-1}(A)$ by sending the homology class of $c$ to the homology class of $a$, $\partial[c] = [a]$. This is well-defined since:



\begin{itemize}
\item The element $a$ is uniquely determined by $\partial b$ since $i$ is injective.
\item A different choice $b'$ for $b$ would have $j(b') = j(b)$, so $b'-b$ is in $\ker j = \operatorname{Im}i$. Thus $b' - b = i(a')$ for some $a'$, hence $b' = b + i(a')$. The effect of replacing $b$ by $b + i(a')$ is to change $a$ to the homologous element $a + \partial a'$ since $i(a + \partial a') = i(a) + i(\partial a') = \partial b + \partial i(a') = \partial(b + i(a'))$.
\item A different choice of $c$ within its homology class would have the form $c + \partial c'$. Since $c' = j(b')$ for some $b'$, we then have $c + \partial c' = c + \partial j(b') = c + j(\partial b') = j(b + \partial b')$, so $b$ is replaced by $b + \partial b'$, which leaves $\partial b$ and therefore also $a$ unchanged.
\end{itemize}

%% --- END page-0125.tex ---
%% --- BEGIN page-0126.tex ---
The map $\partial:H_n(C)\to H_{n-1}(A)$ is a homomorphism since if $\partial[c_1] = [a_1]$ and $\partial[c_2] = [a_2]$ via elements $b_1$ and $b_2$ as above, then $j(b_1 + b_2) = j(b_1) + j(b_2) = c_1 + c_2$ and $i(a_1 + a_2) = i(a_1) + i(a_2) = \partial b_1 + \partial b_2 = \partial(b_1 + b_2)$, so $\partial([c_1] + [c_2]) = [a_1] + [a_2]$.

\begin{theorem}[The sequence of homology groups]
\label{thm:long-exact-sequence-pair}
\dcref{ch2:2.16}
\group{Singular Homology}
\uses{lem:boundary-composition-zero}
\[
\cdots \to H_n(A) \xrightarrow{i_*} H_n(B) \xrightarrow{j_*} H_n(C) \xrightarrow{\partial} H_{n-1}(A) \xrightarrow{i_*} H_{n-1}(B) \to \cdots
\]
\textbf{is exact.}
\end{theorem}
\begin{proof}
: There are six things to verify:

$\operatorname{Im}i_* \subset \ker j_*$. This is immediate since $ji = 0$ implies $j_* i_* = 0$.

$\operatorname{Im}j_* \subset \ker\partial$. We have $\partial j_* = 0$ since in this case $\partial b = 0$ in the definition of $\partial$.

$\operatorname{Im}\partial \subset \ker i_*$. Here $i_*\partial = 0$ since $i_*\partial$ takes $[c]$ to $[\partial b] = 0$.

$\ker j_* \subset \operatorname{Im}i_*$. A homology class in $\ker j_*$ is represented by a cycle $b \in B_n$ with $j(b)$ a boundary, so $j(b) = \partial c'$ for some $c' \in C_{n+1}$. Since $j$ is surjective, $c' = j(b')$ for some $b' \in B_{n+1}$. We have $j(b - \partial b') = j(b) - j(\partial b') = j(b) - \partial j(b') = 0$ since $\partial j(b') = \partial c' = j(b)$. So $b - \partial b' = i(a)$ for some $a \in A_n$. This $a$ is a cycle since $i(\partial a) = \partial i(a) = \partial(b - \partial b') = \partial b = 0$ and $i$ is injective. Thus $i_*[a] = [b - \partial b'] = [b]$, showing that $i_*$ maps onto $\ker j_*$.

$\ker\partial \subset \operatorname{Im}j_*$. In the notation used in the definition of $\partial$, if $c$ represents a homology class in $\ker\partial$, then $a = \partial a'$ for some $a' \in A_n$. The element $b - i(a')$ is a cycle since $\partial(b - i(a')) = \partial b - \partial i(a') = \partial b - i(\partial a') = \partial b - i(a) = 0$. And $j(b - i(a')) = j(b) - ji(a') = j(b) = c$, so $j_*$ maps $[b - i(a')]$ to $[c]$.

$\ker i_* \subset \operatorname{Im}\partial$. Given a cycle $a \in A_{n-1}$ such that $i(a) = \partial b$ for some $b \in B_n$, then $j(b)$ is a cycle since $\partial j(b) = j(\partial b) = ji(a) = 0$, and $\partial$ takes $[j(b)]$ to $[a]$.
\end{proof}
This theorem represents the beginnings of the subject of homological algebra. The method of proof is sometimes called \textit{diagram chasing}.

Returning to topology, the preceding algebraic theorem yields a long exact sequence of homology groups:

\[
\cdots \to H_n(A) \xrightarrow{i_*} H_n(X) \xrightarrow{j_*} H_n(X,A) \xrightarrow{\partial} H_{n-1}(A) \xrightarrow{i_*} H_{n-1}(X) \to \cdots
\]
\[
\cdots \to H_0(X,A) \to 0
\]

The boundary map $\partial:H_n(X,A)\to H_{n-1}(A)$ has a very simple description: If a class $[\alpha] \in H_n(X,A)$ is represented by a relative cycle $\alpha$, then $\partial[\alpha]$ is the class of the cycle $\partial\alpha$ in $H_{n-1}(A)$. This is immediate from the algebraic definition of the boundary homomorphism in the long exact sequence of homology groups associated to a short exact sequence of chain complexes.

This long exact sequence makes precise the idea that the groups $H_n(X,A)$ measure the difference between the groups $H_n(X)$ and $H_n(A)$. In particular, exactness

%% --- END page-0126.tex ---
%% --- BEGIN page-0127.tex ---
implies that if $H_n(X,A) = 0$ for all $n$, then the inclusion $A\hookrightarrow X$ induces isomorphisms $H_n(A) \simeq H_n(X)$ for all $n$, by the remark (iii) following the definition of exactness. The converse is also true according to an exercise at the end of this section.

There is a completely analogous long exact sequence of reduced homology groups for a pair $(X,A)$ with $A \ne \emptyset$. This comes from applying the preceding algebraic machinery to the short exact sequence of chain complexes formed by the short exact sequences $0\to C_n(A)\to C_n(X)\to C_n(X,A)\to 0$ in nonnegative dimensions, augmented by the short exact sequence $0 \to \Z \xrightarrow{1} \Z \to 0 \to 0$ in dimension $-1$. In particular, this means that $\tilde H_n(X,A)$ is the same as $H_n(X,A)$ for all $n$, when $A \ne \emptyset$.

\begin{example}[Long exact sequence of disk and sphere pair]
\label{ex:long-exact-sequence-disk-sphere}
\dcref{ch2:2.17}
\group{Singular Homology}
\uses{thm:long-exact-sequence-pair}
In the long exact sequence of reduced homology groups for the pair $(D^n,\partial D^n)$, the maps $H_i(D^n,\partial D^n) \xrightarrow{\partial} \tilde H_{i-1}(S^{n-1})$ are isomorphisms for all $i > 0$ since the remaining terms $\tilde H_i(D^n)$ are zero for all $i$. Thus we obtain the calculation
\[
H_i(D^n,\partial D^n) \simeq \begin{cases}
\Z & \text{for } i = n \\
0 & \text{otherwise}
\end{cases}
\]
\end{example}

\begin{example}[Reduced homology isomorphic to relative homology of basepoint]
\label{ex:reduced-homology-relative}
\dcref{ch2:2.18}
\group{Singular Homology}
\uses{thm:long-exact-sequence-pair}
Applying the long exact sequence of reduced homology groups to a pair $(X,x_0)$ with $x_0 \in X$ yields isomorphisms $H_n(X,x_0) \simeq \tilde H_n(X)$ for all $n$ since $\tilde H_n(x_0) = 0$ for all $n$.
\end{example}

There are induced homomorphisms for relative homology just as there are in the nonrelative, or `absolute', case. A map $f:X\to Y$ with $f(A) \subset B$, or more concisely $f:(X,A)\to(Y,B)$, induces homomorphisms $f_\sharp:C_n(X,A)\to C_n(Y,B)$ since the chain map $f_\sharp:C_n(X)\to C_n(Y)$ takes $C_n(A)$ to $C_n(B)$, so we get a well-defined map on quotients, $f_\sharp:C_n(X,A)\to C_n(Y,B)$. The relation $f_\sharp\partial = \partial f_\sharp$ holds for relative chains since it holds for absolute chains. By Proposition 2.9 we then have induced homomorphisms $f_*:H_n(X,A)\to H_n(Y,B)$.

\begin{proposition}[If two maps $f,g:(X,A)\to(Y,B)$ are homotopic through maps of pairs $(X,A)\to(Y,B)$, then $f_* = g_*:H_n(X,A)\to H_n(Y,B)$]
\label{prop:homotopic-maps-pairs}
\dcref{ch2:2.19}
\group{Singular Homology}
\uses{def:homotopy,thm:homotopic-maps-induce-same-homology-map}

\end{proposition}\textbf{Proof}: The prism operator $P$ from the proof of Theorem 2.10 takes $C_n(A)$ to $C_{n+1}(B)$, hence induces a relative prism operator $P:C_n(X,A)\to C_{n+1}(Y,B)$. Since we are just passing to quotient groups, the formula $\partial P + P\partial = g_\sharp - f_\sharp$ remains valid. Thus the maps $f_\sharp$ and $g_\sharp$ on relative chain groups are chain homotopic, and hence they induce the same homomorphism on relative homology groups. \qed

An easy generalization of the long exact sequence of a pair $(X,A)$ is the long exact sequence of a triple $(X,A,B)$, where $B \subset A \subset X$:

\[
\cdots \to H_n(A,B) \to H_n(X,B) \to H_n(X,A) \to H_{n-1}(A,B) \to \cdots
\]

This is the long exact sequence of homology groups associated to the short exact sequence of chain complexes formed by the short exact sequences
\[
0 \to C_n(A,B) \to C_n(X,B) \to C_n(X,A) \to 0
\]

%% --- END page-0127.tex ---
%% --- BEGIN page-0128.tex ---
For example, taking $B$ to be a point, the long exact sequence of the triple $(X,A,B)$ becomes the long exact sequence of reduced homology for the pair $(X,A)$.

\textbf{Excision}

A fundamental property of relative homology groups is given by the following \textbf{Excision Theorem}, describing when the relative groups $H_n(X,A)$ are unaffected by deleting, or excising, a subset $Z \subset A$.

\begin{theorem}[Given subspaces $Z \subset A \subset X$ such that the closure of $Z$ is contained in the interior of $A$, then the inclusion $(X-Z,A-Z) \hookrightarrow (X,A)$ induces isomorphisms $H_n(X-Z,A-Z)\to H_n(X,A)$ for all $n$. Equivalently, for subspaces $A, B \subset X$ whose interiors cover $X$, the inclusion $(B,A \cap B) \hookrightarrow (X,A)$ induces isomorphisms $H_n(B,A\cap B)\to H_n(X,A)$ for all $n$]
\label{thm:excision-general}
\dcref{ch2:2.20}
\group{Singular Homology}
\uses{lem:boundary-composition-zero}
The translation between the two versions is obtained by setting $B = X - Z$ and $Z = X - B$. Then $A \cap B = A - Z$ and the condition $\text{cl}Z \subset \operatorname{int}A$ is equivalent to $X = \operatorname{int}A \cup \operatorname{int}B$ since $X - \operatorname{int}B = \text{cl}Z$.
\end{theorem}



The proof of the excision theorem will involve a rather lengthy technical detour involving a construction known as barycentric subdivision, which allows homology groups to be computed using small singular simplices. In a metric space `smallness' can be defined in terms of diameters, but for general spaces it will be defined in terms of covers.

For a space $X$, let $\mathcal{U} = \{U_j\}$ be a collection of subspaces of $X$ whose interiors form an open cover of $X$, and let $C_n^{\mathcal{U}}(X)$ be the subgroup of $C_n(X)$ consisting of chains $\sum_i n_i\sigma_i$ such that each $\sigma_i$ has image contained in some set in the cover $\mathcal{U}$. The boundary map $\partial:C_n(X)\to C_{n-1}(X)$ takes $C_n^{\mathcal{U}}(X)$ to $C_{n-1}^{\mathcal{U}}(X)$, so the groups $C_n^{\mathcal{U}}(X)$ form a chain complex. We denote the homology groups of this chain complex by $H_n^{\mathcal{U}}(X)$.

\begin{proposition}[Inclusion of $\mathcal{U}$-chains is a chain homotopy equivalence]
\label{prop:inclusion-u-chains-homotopy-equivalence}
\dcref{ch2:2.21}
\group{Singular Homology}
\uses{prop:chain-homotopic-maps-induce-same-homology-map}
The inclusion $\iota:C_n^{\mathcal{U}}(X) \hookrightarrow C_n(X)$ is a chain homotopy equivalence, that is, there is a chain map $\rho:C_n(X)\to C_n^{\mathcal{U}}(X)$ such that $\iota\rho$ and $\rho\iota$ are chain homotopic to the identity. Hence $\iota$ induces isomorphisms $H_n^{\mathcal{U}}(X) \simeq H_n(X)$ for all $n$.
\end{proposition}
\begin{proof}
: The barycentric subdivision process will be performed at four levels, beginning with the most geometric and becoming increasingly algebraic.

\textbf{(1) Barycentric Subdivision of Simplices.} The points of a simplex $[v_0,\ldots,v_n]$ are the linear combinations $\sum_i t_i v_i$ with $\sum_i t_i = 1$ and $t_i \ge 0$ for each $i$. The \textbf{barycenter} or `center of gravity' of the simplex $[v_0,\ldots,v_n]$ is the point $b = \sum_i t_i v_i$ whose barycentric coordinates $t_i$ are all equal, namely $t_i = 1/(n+1)$ for each $i$. The \textbf{barycentric subdivision} of $[v_0,\ldots,v_n]$ is the decomposition of $[v_0,\ldots,v_n]$ into the $n$-simplices $[b,w_0,\ldots,w_{n-1}]$ where, inductively, $[w_0,\ldots,w_{n-1}]$ is an $(n-1)$-simplex in the

%% --- END page-0128.tex ---
%% --- BEGIN page-0129.tex ---
barycentric subdivision of a face $[v_0,\ldots,\hat{v}_i,\ldots,v_n]$. The induction starts with the case $n=0$ when the barycentric subdivision of $[v_0]$ is defined to be just $[v_0]$ itself. The next two cases $n=1,2$ and part of the case $n=3$ are shown in the figure. It follows from the inductive definition that the vertices of simplices in the barycentric subdivision of $[v_0,\ldots,v_n]$ are exactly the barycenters of all the $k$-dimensional faces $[v_{i_0},\ldots,v_{i_k}]$ of $[v_0,\ldots,v_n]$ for $0 \le k \le n$. When $k=0$ this gives the original vertices $v_i$ since the barycenter of a $0$-simplex is itself. The barycenter of $[v_{i_0},\ldots,v_{i_k}]$ has barycentric coordinates $t_i = 1/(k+1)$ for $i = i_0,\ldots,i_k$ and $t_i = 0$ otherwise.



The $n$-simplices of the barycentric subdivision of $\Delta^n$, together with all their faces, do in fact form a $\Delta$-complex structure on $\Delta^n$, indeed a simplicial complex structure, though we shall not need to know this in what follows.

A fact we will need is that the diameter of each simplex of the barycentric subdivision of $[v_0,\ldots,v_n]$ is at most $n/(n+1)$ times the diameter of $[v_0,\ldots,v_n]$. Here the diameter of a simplex is by definition the maximum distance between any two of its points, and we are using the metric from the ambient Euclidean space $\R^m$ containing $[v_0,\ldots,v_n]$. The diameter of a simplex equals the maximum distance between any of its vertices because the distance between two points $v$ and $\sum_i t_i v_i$ of $[v_0,\ldots,v_n]$ satisfies the inequality

\[
|v - \sum_i t_i v_i| = |\sum_i t_i(v - v_i)| \le \sum_i t_i|v - v_i| \le \sum_i t_i \max_j |v - v_j| = \max_j|v - v_j|
\]

To obtain the bound $n/(n+1)$ on the ratio of diameters, we therefore need to verify that the distance between any two vertices $w_j$ and $w_k$ of a simplex $[w_0,\ldots,w_n]$ of the barycentric subdivision of $[v_0,\ldots,v_n]$ is at most $n/(n+1)$ times the diameter of $[v_0,\ldots,v_n]$. If neither $w_j$ nor $w_k$ is the barycenter $b$ of $[v_0,\ldots,v_n]$, then these two points lie in a proper face of $[v_0,\ldots,v_n]$ and we are done by induction on $n$. So we may suppose $w_j$, say, is the barycenter $b$, and then by the previous displayed inequality we may take $w_k$ to be a vertex $v_i$. Let $b_i$ be the barycenter of $[v_0,\ldots,\hat{v}_i,\ldots,v_n]$, with all barycentric coordinates equal to $1/n$ except for $t_i = 0$. Then we have $b = \frac{1}{n+1}v_i + \frac{n}{n+1}b_i$. The sum of the two coefficients is $1$, so $b$ lies on the line segment $[v_i,b_i]$ from $v_i$ to $b_i$, and the distance from $b$ to $v_i$ is $n/(n+1)$ times the length of $[v_i,b_i]$. Hence the distance from $b$ to $v_i$ is bounded by $n/(n+1)$ times the diameter of $[v_0,\ldots,v_n]$.

The significance of the factor $n/(n+1)$ is that by repeated barycentric subdivision we can produce simplices of arbitrarily small diameter since $(n/(n+1))^r$ approaches

%% --- END page-0129.tex ---
%% --- BEGIN page-0130.tex ---
0 as $r$ goes to infinity. It is important that the bound $n/(n+1)$ does not depend on the shape of the simplex since repeated barycentric subdivision produces simplices of many different shapes.

\textbf{(2) Barycentric Subdivision of Linear Chains.} The main part of the proof will be to construct a subdivision operator $S:C_n(X)\to C_n(X)$ and show this is chain homotopic to the identity map. First we will construct $S$ and the chain homotopy in a more restricted linear setting.

For a convex set $Y$ in some Euclidean space, the linear maps $\Delta^n\to Y$ generate a subgroup of $C_n(Y)$ that we denote $LC_n(Y)$, the linear chains. The boundary map $\partial:C_n(Y)\to C_{n-1}(Y)$ takes $LC_n(Y)$ to $LC_{n-1}(Y)$, so the linear chains form a subcomplex of the singular chain complex of $Y$. We can uniquely designate a linear map $\lambda:\Delta^n\to Y$ by $[w_0,\ldots,w_n]$ where $w_i$ is the image under $\lambda$ of the $i^{\text{th}}$ vertex of $\Delta^n$. To avoid having to make exceptions for $0$-simplices it will be convenient to augment the complex $LC(Y)$ by setting $LC_{-1}(Y) = \Z$ generated by the empty simplex $[\varnothing]$, with $\partial[w_0] = [\varnothing]$ for all $0$-simplices $[w_0]$.

Each point $b \in Y$ determines a homomorphism $b:LC_n(Y)\to LC_{n+1}(Y)$ defined on basis elements by $b([w_0,\ldots,w_n]) = [b,w_0,\ldots,w_n]$. Geometrically, the homomorphism $b$ can be regarded as a cone operator, sending a linear chain to the cone having the linear chain as the base of the cone and the point $b$ as the tip of the cone. Applying the usual formula for $\partial$, we obtain the relation $\partial b([w_0,\ldots,w_n]) = [w_0,\ldots,w_n] - b(\partial[w_0,\ldots,w_n])$. By linearity it follows that $\partial b(\alpha) = \alpha - b(\partial\alpha)$ for all $\alpha \in LC_n(Y)$. This expresses algebraically the geometric fact that the boundary of a cone consists of its base together with the cone on the boundary of its base. The relation $\partial b(\alpha) = \alpha - b(\partial\alpha)$ can be rewritten as $\partial b + b\partial = \mathbb{1}$, so $b$ is a chain homotopy between the identity map and the zero map on the augmented chain complex $LC(Y)$.

Now we define a subdivision homomorphism $S:LC_n(Y)\to LC_n(Y)$ by induction on $n$. Let $\lambda:\Delta^n\to Y$ be a generator of $LC_n(Y)$ and let $b_\lambda$ be the image of the barycenter of $\Delta^n$ under $\lambda$. Then the inductive formula for $S$ is $S(\lambda) = b_\lambda(S\partial\lambda)$ where $b_\lambda:LC_{n-1}(Y)\to LC_n(Y)$ is the cone operator defined in the preceding paragraph. The induction starts with $S([\varnothing]) = [\varnothing]$, so $S$ is the identity on $LC_{-1}(Y)$. It is also the identity on $LC_0(Y)$, since when $n=0$ the formula for $S$ becomes $S([w_0]) = w_0(S\partial[w_0]) = w_0(S([\varnothing])) = w_0([\varnothing]) = [w_0]$. When $\lambda$ is an embedding, with image a genuine $n$-simplex $[w_0,\ldots,w_n]$, then $S(\lambda)$ is the sum of the $n$-simplices in the barycentric subdivision of $[w_0,\ldots,w_n]$, with certain signs that could be computed explicitly. This is apparent by comparing the inductive definition of $S$ with the inductive definition of the barycentric subdivision of a simplex.

Let us check that the maps $S$ satisfy $\partial S = S\partial$, and hence give a chain map from the chain complex $LC(Y)$ to itself. Since $S = \mathbb{1}$ on $LC_0(Y)$ and $LC_{-1}(Y)$, we certainly have $\partial S = S\partial$ on $LC_0(Y)$. The result for larger $n$ is given by the following calculation, in which we omit some parentheses to unclutter the formulas:

%% --- END page-0130.tex ---
%% --- BEGIN ALIGNED PAGES 131-136 ---
% Source: Hatcher, "Algebraic Topology", Chapter 2 "Homology"
% (2002 Cambridge University Press edition), book pages 122--127.
% These pages are entirely within the "Exact Sequences and Excision" subsection.
% Content: chain homotopy formulas, barycentric subdivision, proof of excision theorem,
% good pairs, local homology, naturality of exact sequences.

We now establish formulas needed in the proof. For a singular $n$-chain $\Lambda$ we have
\begin{align*}
\partial S\Lambda &= \partial b_\lambda(S\partial\Lambda)\\
&= S\partial\Lambda - b_\lambda\partial(S\partial\Lambda) \quad \text{since } \partial b_\lambda = \mathbb{I} - b_\lambda\partial\\
&= S\partial\Lambda - b_\lambda S(\partial\partial\Lambda) \quad \text{since } \partial S(\partial\Lambda) = S\partial(\partial\Lambda) \text{ by induction on } n\\
&= S\partial\Lambda \quad \text{since } \partial\partial = 0
\end{align*}

We next build a chain homotopy $T:LC_n(Y)\to LC_{n+1}(Y)$ between $S$ and the identity, fitting into a diagram


We define $T$ on $LC_n(Y)$ inductively by setting $T = 0$ for $n = -1$ and letting $T\Lambda = b_\lambda(\Lambda - T\partial\Lambda)$ for $n \geq 0$. The geometric motivation for this formula is an inductively defined subdivision of $\Delta^n\times I$ obtained by joining all simplices in $\Delta^n\times\{0\}\cup\partial\Delta^n\times I$ to the barycenter of $\Delta^n\times\{1\}$, as indicated in the figure in the case $n = 2$. What $T$ actually does is take the image of this subdivision under the projection $\Delta^n\times I\to\Delta^n$.

The chain homotopy formula $\partial T + T\partial = \mathbb{I} - S$ is trivial on $LC_{-1}(Y)$ where $T = 0$ and $S = \mathbb{I}$. Verifying the formula on $LC_n(Y)$ with $n \geq 0$ is done by the calculation
\begin{align*}
\partial T\Lambda &= \partial b_\lambda(\Lambda - T\partial\Lambda)\\
&= \Lambda - T\partial\Lambda - b_\lambda\partial(\Lambda - T\partial\Lambda) \quad \text{since } \partial b_\lambda = \mathbb{I} - b_\lambda\partial\\
&= \Lambda - T\partial\Lambda - b_\lambda[\partial\Lambda - \partial T(\partial\Lambda)]\\
&= \Lambda - T\partial\Lambda - b_\lambda[S(\partial\Lambda) + T\partial(\partial\Lambda)] \quad \text{by induction on } n\\
&= \Lambda - T\partial\Lambda - S\Lambda \quad \text{since } \partial\partial = 0 \text{ and } S\Lambda = b_\lambda(S\partial\Lambda)
\end{align*}

Now we can discard the group $LC_{-1}(Y)$ and the relation $\partial T + T\partial = \mathbb{I} - S$ still holds since $T$ was zero on $LC_{-1}(Y)$.

\begin{description}
\item[(3) Barycentric Subdivision of General Chains.]
\end{description}

Define $S:C_n(X)\to C_n(X)$ by setting $S\sigma = \sigma_* S\Delta^n$ for a singular $n$-simplex $\sigma:\Delta^n\to X$. Since $S\Delta^n$ is the sum of the $n$-simplices in the barycentric subdivision of $\Delta^n$, with certain signs, $S\sigma$ is the corresponding signed sum of the restrictions of $\sigma$ to the $n$-simplices of the barycentric subdivision of $\Delta^n$. The operator $S$ is a chain map since
\begin{align*}
\partial S\sigma &= \partial\sigma_* S\Delta^n = \sigma_*\partial S\Delta^n = \sigma_*S\partial\Delta^n\\
&= \sigma_*\left(S\left(\sum_i(-1)^i\Delta^n_i\right)\right) \quad \text{where } \Delta^n_i \text{ is the } i\text{-th face of } \Delta^n\\
&= \sum_i(-1)^i\sigma_*S\Delta^n_i\\
&= \sum_i(-1)^iS(\sigma|\Delta^n_i)\\
&= S\left(\sum_i(-1)^i\sigma|\Delta^n_i\right) = S(\partial\sigma)
\end{align*}

In similar fashion we define $T:C_n(X)\to C_{n+1}(X)$ by $T\sigma = \sigma_* T\Delta^n$, and this gives a chain homotopy between $S$ and the identity, since the formula $\partial T + T\partial = \mathbb{I} - S$ holds by the calculation
\begin{align*}
\partial T\sigma &= \partial\sigma_* T\Delta^n = \sigma_*\partial T\Delta^n = \sigma_*(\Delta^n - S\Delta^n - T\partial\Delta^n)\\
&= \sigma - S\sigma - \sigma_* T\partial\Delta^n = \sigma - S\sigma - T(\partial\sigma)
\end{align*}
where the last equality follows just as in the previous displayed calculation, with $S$ replaced by $T$.

\begin{description}
\item[(4) Iterated Barycentric Subdivision.]
\end{description}

A chain homotopy between $\mathbb{I}$ and the iterate $S^m$ is given by the operator $D_m = \sum_{0\leq i<m} TS^i$ since
\[\partial D_m + D_m\partial = \sum_{0\leq i<m}(\partial TS^i + TS^i\partial) = \sum_{0\leq i<m}(\partial TS^i + T\partial S^i) = \sum_{0\leq i<m}(\partial T + T\partial)S^i = \sum_{0\leq i<m}(\mathbb{I} - S)S^i = \sum_{0\leq i<m}(S^i - S^{i+1}) = \mathbb{I} - S^m\]

For each singular $n$-simplex $\sigma:\Delta^n\to X$ there exists an $m$ such that $S^m(\sigma)$ lies in $C_n^{\mathcal{U}}(X)$ since the diameter of the simplices of $S^m(\Delta^n)$ will be less than a Lebesgue number of the cover of $\Delta^n$ by the open sets $\sigma^{-1}(\text{int } U_j)$ if $m$ is large enough. (Recall that a Lebesgue number for an open cover of a compact metric space is a number $\varepsilon > 0$ such that every set of diameter less than $\varepsilon$ lies in some set of the cover; such a number exists by an elementary compactness argument.) We cannot expect the same number $m$ to work for all $\sigma$'s, so let us define $m(\sigma)$ to be the smallest $m$ such that $S^m\sigma$ is in $C_n^{\mathcal{U}}(X)$.

We now define $D:C_n(X)\to C_{n+1}(X)$ by setting $D\sigma = D_{m(\sigma)}\sigma$ for each singular $n$-simplex $\sigma:\Delta^n\to X$. For this $D$ we would like to find a chain map $\rho:C_n(X)\to C_n(X)$ with image in $C_n^{\mathcal{U}}(X)$ satisfying the chain homotopy equation
\[(*) \quad \partial D + D\partial = \mathbb{I} - \rho\]

A quick way to do this is simply to regard this equation as defining $\rho$, so we let $\rho = \mathbb{I} - \partial D - D\partial$. It follows easily that $\rho$ is a chain map since
\[\partial\rho(\sigma) = \partial\sigma - \partial^2D\sigma - \partial D\partial\sigma = \partial\sigma - \partial D\partial\sigma\]
and
\[\rho(\partial\sigma) = \partial\sigma - \partial D\partial\sigma - D\partial^2\sigma = \partial\sigma - \partial D\partial\sigma\]

To check that $\rho$ takes $C_n(X)$ to $C_n^{\mathcal{U}}(X)$ we compute $\rho(\sigma)$ more explicitly:
\begin{align*}
\rho(\sigma) &= \sigma - \partial D\sigma - D(\partial\sigma)\\
&= \sigma - \partial D_{m(\sigma)}\sigma - D(\partial\sigma)\\
&= S^{m(\sigma)}\sigma + D_{m(\sigma)}(\partial\sigma) - D(\partial\sigma) \quad \text{since } \partial D_m + D_m\partial = \mathbb{I} - S^m
\end{align*}

The term $S^{m(\sigma)}\sigma$ lies in $C_n^{\mathcal{U}}(X)$ by the definition of $m(\sigma)$. The remaining terms $D_{m(\sigma)}(\partial\sigma) - D(\partial\sigma)$ are linear combinations of terms $D_{m(\sigma)}(\sigma_j) - D_{m(\sigma_j)}(\sigma_j)$ for $\sigma_j$ the restriction of $\sigma$ to a face of $\Delta^n$, so $m(\sigma_j) \leq m(\sigma)$ and hence the difference $D_{m(\sigma)}(\sigma_j) - D_{m(\sigma_j)}(\sigma_j)$ consists of terms $TS^i(\sigma_j)$ with $i \geq m(\sigma_j)$, and these terms lie in $C_n^{\mathcal{U}}(X)$ since $T$ takes $C_{n-1}^{\mathcal{U}}(X)$ to $C_n^{\mathcal{U}}(X)$.

Viewing $\rho$ as a chain map $C_n(X)\to C_n^{\mathcal{U}}(X)$, the equation $(*)$ says that $\partial D + D\partial = \mathbb{I} - \rho$ for $\iota:C_n^{\mathcal{U}}(X)\to C_n(X)$ the inclusion. Furthermore, $\rho\iota = \mathbb{I}$ since $D$ is identically zero on $C_n^{\mathcal{U}}(X)$, as $m(\sigma) = 0$ if $\sigma$ is in $C_n^{\mathcal{U}}(X)$, hence the summation defining $D\sigma$ is empty. Thus we have shown that $\rho$ is a chain homotopy inverse for $\iota$.
\end{proof}
\begin{proof}[the Excision Theorem]
 We prove the second version, involving a decomposition $X = A \cup B$. For the cover $\mathcal{U} = \{A, B\}$ we introduce the suggestive notation $C_n(A + B)$ for $C_n^{\mathcal{U}}(X)$, the sums of chains in $A$ and chains in $B$. At the end of the preceding proof we had formulas $\partial D + D\partial = \mathbb{I} - \rho$ and $\rho\iota = \mathbb{I}$. All the maps appearing in these formulas take chains in $A$ to chains in $A$, so they induce quotient maps when we factor out chains in $A$. These quotient maps automatically satisfy the same two formulas, so the inclusion $C_n(A + B)/C_n(A) \to C_n(X)/C_n(A)$ induces an isomorphism on homology. The map $C_n(B)/C_n(A \cap B)\to C_n(A + B)/C_n(A)$ induced by inclusion is obviously an isomorphism since both quotient groups are free with basis the singular $n$-simplices in $B$ that do not lie in $A$. Hence we obtain the desired isomorphism $H_n(B, A \cap B) = H_n(X, A)$ induced by inclusion.
\end{proof}
All that remains in the proof of Theorem 2.13 is to replace relative homology groups with absolute homology groups. This is achieved by the following result.

\begin{proposition}[Good pairs and absolute homology]
\label{prop:good-pairs-absolute-homology}
\dcref{ch2:2.22}
\group{Singular Homology}
\uses{thm:excision-general,def:deformation-retraction}
For good pairs $(X, A)$, the quotient map $q:(X,A)\to(X/A, A/A)$ induces isomorphisms
\[q_*:H_n(X,A)\to H_n(X/A, A/A) \simeq \tilde{H}_n(X/A)\]
for all $n$.
\end{proposition}

\begin{proof}
Let $V$ be a neighborhood of $A$ in $X$ that deformation retracts onto $A$. We have a commutative diagram


The upper left horizontal map is an isomorphism since in the long exact sequence of the triple $(X, V, A)$ the groups $H_n(V, A)$ are zero for all $n$, because a deformation retraction of $V$ onto $A$ gives a homotopy equivalence of pairs $(V, A) = (A, A)$, and $H_n(A,A) = 0$. The deformation retraction of $V$ onto $A$ induces a deformation retraction of $V/A$ onto $A/A$, so the same argument shows that the lower left horizontal map is an isomorphism as well. The other two horizontal maps are isomorphisms directly from excision. The right-hand vertical map $q_*$ is an isomorphism since $q$ restricts to a homeomorphism on the complement of $A$. From the commutativity of the diagram it follows that the left-hand $q_*$ is an isomorphism.
\end{proof}
This proposition shows that relative homology can be expressed as reduced absolute homology in the case of good pairs $(X, A)$, but in fact there is a way of doing this for arbitrary pairs. Consider the space $X \cup CA$ where $CA$ is the cone $(A\times I)/(A\times\{0\})$ whose base $A\times\{1\}$ we identify with $A \subset X$. Using terminology introduced in Chapter 0, $X\cup CA$ can also be described as the mapping cone of the inclusion $A \to X$. The assertion is that $H_n(X, A)$ is isomorphic to $\tilde{H}_n(X \cup CA)$ for all $n$ via the sequence of isomorphisms
\[\tilde{H}_n(X \cup CA) \simeq H_n(X \cup CA, CA) \simeq H_n(X \cup CA - \{p\}, CA - \{p\}) \simeq H_n(X, A)\]
where $p \in CA$ is the tip of the cone. The first isomorphism comes from the exact sequence of the pair, using the fact that $CA$ is contractible. The second isomorphism is excision, and the third comes from a deformation retraction of $CA - \{p\}$ onto $A$.

% NEEDS-USES: homology of spheres, homology of disks, long exact sequence of pair
\begin{example}[Explicit cycles generating homology of spheres and disks]
\label{ex:cycles-spheres-disks}
\dcref{ch2:2.23}
\group{Singular Homology}
\uses{cor:homology-spheres}
Let us find explicit cycles representing generators of the infinite cyclic groups $H_n(D^n,\partial D^n)$ and $\tilde{H}_n(S^n)$. Replacing $(D^n, \partial D^n)$ by the equivalent pair $(\Delta^n, \partial\Delta^n)$, we will show by induction on $n$ that the identity map $\iota_n:\Delta^n\to\Delta^n$, viewed as a singular $n$-simplex, is a cycle generating $H_n(\Delta^n, \partial\Delta^n)$. That it is a cycle is clear since we are considering relative homology. When $n = 0$ it certainly represents a generator. For the induction step, let $\Lambda \subset \Delta^n$ be the union of all but one of the $(n-1)$-dimensional faces of $\Delta^n$. Then we claim there are isomorphisms
\[H_n(\Delta^n, \partial\Delta^n) \xrightarrow{\simeq} H_{n-1}(\partial\Delta^n, \Lambda) \xrightarrow{\simeq} H_{n-1}(\Delta^{n-1}, \partial\Delta^{n-1})\]

The first isomorphism is a boundary map in the long exact sequence of the triple $(\Delta^n, \partial\Delta^n, \Lambda)$, whose third terms $H_i(\Delta^n, \Lambda)$ are zero since $\Delta^n$ deformation retracts onto $\Lambda$, hence $(\Delta^n, \Lambda) = (\Lambda, \Lambda)$. The second isomorphism is induced by the inclusion $i:\Delta^{n-1}\to\partial\Delta^n$ as the face not contained in $\Lambda$. When $n = 1$, $i$ induces an isomorphism on relative homology since this is true already at the chain level. When $n > 1$, $\partial\Delta^{n-1}$ is nonempty so we are dealing with good pairs and $i$ induces a homeomorphism of quotients $\Delta^{n-1}/\partial\Delta^{n-1} = \partial\Delta^n/\Lambda$. The induction step then follows since the cycle $i_n$ is sent under the first isomorphism to the cycle $\partial i_n$, which equals $\pm i_{n-1}$ in $C_{n-1}(\partial\Delta^n, \Lambda)$.

To find a cycle generating $\tilde{H}_n(S^n)$ let us regard $S^n$ as two $n$-simplices $\Delta_1^n$ and $\Delta_2^n$ with their boundaries identified in the obvious way, preserving the ordering of vertices. The difference $\Delta_1^n - \Delta_2^n$, viewed as a singular $n$-chain, is then a cycle, and we claim it represents a generator of $\tilde{H}_n(S^n)$. To see this, consider the isomorphisms
\[\tilde{H}_n(S^n) \xrightarrow{\simeq} H_n(S^n, \Delta_2^n) \xrightarrow{\simeq} H_n(\Delta_1^n, \partial\Delta_1^n)\]

where the first isomorphism comes from the long exact sequence of the pair $(S^n, \Delta_2^n)$ and the second isomorphism is justified in the nontrivial cases $n > 0$ by passing to quotients as before. Under these isomorphisms the cycle $\Delta_1^n - \Delta_2^n$ in the first group corresponds to the cycle $\Delta_1^n$ in the third group, which represents a generator of this group as we have seen, so $\Delta_1^n - \Delta_2^n$ represents a generator of $\tilde{H}_n(S^n)$.
\end{example}

The preceding proposition implies that the excision property holds also for subcomplexes of CW complexes:

% NEEDS-USES: CW pairs, excision
\begin{corollary}[Excision for CW subcomplexes]
\label{cor:excision-cw-subcomplexes}
\dcref{ch2:2.24}
\group{Singular Homology}
\uses{thm:excision-general,def:cw-pair,prop:cw-pairs-have-hep}
If the CW complex $X$ is the union of subcomplexes $A$ and $B$, then the inclusion $(B, A \cap B) \to (X, A)$ induces isomorphisms
\[H_n(B, A \cap B) \to H_n(X, A)\]
for all $n$.
\end{corollary}

\begin{proof}
Since CW pairs are good, \cref{prop:good-pairs-absolute-homology} allows us to pass to the quotient spaces $B/(A \cap B)$ and $X/A$ which are homeomorphic, assuming we are not in the trivial case $A \cap B = \emptyset$.
\end{proof}
Here is another application of the preceding proposition:

% NEEDS-USES: reduced homology, basepoint, good pairs
\begin{corollary}[Wedge sums and direct sum of homology]
\label{cor:wedge-sum-direct-sum-homology}
\dcref{ch2:2.25}
\group{Singular Homology}
\uses{prop:good-pairs-absolute-homology}
For a wedge sum $\bigvee_\alpha X_\alpha$, the inclusions $i_\alpha: X_\alpha \to \bigvee_\alpha X_\alpha$ induce an isomorphism
\[\bigoplus_\alpha i_{\alpha*}:\bigoplus_\alpha \tilde{H}_n(X_\alpha) \to \tilde{H}_n\left(\bigvee_\alpha X_\alpha\right)\]
provided that the wedge sum is formed at basepoints $x_\alpha \in X_\alpha$ such that the pairs $(X_\alpha, x_\alpha)$ are good.
\end{corollary}

\begin{proof}
Since reduced homology is the same as homology relative to a basepoint, this follows from \cref{prop:good-pairs-absolute-homology} by taking $(X, A) = (\bigcup_\alpha X_\alpha, \bigcup_\alpha \{x_\alpha\})$.
\end{proof}
Here is an application of the machinery we have developed, a classical result of Brouwer from around 1910 known as `invariance of dimension', which says in particular that $\R^m$ is not homeomorphic to $\R^n$ if $m \ne n$.

% NEEDS-USES: homology of spheres, excision, long exact sequence
\begin{theorem}[Invariance of dimension]
\label{thm:invariance-of-dimension}
\dcref{ch2:2.26}
\group{Singular Homology}
\uses{prop:good-pairs-absolute-homology}
If nonempty open sets $U \subset \R^m$ and $V \subset \R^n$ are homeomorphic, then $m = n$.
\end{theorem}

\begin{proof}
For $x \in U$ we have $H_k(U, U - \{x\}) \simeq H_k(\R^m, \R^m - \{x\})$ by excision. From the long exact sequence for the pair $(\R^m, \R^m - \{x\})$ we get $H_k(\R^m, \R^m - \{x\}) \simeq \tilde{H}_{k-1}(\R^m - \{x\})$. Since $\R^m - \{x\}$ deformation retracts onto a sphere $S^{m-1}$, we conclude that $H_k(U, U - \{x\})$ is $\Z$ for $k = m$ and 0 otherwise. By the same reasoning, $H_k(V, V - \{y\})$ is $\Z$ for $k = n$ and 0 otherwise. Since a homeomorphism $h:U\to V$ induces isomorphisms $H_k(U, U - \{x\})\to H_k(V, V - \{h(x)\})$ for all $k$, we must have $m = n$.
\end{proof}
Generalizing the idea of this proof, the \textbf{local homology groups} of a space $X$ at a point $x \in X$ are defined to be the groups $H_n(X, X - \{x\})$. For any open neighborhood $U$ of $x$, excision gives isomorphisms $H_n(X, X - \{x\}) \simeq H_n(U, U - \{x\})$ assuming points are closed in $X$, and thus the groups $H_n(X, X - \{x\})$ depend only on the local topology of $X$ near $x$. A homeomorphism $f:X\to Y$ must induce isomorphisms $H_n(X, X - \{x\}) \simeq H_n(Y, Y - \{f(x)\})$ for all $x$ and $n$, so the local homology groups can be used to tell when spaces are not locally homeomorphic at certain points, as in the preceding proof. The exercises give some further examples of this.

\subsection*{Naturality}

The exact sequences we have been constructing have an extra property that will become important later at key points in many arguments, though at first glance this property may seem just an idle technicality, not very interesting. We shall discuss the property now rather than interrupting later arguments to check it when it is needed, but the reader may prefer to postpone a careful reading of this discussion.

The property is called \textbf{naturality}. For example, to say that the long exact sequence of a pair is natural means that for a map $f:(X,A)\to(Y,B)$, the diagram

is commutative. Commutativity of the squares involving $i_*$ and $j_*$ follows from the obvious commutativity of the corresponding squares of chain groups, with $C_n$ in place of $H_n$. For the other square, when we defined induced homomorphisms we saw that $f_*\partial = \partial f_*$ at the chain level. Then for a class $[\alpha] \in H_n(X,A)$ represented by a relative cycle $\alpha$, we have
\[f_*\partial[\alpha] = f_*[\partial\alpha] = [f_*(\partial\alpha)] = [\partial f_*(\alpha)] = \partial[f_*\alpha] = \partial f_*[\alpha].\]

Alternatively, we could appeal to the general algebraic fact that the long exact sequence of homology groups associated to a short exact sequence of chain complexes is natural: For a commutative diagram of short exact sequences of chain complexes

the induced diagram

is commutative. Commutativity of the first two squares is obvious since $\beta i = i'\alpha$ implies $\beta_* i_* = i'_*\alpha_*$ and $\gamma j = j'\beta$ implies $\gamma_* j_* = j'_*\beta_*$. For the third square, recall that the map $\partial:H_n(C)\to H_{n-1}(A)$ was defined by $\partial[c] = [a]$ where $c = j(b)$ and $i(a) = \partial b$. Then $\partial(\gamma_*[c]) = [\alpha(a)]$ since $\gamma_*(c) = \gamma(j(b)) = j'(\beta(b))$ and $i'(\alpha(a)) = \beta(i(a)) = \beta(\partial b) = \partial(\beta(b))$. Hence $\partial\gamma_*[c] = \alpha_*[a] = \alpha_*\partial[c]$.

%% --- END ALIGNED PAGES 131-136 ---
%% --- BEGIN page-0137.tex ---
This algebraic fact also implies naturally of the long exact sequence of a triple
and the long exact sequence of reduced homology of a pair.

Finally, there is the naturality of the long exact sequence in Theorem 2.13, that
is, commutativity of the diagram
\[
\cdots \to \tilde H_n(A) \xrightarrow{i_*} \tilde H_n(X) \xrightarrow{q_*} \tilde H_n(X/A) \xrightarrow{\partial} \tilde H_{n-1}(A) \to \cdots
\]
with vertical maps $f_*$ on each of the four groups and $\bar f: X/A \to Y/B$ induced
by $f$. The first two squares commute since $fi = if$ and $\bar f q = qf$. The third square
expands into
\[
\begin{array}{ccccccc}
\tilde H_n(X/A) & \xrightarrow{j_*} & H_n(X/A, A/A) & \xrightarrow{q_*} & H_n(X,A) & \xrightarrow{\partial} & \tilde H_{n-1}(A) \\
\bar f_* \downarrow & & \bar f_* \downarrow & & f_* \downarrow & & f_* \downarrow \\
\tilde H_n(Y/B) & \xrightarrow{j_*} & H_n(Y/B, B/B) & \xrightarrow{q_*} & H_n(Y,B) & \xrightarrow{\partial} & \tilde H_{n-1}(B)
\end{array}
\]
We have already shown commutativity of the first and third squares, and the second
square commutes since $\bar f q = qf$.

\subsection*{The Equivalence of Simplicial and Singular Homology}

We can use the preceding results to show that the simplicial and singular homol-
ogy groups of $\Delta$-complexes are always isomorphic. For the proof it will be convenient
to consider the relative case as well, so let $X$ be a $\Delta$-complex with $A \subset X$ a sub-
complex. Thus $A$ is the $\Delta$-complex formed by any union of simplices of $X$. Relative
groups $H_n^\Delta(X,A)$ can be defined in the same way as for singular homology, via relative
chains $\Delta_n(X,A) = \Delta_n(X)/\Delta_n(A)$, and this yields a long exact sequence of simplicial
homology groups for the pair $(X,A)$ by the same algebraic argument as for singular
homology. There is a canonical homomorphism $H_n^\Delta(X,A) \to H_n(X,A)$ induced by the
chain map $\Delta_n(X,A) \to C_n(X,A)$ sending each $n$-simplex of $X$ to its characteristic
map $\sigma: \Delta^n \to X$. The possibility $A = \emptyset$ is not excluded, in which case the relative
groups reduce to absolute groups.

\begin{theorem}[Isomorphism of simplicial and singular homology]
\label{thm:simplicial-singular-homology-isomorphic}
\dcref{ch2:2.27}
\group{Singular Homology}
\uses{lem:five-lemma}
The homomorphisms $H_n^\Delta(X,A) \to H_n(X,A)$ are isomorphisms for
all $n$ and all $\Delta$-complex pairs $(X,A)$.
\end{theorem}
\begin{proof}
 First we do the case that $X$ is finite-dimensional and $A$ is empty. For $X^k$
the $k$-skeleton of $X$, consisting of all simplices of dimension $k$ or less, we have a
commutative diagram of exact sequences:
\[
\begin{array}{ccccc}
H_{n+1}^\Delta(X^k, X^{k-1}) & \to H_n^\Delta(X^{k-1}) & \to H_n^\Delta(X^k) & \to H_n^\Delta(X^k, X^{k-1}) & \to H_{n-1}^\Delta(X^{k-1}) \\
\downarrow & \downarrow & \downarrow & \downarrow & \downarrow \\
H_{n+1}(X^k, X^{k-1}) & \to H_n(X^{k-1}) & \to H_n(X^k) & \to H_n(X^k, X^{k-1}) & \to H_{n-1}(X^{k-1})
\end{array}
\]

%% --- END page-0137.tex ---
%% --- BEGIN page-0138.tex ---
Let us first show that the first and fourth vertical maps are isomorphisms for all $n$.
The simplicial chain group $\Delta_n(X^k, X^{k-1})$ is zero for $n \ne k$, and is free abelian with
basis the $k$-simplices of $X$ when $n = k$. Hence $H_n^\Delta(X^k, X^{k-1})$ has exactly the same
description. The corresponding singular homology groups $H_n(X^k, X^{k-1})$ can be com-
puted by considering the map $\Phi: \sqcup_\alpha(\Delta_k^\alpha, \partial \Delta_k^\alpha) \to (X^k, X^{k-1})$ formed by the character-
istic maps $\Delta^k \to X$ for all the $k$-simplices of $X$. Since $\Phi$ induces a homeomorphism
of quotient spaces $\sqcup_\alpha \Delta_k^\alpha / \sqcup_\alpha \partial \Delta_k^\alpha \cong X^k/X^{k-1}$, it induces isomorphisms on all singu-
lar homology groups. Thus $H_n(X^k, X^{k-1})$ is zero for $n \ne k$, while for $n = k$ this
group is free abelian with basis represented by the relative cycles given by the char-
acteristic maps of all the $k$-simplices of $X$, in view of the fact that $H_k(\Delta^k, \partial\Delta^k)$ is
generated by the identity map $\Delta^k \to \Delta^k$, as we showed in Example 2.23. Therefore the
map $H_n^\Delta(X^k, X^{k-1}) \to H_k(X^k, X^{k-1})$ is an isomorphism.

By induction on $k$ we may assume the second and fifth vertical maps in the pre-
ceding diagram are isomorphisms as well. The following frequently quoted basic alge-
braic lemma will then imply that the middle vertical map is an isomorphism, finishing
the proof when $X$ is finite-dimensional and $A = \emptyset$. \end{proof}

\begin{lemma}[The Five-Lemma]
\label{lem:five-lemma}
\group{Singular Homology}
\uses{}
In a commutative diagram
\[
\begin{array}{ccccccccc}
A & \xrightarrow{i} & B & \xrightarrow{j} & C & \xrightarrow{k} & D & \xrightarrow{\ell} & E \\
\alpha \downarrow & & \beta \downarrow & & \gamma \downarrow & & \delta \downarrow & & \varepsilon \downarrow \\
A' & \xrightarrow{i'} & B' & \xrightarrow{j'} & C' & \xrightarrow{k'} & D' & \xrightarrow{\ell'} & E'
\end{array}
\]
of abelian groups, if the two rows are exact and $\alpha, \beta, \delta,$ and $\varepsilon$ are isomorphisms, then $\gamma$ is an isomorphism also.
\end{lemma}

\begin{proof}
It suffices to show:
\begin{enumerate}
\item[(a)] $\gamma$ is surjective if $\beta$ and $\delta$ are surjective and $\varepsilon$ is injective.
\item[(b)] $\gamma$ is injective if $\beta$ and $\delta$ are injective and $\alpha$ is surjective.
\end{enumerate}

The proofs of these two statements are straightforward diagram chasing. There is really no choice about how the argument can proceed, and it would be a good exercise for the reader to close the book now and reconstruct the proofs without looking.

To prove (a), start with an element $c' \in C'$. Then $k'(c') = \delta(d)$ for some $d \in D$ since $\delta$ is surjective. Since $\varepsilon$ is injective and $\varepsilon \ell(d) = \ell' \delta(d) = \ell' k'(c') = 0$, we deduce that $\ell(d) = 0$, hence $d = k(c)$ for some $c \in C$ by exactness of the upper row. The difference $c' - \gamma(c)$ maps to $0$ under $k'$ since $k'(c') - k'\gamma(c) = k'(c') - \delta k(c) = k'(c') - \delta(d) = 0$. Therefore $c' - \gamma(c) = j'(b')$ for some $b' \in B'$ by exactness. Since $\beta$ is surjective, $b' = \beta(b)$ for some $b \in B$, and then $\gamma(c + j(b)) = \gamma(c) + \gamma j(b) = \gamma(c) + j'(b) = \gamma(c) + j'(\beta(b)) = \gamma(c) + j'(\beta(b)) = c'$, showing that $\gamma$ is surjective.

To prove (b), suppose that $\gamma(c) = 0$. Since $\delta$ is injective, $\delta k(c) = k' \gamma(c) = 0$ implies $k(c) = 0$, so $c = j(b)$ for some $b \in B$. The element $\beta(b)$ satisfies $j' \beta(b) = \gamma j(b) = \gamma(c) = 0$, so $\beta(b) = i'(a')$ for some $a' \in A'$. Since $\alpha$ is surjective, $a' = \alpha(a)$ for some $a \in A$. Since $\beta$ is injective, $\beta(i(a) - b) = \beta i(a) - \beta(b) = i'(\alpha(a)) - \beta(b) = i'(\alpha(a)) - \beta(b) = 0$ implies $i(a) - b = 0$. Thus $b = i(a)$, and hence $c = j(b) = ji(a) = 0$ since $ji = 0$. This shows $\gamma$ has trivial kernel.
\end{proof}

%% --- END page-0138.tex ---
%% --- BEGIN page-0139.tex ---
\begin{proof}[Proof of Theorem 2.27 (continued)]
Returning to the proof of the theorem, we next consider the case that $X$ is infinite-dimensional, where we will use the following fact: A compact set in $X$ can meet only finitely many open simplices of $X$, that is, simplices with their proper faces deleted. This is a general fact about CW complexes proved in the Appendix, but here is a direct proof for $\Delta$-complexes. If a compact set $C$ intersected infinitely many open simplices, it would contain an infinite sequence of points $x_i$ each lying in a different open simplex. Then the sets $U_i = X - \bigcup_{j \ne i} \{x_j\}$, which are open since their preimages under the characteristic maps of all the simplices are clearly open, form an open cover of $C$ with no finite subcover.

This can be applied to show the map $H_n^\Delta(X) \to H_n(X)$ is surjective. Represent a given element of $H_n(X)$ by a singular $n$-cycle $z$. This is a linear combination of finitely many singular simplices with compact images, meeting only finitely many open simplices of $X$, hence contained in $X^k$ for some $k$. We have shown that $H_n^\Delta(X^k) \to H_n(X^k)$ is an isomorphism, in particular surjective, so $z$ is homologous in $X^k$ (hence in $X$) to a simplicial cycle. This gives surjectivity. Injectivity is similar: If a simplicial $n$-cycle $z$ is the boundary of a singular chain in $X$, this chain has compact image and hence must lie in some $X^k$, so $z$ represents an element of the kernel of $H_n^\Delta(X^k) \to H_n(X^k)$. But we know this map is injective, so $z$ is a simplicial boundary in $X^k$, and therefore in $X$.

It remains to do the case of arbitrary $X$ with $A \ne \emptyset$, but this follows from the absolute case by applying the five-lemma to the canonical map from the long exact sequence of simplicial homology groups for the pair $(X, A)$ to the corresponding long exact sequence of singular homology groups.
\end{proof}

We can deduce from this theorem that $H_n(X)$ is finitely generated whenever $X$
is a $\Delta$-complex with finitely many $n$-simplices, since in this case the simplicial chain
group $\Delta_n(X)$ is finitely generated, hence also its subgroup of cycles and therefore
also the latter group's quotient $H_n^\Delta(X)$. If we write $H_n(X)$ as the direct sum of cyclic
groups, then the number of $\Z$ summands is known traditionally as the $n^{\text{th}}$ \textbf{Betti
number} of $X$, and integers specifying the orders of the finite cyclic summands are
called \textbf{torsion coefficients}.

It is a curious historical fact that homology was not thought of originally as a
sequence of groups, but rather as Betti numbers and torsion coefficients. One can
after all compute Betti numbers and torsion coefficients from the simplicial boundary
maps without actually mentioning homology groups. This computational viewpoint,
with homology being numbers rather than groups, prevailed from when Poincar\'e first
started serious work on homology around 1900, up until the 1920s when the more
abstract viewpoint of groups entered the picture. During this period `homology' meant
primarily `simplicial homology', and it was another 20 years before the shift to singular
homology was complete, with the final definition of singular homology emerging only

%% --- END page-0139.tex ---
%% --- BEGIN page-0140.tex ---
in a 1944 paper of Eilenberg, after contributions from quite a few others, particularly
Alexander and Lefschertz. Within the next few years the rest of the basic structure
of homology theory as we have presented it fell into place, and the first definitive
treatment appeared in the classic book [Eilenberg \& Steenrod 1952].

\subsection*{Exercises}

\begin{enumerate}

\item What familiar space is the quotient $\Delta$-complex of a 2-simplex $[v_0, v_1, v_2]$ obtained
by identifying the edges $[v_0, v_1]$ and $[v_1, v_2]$, preserving the ordering of vertices?

\item Show that the $\Delta$-complex obtained from $\Delta^3$ by performing the order-preserving
edge identifications $[v_0, v_1] \sim [v_1, v_3]$ and $[v_0, v_2] \sim [v_2, v_3]$ deformation retracts
onto a Klein bottle. Also, find other pairs of identifications of edges that produce
$\Delta$-complexes deformation retracting onto a torus, a 2-sphere, and $\RP^2$.

\item Construct a $\Delta$-complex structure on $\RP^n$ as a quotient of a $\Delta$-complex structure
on $S^n$ having vertices the two vectors of length 1 along each coordinate axis in $\R^{n+1}$.

\item Compute the simplicial homology groups of the triangular parachute obtained from
$\Delta^2$ by identifying its three vertices to a single point.

\item Compute the simplicial homology groups of the Klein bottle using the $\Delta$-complex
structure described at the beginning of this section.

\item Compute the simplicial homology groups of the $\Delta$-complex obtained from $n + 1$
2-simplices $\Delta_0^2, \ldots, \Delta_n^2$ by identifying all three edges of $\Delta_0^2$ to a single edge, and for
$i > 0$ identifying the edges $[v_0, v_1]$ and $[v_1, v_2]$ of $\Delta_i^2$ to a single edge and the edge
$[v_0, v_2]$ to the edge $[v_0, v_1]$ of $\Delta_{i-1}^2$.

\item Find a way of identifying pairs of faces of $\Delta^3$ to produce a $\Delta$-complex structure
on $S^3$ having a single 3-simplex, and compute the simplicial homology groups of this
$\Delta$-complex.

\item Construct a 3-dimensional $\Delta$-complex $X$ from $n$ tetrahedra $T_1, \ldots, T_n$ by the following two steps. First arrange the
tetrahedra in a cyclic pattern as in the figure, so that each $T_i$
shares a common vertical face with its two neighbors $T_{i-1}$
and $T_{i+1}$, subscripts being taken mod $n$. Then identify the
bottom face of $T_i$ with the top face of $T_{i+1}$ for each $i$. Show the simplicial homology
groups of $X$ in dimensions $0, 1, 2, 3$ are $\Z, \Z_n, 0, \Z$, respectively. [The space $X$ is
an example of a lens space; see Example 2.43 for the general case.]

\item Compute the homology groups of the $\Delta$-complex $X$ obtained from $\Delta^n$ by identi-
fying all faces of the same dimension. Thus $X$ has a single $k$-simplex for each $k \le n$.

\item (a) Show the quotient space of a finite collection of disjoint 2-simplices obtained
by identifying pairs of edges is always a surface, locally homeomorphic to $\R^2$.
(b) Show the edges can always be oriented so as to define a $\Delta$-complex structure on
the quotient surface. [This is more difficult.]

%% --- END page-0140.tex ---
%% --- BEGIN page-0141.tex ---
\item Show that if $A$ is a retract of $X$ then the map $H_n(A) \to H_n(X)$ induced by the
inclusion $A \subset X$ is injective.

\item Show that chain homotopy of chain maps is an equivalence relation.

\item Verify that $f \simeq g$ implies $f_* = g_*$ for induced homomorphisms of reduced
homology groups.

\item Determine whether there exists a short exact sequence $0 \to \Z_4 \to \Z_8 \oplus \Z_2 \to \Z_4 \to 0$.
More generally, determine which abelian groups $A$ fit into a short exact sequence
$0 \to \Z_{pm} \to A \to \Z_{pn} \to 0$ with $p$ prime. What about the case of short exact sequences
$0 \to \Z \to A \to \Z_n \to 0$?

\item For an exact sequence $A \to B \to C \to D \to E$ show that $C = 0$ iff the map $A \to B$
is surjective and $D \to E$ is injective. Hence for a pair of spaces $(X, A)$, the inclusion
$A \hookrightarrow X$ induces isomorphisms on all homology groups iff $H_n(X, A) = 0$ for all $n$.

\item (a) Show that $H_0(X, A) = 0$ iff $A$ meets each path-component of $X$.
(b) Show that $H_1(X, A) = 0$ iff $H_1(A) \to H_1(X)$ is surjective and each path-component
of $X$ contains at most one path-component of $A$.

\item (a) Compute the homology groups $H_n(X, A)$ when $X$ is $S^2$ or $S^1 \times S^1$ and $A$ is a
finite set of points in $X$.
(b) Compute the groups $H_n(X, A)$ and $H_n(X, B)$ for $X$
a closed orientable surface of genus two with $A$ and $B$
the circles shown. [What are $X/A$ and $X/B$?]

\item Show that for the subspace $\Q \subset \R$, the relative homology group $H_1(\R, \Q)$ is free
abelian and find a basis.

\item Compute the homology groups of the subspace of $I \times I$ consisting of the four
boundary edges plus all points in the interior whose first coordinate is rational.

\item Show that $\tilde H_n(X) \cong \tilde H_{n+1}(SX)$ for all $n$, where $SX$ is the suspension of $X$. More
generally, thinking of $SX$ as the union of two cones $CX$ with their bases identified,
compute the reduced homology groups of the union of any finite number of cones
$CX$ with their bases identified.

\item Making the preceding problem more concrete, construct explicit chain maps
$s: C_n(X) \to C_{n+1}(SX)$ inducing isomorphisms $\tilde H_n(X) \to \tilde H_{n+1}(SX)$.

\item Prove by induction on dimension the following facts about the homology of a
finite-dimensional CW complex $X$, using the observation that $X^n/X^{n-1}$ is a wedge
sum of $n$-spheres:
\begin{enumerate}
\item[(a)] If $X$ has dimension $n$ then $H_i(X) = 0$ for $i > n$ and $H_n(X)$ is free.
\item[(b)] $H_n(X)$ is free with basis in bijective correspondence with the $n$-cells if there are
no cells of dimension $n - 1$ or $n + 1$.
\item[(c)] If $X$ has $k$ $n$-cells, then $H_n(X)$ is generated by at most $k$ elements.
\end{enumerate}

%% --- END page-0141.tex ---
%% --- BEGIN page-0142.tex ---
\item Show that the second barycentric subdivision of a $\Delta$-complex is a simplicial
complex. Namely, show that the first barycentric subdivision produces a $\Delta$-complex
with the property that each simplex has all its vertices distinct, then show that for a
$\Delta$-complex with this property, barycentric subdivision produces a simplicial complex.

\item Show that each $n$-simplex in the barycentric subdivision of $\Delta^n$ is defined by $n$
inequalities $t_{i_0} \le t_{i_1} \le \cdots \le t_{i_n}$ in its barycentric coordinates, where $(i_0, \ldots, i_n)$ is a
permutation of $(0, \ldots, n)$.

\item Find an explicit, noninductive formula for the barycentric subdivision operator
$S: C_n(X) \to C_n(X)$.

\item Show that $H_1(X, A)$ is not isomorphic to $\tilde H_1(X/A)$ if $X = [0,1]$ and $A$ is the
sequence $1, 1/2, 1/3, \ldots$ together with its limit $0$. [See Example 1.25.]

\item Let $f:(X,A) \to (Y,B)$ be a map such that both $f: X \to Y$ and the restriction
$f: A \to B$ are homotopy equivalences.
\begin{enumerate}
\item[(a)] Show that $f_*: H_n(X,A) \to H_n(Y,B)$ is an isomorphism for all $n$.
\item[(b)] For the case of the inclusion $f:(D^n, S^{n-1}) \hookrightarrow (D^n, D^n - \{0\})$, show that $f$ is not
a homotopy equivalence of pairs --- there is no $g:(D^n, D^n - \{0\}) \to (D^n, S^{n-1})$ such
that $fg$ and $gf$ are homotopic to the identity through maps of pairs. [Observe that
a homotopy equivalence of pairs $(X,A) \to (Y,B)$ is also a homotopy equivalence for
the pairs obtained by replacing $A$ and $B$ by their closures.]
\end{enumerate}

\item Let $X$ be the cone on the 1-skeleton of $\Delta^3$, the union of all line segments joining
points in the six edges of $\Delta^3$ to the barycenter of $\Delta^3$. Compute the local homology
groups $H_n(X, X - \{x\})$ for all $x \in X$. Define $\partial X$ to be the subspace of points $x$
such that $H_n(X, X - \{x\}) = 0$ for all $n$, and compute the local homology groups
$H_n(\partial X, \partial X - \{x\})$. Use these calculations to determine which subsets $A \subset X$ have the
property that $f(A) \subset A$ for all homeomorphisms $f: X \to X$.

\item Show that $S^1 \times S^1$ and $S^1 \vee S^1 \vee S^2$ have isomorphic homology groups in all
dimensions, but their universal covering spaces do not.

\item In each of the following commutative diagrams assume that all maps but one are
isomorphisms. Show that the remaining map must be an isomorphism as well.
\[
A \to B \quad A \to B \quad A \to B
\]
\[
\searrow \quad \downarrow \quad \downarrow \quad \uparrow
\]
\[
C \quad C \to D \quad C \to D
\]

\item Using the notation of the five-lemma, give an example where the maps $\alpha, \beta, \delta,$
and $\varepsilon$ are zero but $\gamma$ is nonzero. This can be done with short exact sequences in
which all the groups are either $\Z$ or $0$.

\end{enumerate}

%% --- END page-0142.tex ---
%% --- BEGIN page-0143.tex ---
\section{Computations and Applications}

Now that the basic properties of homology have been established, we can begin to move a little more freely. Our first topic, exploiting the calculation of $H_n(S^n)$, is Brouwer's notion of degree for maps $S^n \to S^n$. Historically, Brouwer's introduction of this concept in the years 1910-12 preceded the rigorous development of homology, so his definition was rather different, using the technique of simplicial approximation which we explain in \S 2.C. The later definition in terms of homology is certainly more elegant, though perhaps with some loss of geometric intuition. More in the spirit of Brouwer's definition is a third approach using differential topology, presented very lucidly in [Milnor 1965].

\textbf{Degree}

For a map $f : S^n \to S^n$ with $n > 0$, the induced map $f_* : H_n(S^n) \to H_n(S^n)$ is a homomorphism from an infinite cyclic group to itself and so must be of the form $f_*(\alpha) = d\alpha$ for some integer $d$ depending only on $f$. This integer is called the \textbf{degree} of $f$, with the notation $\deg f$. Here are some basic properties of degree:

\begin{enumerate}
\item[(a)] $\deg \mathbb{1} = 1$, since $\mathbb{1}_* = \mathbb{1}$.

\item[(b)] $\deg f = 0$ if $f$ is not surjective. For if we choose a point $x_0 \in S^n - f(S^n)$ then $f$ can be factored as a composition $S^n \to S^n - \{x_0\} \hookrightarrow S^n$ and $H_n(S^n - \{x_0\}) = 0$ since $S^n - \{x_0\}$ is contractible. Hence $f_* = 0$.

\item[(c)] If $f \simeq g$ then $\deg f = \deg g$ since $f_* = g_*$. The converse statement, that $f \simeq g$ if $\deg f = \deg g$, is a fundamental theorem of Hopf from around 1925 which we prove in Corollary 4.25.

\item[(d)] $\deg f g = \deg f \deg g$, since $(fg)_* = f_* g_*$. As a consequence, $\deg f = \pm 1$ if $f$ is a homotopy equivalence since $fg \simeq \mathbb{1}$ implies $\deg f \deg g = \deg \mathbb{1} = 1$.

\item[(e)] $\deg f = -1$ if $f$ is a reflection of $S^n$, fixing the points in a subsphere $S^{n-1}$ and interchanging the two complementary hemispheres. For we can give $S^n$ a $\Delta$-complex structure with these two hemispheres as its two $n$-simplices $\Delta_1^n$ and $\Delta_2^n$, and the $n$-chain $\Delta_1^n - \Delta_2^n$ represents a generator of $H_n(S^n)$ as we saw in Example 2.23, so the reflection interchanging $\Delta_1^n$ and $\Delta_2^n$ sends this generator to its negative.

\item[(f)] The antipodal map $-\mathbb{1} : S^n \to S^n$, $x \mapsto -x$, has degree $(-1)^{n+1}$ since it is the composition of $n+1$ reflections, each changing the sign of one coordinate in $\mathbb{R}^{n+1}$.

\item[(g)] If $f : S^n \to S^n$ has no fixed points then $\deg f = (-1)^{n+1}$. For if $f(x) \neq x$ then the line segment from $f(x)$ to $-x$, defined by $t \mapsto (1-t)f(x) - tx$ for $0 \le t \le 1$, does not pass through the origin. Hence if $f$ has no fixed points, the formula $f_t(x) = [(1-t)f(x) - tx]/|(1-t)f(x) - tx|$ defines a homotopy from $f$ to
\end{enumerate}

%% --- END page-0143.tex ---
%% --- BEGIN page-0144.tex ---
the antipodal map. Note that the antipodal map has no fixed points, so the fact that maps without fixed points are homotopic to the antipodal map is a sort of converse statement.

Here is an interesting application of degree:

\begin{theorem}[Tangent vector fields on spheres]
\label{thm:tangent-vector-fields-spheres}
\dcref{ch2:2.28}
\group{Cellular Homology}
\uses{cor:homology-spheres}
$S^n$ has a continuous field of nonzero tangent vectors iff $n$ is odd.
\end{theorem}
\begin{proof}
 Suppose $x \mapsto v(x)$ is a tangent vector field on $S^n$, assigning to a vector $x \in S^n$ the vector $v(x)$ tangent to $S^n$ at $x$. Regarding $v(x)$ as a vector at the origin instead of at $x$, tangency just means that $x$ and $v(x)$ are orthogonal in $\mathbb{R}^{n+1}$. If $v(x) \ne 0$ for all $x$, we may normalize so that $|v(x)| = 1$ for all $x$ by replacing $v(x)$ by $v(x)/|v(x)|$. Assuming this has been done, the vectors $(\cos t)x + (\sin t)v(x)$ lie in the unit circle in the plane spanned by $x$ and $v(x)$. Letting $\mu$ go from $0$ to $\pi$, we obtain a homotopy $f_t(x) = (\cos t)x + (\sin t)v(x)$ from the identity map of $S^n$ to the antipodal map $-\mathbb{1}$. This implies that $\deg(-\mathbb{1}) = \deg \mathbb{1}$, hence $(-1)^{n+1} = 1$ and $n$ must be odd.

Conversely, if $n$ is odd, say $n = 2k-1$, we can define $v(x_1, x_2, \ldots, x_{2k-1}, x_{2k}) = (-x_2, x_1, \ldots, -x_{2k}, x_{2k-1})$. Then $v(x)$ is orthogonal to $x$, so $v$ is a tangent vector field on $S^n$, and $|v(x)| = 1$ for all $x \in S^n$.
\end{proof}
For the much more difficult problem of finding the maximum number of tangent vector fields on $S^n$ that are linearly independent at each point, see [VBKT] or [Husemoller 1966].

Another nice application of degree, giving a partial answer to a question raised in Example 1.43, is the following result:

\begin{proposition}[Only Z2 acts freely on even spheres]
\label{prop:z2-acts-freely-even-spheres}
\dcref{ch2:2.29}
\group{Cellular Homology}
\uses{thm:tangent-vector-fields-spheres}
$\mathbb{Z}_2$ is the only nontrivial group that can act freely on $S^n$ if $n$ is even.
\end{proposition}

Recall that an action of a group $G$ on a space $X$ is a homomorphism from $G$ to the group Homeo$(X)$ of homeomorphisms $X \to X$, and the action is free if the homeomorphism corresponding to each nontrivial element of $G$ has no fixed points. In the case of $S^n$, the antipodal map $x \mapsto -x$ generates a free action of $\mathbb{Z}_2$.
\begin{proof}
 Since homeomorphisms have degree $\pm 1$, an action of a group $G$ on $S^n$ determines a degree function $d: G \to \{\pm 1\}$. This is a homomorphism since $\deg fg = \deg f \deg g$. If the action is free, $d$ sends each nontrivial element of $G$ to $(-1)^{n+1}$ by property (g) above. Thus when $n$ is even, $d$ has trivial kernel, so $G \subset \mathbb{Z}_2$.
\end{proof}
Next we describe a technique for computing degrees which can be applied to most maps that arise in practice. Suppose $f: S^n \to S^n$, $n > 0$, has the property that for some point $y \in S^n$, the preimage $f^{-1}(y)$ consists of only finitely many points, say

%% --- END page-0144.tex ---
%% --- BEGIN page-0145.tex ---
$x_1, \ldots, x_m$. Let $U_1, \ldots, U_m$ be disjoint neighborhoods of these points, mapped by $f$
into a neighborhood $V$ of $y$. Then $f(U_i - x_i) \subset V - y$ for each $i$, and we have a
diagram
\[
\begin{array}{ccccc}
  & & H_n(U_i, U_i - x_i) & \xrightarrow{f_*} & H_n(V, V - y) \\
  & \simeq \nearrow & \downarrow k_i & & \uparrow \simeq \\
  H_n(S^n, S^n - x_i) & \xrightarrow{p_i} & H_n(S^n, S^n - f^{-1}(y)) & \xrightarrow{f_*} & H_n(S^n, S^n - y) \\
  \simeq \searrow & & \uparrow j & & \uparrow \simeq \\
  & & H_n(S^n) & \xrightarrow{f_*} & H_n(S^n)
\end{array}
\]

where all the maps are the obvious ones, and in particular $k_i$ and $p_i$ are induced
by inclusions, so the triangles and squares commute. The two isomorphisms in the
upper half of the diagram come from excision, while the lower two isomorphisms come
from exact sequences of pairs. Via these four isomorphisms, the top two groups in the
diagram can be identified with $H_n(S^n) \cong \mathbb{Z}$, and the top homomorphism $f_*$ becomes
multiplication by an integer called the \textbf{local degree} of $f$ at $x_i$, written $\deg f|_{x_i}$.

For example, if $f$ is a homeomorphism, then $y$ can be any point and there is
only one corresponding $x_i$, so all the maps in the diagram are isomorphisms and
$\deg f|_{x_i} = \deg f = \pm 1$. More generally, if $f$ maps each $U_i$ homeomorphically onto
$V$, then $\deg f|_{x_i} = \pm 1$ for each $i$. This situation occurs quite often in applications,
and it is usually not hard to determine the correct signs.

Here is the formula that reduces degree calculations to computing local degrees:

\begin{proposition}[Local degree formula]
\label{prop:local-degree-formula}
\dcref{ch2:2.30}
\group{Cellular Homology}
\uses{cor:homology-spheres}
$\deg f = \sum_i \deg f|_{x_i}$.
\end{proposition}
\begin{proof}
 By excision, the central term $H_n(S^n, S^n - f^{-1}(y))$ in the preceding diagram
is the direct sum of the groups $H_n(U_i, U_i - x_i) \cong \mathbb{Z}$, with $k_i$ the inclusion of the $i^{\text{th}}$
summand. The map $p_j$ is projection to the $i^{\text{th}}$ summand since the upper triangle
commutes and $p_j k_j = 0$ for $j \ne i$, as $p_j k_j$ factors through $H_n(U_j, U_j) = 0$. Identifying the outer groups in the diagram with $\mathbb{Z}$ as before, commutativity of the lower
triangle says that $p_j j(1) = 1$, hence $j(1) = (1, \ldots, 1) = \sum_i k_i(1)$. Commutativity of
the upper square says that the middle $f_*$ takes $k_i(1)$ to $\deg f|_{x_i}$, hence the sum
$\sum_i k_i(1) = j(1)$ is taken to $\sum_i \deg f|_{x_i}$. Commutativity of the lower square then
gives the formula $\deg f = \sum_i \deg f|_{x_i}$.
\end{proof}
\begin{example}[Degree of reflection and antipodal map]
\label{ex:degree-reflection-antipodal}
\dcref{ch2:2.31}
\group{Cellular Homology}
\uses{cor:homology-spheres}
We can use this result to construct a map $S^n \to S^n$ of any given degree,
for each $n \ge 1$. Let $q: S^n \to \bigvee_k S^n$ be the quotient map obtained by collapsing the
complement of $k$ disjoint open balls $B_i$ in $S^n$ to a point, and let $p: \bigvee_k S^n \to S^n$ identify
all the summands to a single sphere. Consider the composition $f = pq$. For almost all
$y \in S^n$ we have $f^{-1}(y)$ consisting of one point $x_i$ in each $B_i$. The local degree of $f$
at $x_i$ is $\pm 1$ since $f$ is a homeomorphism near $x_i$. By precomposing $p$ with reflections
of the summands of $\bigvee_k S^n$ if necessary, we can make each local degree either $+1$ or
$-1$, whichever we wish. Thus we can produce a map $S^n \to S^n$ of degree $\pm k$.
\end{example}

%% --- END page-0145.tex ---
%% --- BEGIN page-0146.tex ---
\begin{example}[Degree of suspension map]
\label{ex:degree-suspension}
\dcref{ch2:2.32}
\group{Cellular Homology}
\uses{def:suspension,cor:homology-spheres}
In the case of $S^1$, the map $f(z) = z^k$, where we view $S^1$ as the unit circle in $\mathbb{C}$, has degree $k$. This is evident in the case $k = 0$ since $f$ is then constant. The case $k < 0$ reduces to the case $k > 0$ by composing with $z \mapsto z^{-1}$, which is a reflection, of degree $-1$. To compute the degree when $k > 0$, observe first that for any $y \in S^1$, $f^{-1}(y)$ consists of $k$ points $x_1, \cdots, x_k$ near each of which $f$ is a local homeomorphism, stretching a circular arc by a factor of $k$. This local stretching can be eliminated by a deformation of $f$ near $x_i$ that does not change local degree, so the local degree at $x_i$ is the same as for a rotation of $S^1$. A rotation is a homeomorphism so its local degree at any point equals its global degree, which is $+1$ since a rotation is homotopic to the identity. Hence $\deg f |_{x_i} = 1$ and $\deg f = k$.
\end{example}

Another way of obtaining a map $S^n \to S^n$ of degree $k$ is to take a repeated suspension of the map $z \mapsto z^k$ in Example 2.32, since suspension preserves degree:

\begin{proposition}[Homology of CW complexes]
\label{prop:homology-cw-complexes}
\dcref{ch2:2.33}
\group{Cellular Homology}
\uses{def:cell-complex}
$\deg Sf = \deg f$, where $Sf : S^{n+1} \to S^{n+1}$ is the suspension of the map $f : S^n \to S^n$.
\end{proposition}
\begin{proof}
 Let $CS^n$ denote the cone $(S^n \times I)/(S^n \times 1)$ with base $S^n = S^n \times 0 \subset CS^n$, so $CS^n/S^n$ is the suspension of $S^n$. The map $f$ induces $Cf : (CS^n, S^n) \to (CS^n, S^n)$ with quotient $Sf$. The naturality of the boundary maps in the long exact sequence of the pair $(CS^n, S^n)$ then gives commutativity of the diagram at the right. Hence if $f_*$ is multiplication by $d$, so is $Sf_*$.
\end{proof}


Note that for $f : S^n \to S^n$, the suspension $Sf$ maps only one point to each of the two `poles' of $S^{n+1}$. This implies that the local degree of $Sf$ at each pole must equal the global degree of $Sf$. Thus the local degree of a map $S^n \to S^n$ can be any integer if $n \ge 2$, just as the degree itself can be any integer when $n \ge 1$.

\textbf{Cellular Homology}

Cellular homology is a very efficient tool for computing the homology groups of CW complexes, based on degree calculations. Before giving the definition of cellular homology, we first establish a few preliminary facts:

\begin{lemma}[Homology of CW skeleton pairs]
\label{lem:homology-cw-skeleton-pairs}
\dcref{ch2:2.34}
\group{Cellular Homology}
\uses{def:cell-complex}
If $X$ is a CW complex, then:
\end{lemma}

(a) $H_k(X^n, X^{n-1})$ is zero for $k \ne n$ and is free abelian for $k = n$, with a basis in one-to-one correspondence with the $n$-cells of $X$.

(b) $H_k(X^n) = 0$ for $k > n$. In particular, if $X$ is finite-dimensional then $H_k(X) = 0$ for $k > \dim X$.

(c) The map $H_k(X^n) \to H_k(X)$ induced by the inclusion $X^n \to X$ is an isomorphism for $k < n$ and surjective for $k = n$.
\begin{proof}
 Statement (a) follows immediately from the observation that $(X^n, X^{n-1})$ is a good pair and $X^n/X^{n-1}$ is a wedge sum of $n$-spheres, one for each $n$-cell of $X$. Here

%% --- END page-0146.tex ---
%% --- BEGIN page-0147.tex ---
we are using Proposition 2.22 and Corollary 2.25. Next consider the following part of
the long exact sequence of the pair $(X^n, X^{n-1})$:
\[
H_{k+1}(X^n, X^{n-1}) \to H_k(X^{n-1}) \to H_k(X^n) \to H_k(X^n, X^{n-1})
\]

If $k \ne n$ the last term is zero by part (a) so the middle map is surjective, while if
$k = n-1$ then the first term is zero so the middle map is injective. Now look at the
inclusion-induced homomorphisms
\[
H_k(X^0) \to H_k(X^1) \to \cdots \to H_k(X^{k-1}) \to H_k(X^k) \to H_k(X^{k+1}) \to \cdots
\]

By what we have just shown these are all isomorphisms except that the map to $H_k(X^k)$
may not be surjective and the map from $H_k(X^k)$ may not be injective. The first part
of the sequence then gives statement (b) since $H_k(X^0) = 0$ when $k > 0$. Also, the last
part of the sequence gives (c) when $X$ is finite-dimensional.

The proof of (c) when $X$ is infinite-dimensional requires more work, and this can
be done in two different ways. The more direct approach is to descend to the chain
level and use the fact that a singular chain in $X$ has compact image, hence meets only
finitely many cells of $X$ by Proposition A.1 in the Appendix. Thus each chain lies
in a finite skeleton $X^m$. So a $k$-cycle in $X$ is a cycle in some $X^m$, and then by the
finite-dimensional case of (c), the cycle is homologous to a cycle in $X^n$ if $n \ge k$, so
$H_k(X^n) \to H_k(X)$ is surjective. Similarly for injectivity, if a $k$-cycle in $X^n$ bounds a
chain in $X$, this chain lies in some $X^m$ with $m \ge n$, so by the finite-dimensional case
the cycle bounds a chain in $X^n$ if $n > k$.

The other approach is more general. From the long exact sequence of the pair
$(X, X^n)$ it suffices to show $H_k(X, X^n) = 0$ for $k \le n$. Since $H_k(X, X^n) \cong \tilde{H}_k(X/X^n)$,
this reduces the problem to showing:
\[
(*) \quad \tilde{H}_k(X) = 0 \text{ for } k \le n \text{ if the } n\text{-skeleton of } X \text{ is a point.}
\]

When $X$ is finite-dimensional, $(*)$ is immediate from the finite-dimensional case
of (c) which we have already shown. It will suffice therefore to reduce the infinite-
dimensional case to the finite-dimensional case. This reduction will be achieved by
stretching $X$ out to a complex that is at least locally finite-dimensional, using a special
case of the `mapping telescope' construction described in greater generality in \S 3.F.

Consider $X \times [0, \infty)$ with its product cell structure,
where we give $[0, \infty)$ the cell structure with the integer
points as 0-cells. Let $T = \bigcup_i X^i \times [i, \infty)$, a subcomplex
of $X \times [0, \infty)$. The figure shows a schematic picture of $T$ with $[0, \infty)$ in the hor-
izontal direction and the subcomplexes $X^i \times [i, i+1]$ as rectangles whose size in-
creases with $i$ since $X^i \subset X^{i+1}$. The line labeled $R$ can be ignored for now. We claim
that $T \simeq X$, hence $H_k(X) \cong H_k(T)$ for all $k$. Since $X$ is a deformation retract of
$X \times [0, \infty)$, it suffices to show that $X \times [0, \infty)$ also deformation retracts onto $T$. Let
$Y_i = T \cup (X \times [i, \infty))$. Then $Y_i$ deformation retracts onto $Y_{i+1}$ since $X \times [i, i+1]$ defor-
mation retracts onto $X^i \times [i, i+1] \cup X \times \{i+1\}$ by Proposition 0.16. If we perform the

%% --- END page-0147.tex ---
%% --- BEGIN page-0148.tex ---
deformation retraction of $Y_i$ onto $Y_{i+1}$ during the $t$-interval $[1 - 1/2^i, 1 - 1/2^{i+1}]$, then this gives a deformation retraction $f_t$ of $X \times [0, \infty)$ onto $T$, with points in $X' \times [0, \infty)$ stationary under $f_t$ for $t \geq 1 - 1/2^{i+1}$. Continuity follows from the fact that CW complexes have the weak topology with respect to their skeleta, so a map is continuous if its restriction to each skeleton is continuous.

Recalling that $X^0$ is a point, let $R \subset T$ be the ray $X^0 \times [0, \infty)$ and let $Z \subset T$ be the union of this ray with all the subcomplexes $X^i \times \{t\}$. Then $Z/R$ is homeomorphic to $\bigvee_i X^i$, a wedge sum of finite-dimensional complexes with $n$-skeleton a point, so the finite-dimensional case of $(*)$ together with Corollary 2.25 describing the homology of wedge sums implies that $\widetilde{H}_k(Z/R) = 0$ for $k \leq n$. The same is therefore true for $Z$, from the long exact sequence of the pair $(Z, R)$, since $R$ is contractible. Similarly, $T/Z$ is a wedge sum of finite-dimensional complexes with $(n+1)$-skeleton a point, since if we first collapse each subcomplex $X^i \times \{t\}$ of $T$ to a point, we obtain the infinite sequence of suspensions $SX^i$ `skewered' along the ray $R$, and then if we collapse $R$ to a point we obtain $\bigvee_i \Sigma X^i$ where $\Sigma X^i$ is the reduced suspension of $X^i$, obtained from $SX^i$ by collapsing the line segment $X^0 \times [t, i+1]$ to a point, so $\Sigma X^i$ has $(n+1)$-skeleton a point. Thus $\widetilde{H}_k(T/Z) = 0$ for $k \leq n+1$. The long exact sequence of the pair $(T,Z)$ then implies that $\widetilde{H}_k(T) = 0$ for $k \leq n$, and we have proved $(*)$.
\end{proof}
Let $X$ be a CW complex. Using Lemma 2.34, portions of the long exact sequences for the pairs $(X^{n+1}, X^n)$, $(X^n, X^{n-1})$, and $(X^{n-1}, X^{n-2})$ fit into a diagram

\[
\begin{array}{ccccc}
& & 0 & & \\
& \swarrow & & \searrow & \\
0 & \to & H_n(X^{n+1}) & \simeq & H_n(X) \\
& & \downarrow & & \downarrow \\
& & j_n & & \\
\cdots \to H_{n+1}(X^{n+1}, X^n) & \xrightarrow{d_{n+1}} & H_n(X^n, X^{n-1}) & \xrightarrow{d_n} & H_{n-1}(X^{n-1}, X^{n-2}) \to \cdots \\
& & \downarrow & & \downarrow \\
& & \delta_n & & j_{n-1} \\
& \swarrow & & \searrow & \\
& & 0 & & H_{n-1}(X^{n-1})
\end{array}
\]

where $d_{n+1}$ and $d_n$ are defined as the compositions $j_n \partial_{n+1}$ and $j_{n-1} \partial_n$, which are just `relativizations' of the boundary maps $\partial_{n+1}$ and $\partial_n$. The composition $d_n d_{n+1}$ includes two successive maps in one of the exact sequences, hence is zero. Thus the horizontal row in the diagram is a chain complex, called the \textbf{cellular chain complex} of $X$ since $H_n(X^n, X^{n-1})$ is free with basis in one-to-one correspondence with the $n$-cells of $X$, so one can think of elements of $H_n(X^n, X^{n-1})$ as linear combinations of $n$-cells of $X$. The homology groups of this cellular chain complex are called the \textbf{cellular homology groups} of $X$. Temporarily we denote them $H_n^{CW}(X)$.

\begin{theorem}[Isomorphism of cellular and singular homology]
\label{thm:cellular-singular-homology-isomorphic}
\dcref{ch2:2.35}
\group{Cellular Homology}
\uses{prop:homology-cw-complexes,lem:homology-cw-skeleton-pairs}
$H_n^{CW}(X) \simeq H_n(X)$.
\end{theorem}

%% --- END page-0148.tex ---
%% --- BEGIN page-0149.tex ---
\begin{proof}
 From the diagram above, $H_n(X)$ can be identified with $H_n(X^n)/\text{Im}\partial_{n+1}$. Since $j_n$ is injective, it maps $\text{Im}\partial_{n+1}$ isomorphically onto $\text{Im}(j_n\partial_{n+1}) = \text{Im}d_{n+1}$ and $H_n(X^n)$ isomorphically onto $\text{Im}j_n = \ker\partial_n$. Since $j_{n-1}$ is injective, $\ker\partial_n = \ker d_n$. Thus $j_n$ induces an isomorphism of the quotient $H_n(X^n)/\text{Im}\partial_{n+1}$ onto $\ker d_n/\text{Im}d_{n+1}$.
\end{proof}
Here are a few immediate applications:

(i) $H_n(X) = 0$ if $X$ is a CW complex with no $n$-cells.

(ii) More generally, if $X$ is a CW complex with $k$ $n$-cells, then $H_n(X)$ is generated by at most $k$ elements. For since $H_n(X^n, X^{n-1})$ is free abelian on $k$ generators, the subgroup $\ker d_n$ must be generated by at most $k$ elements, hence also the quotient $\ker d_n/\text{Im}d_{n+1}$.

(iii) If $X$ is a CW complex having no two of its cells in adjacent dimensions, then $H_n(X)$ is free abelian with basis in one-to-one correspondence with the $n$-cells of $X$. This is because the cellular boundary maps $d_n$ are automatically zero in this case.

This last observation applies for example to $\mathbb{CP}^n$, which has a CW structure with one cell of each even dimension $2k \le 2n$ as we saw in Example 0.6. Thus
\[
H_i(\mathbb{CP}^n) \cong \begin{cases} \mathbb{Z} & \text{for } i = 0,2,4,\ldots,2n \\ 0 & \text{otherwise} \end{cases}
\]

Another simple example is $S^n \times S^n$ with $n > 1$, using the product CW structure consisting of a 0-cell, two $n$-cells, and a $2n$-cell.

It is possible to prove the statements (i)--(iii) for finite-dimensional CW complexes by induction on the dimension, without using cellular homology but only the basic results from the previous section. However, the viewpoint of cellular homology makes (i)--(iii) quite transparent.

Next we describe how the cellular boundary maps $d_n$ can be computed. When $n = 1$ this is easy since the boundary map $d_1 : H_1(X^1, X^0) \to H_0(X^0)$ is the same as the simplicial boundary map $\Delta_1(X) \to \Delta_0(X)$. In case $X$ is connected and has only one 0-cell, then $d_1$ must be 0, otherwise $H_0(X)$ would not be $\mathbb{Z}$. When $n > 1$ we will show that $d_n$ can be computed in terms of degrees:

\textbf{Cellular Boundary Formula.} $d_n(e_\alpha^n) = \sum_\beta d_{\alpha\beta}e_\beta^{n-1}$ where $d_{\alpha\beta}$ is the degree of the map $S^{n-1} \to X^{n-1} \to S_\beta^{n-1}$ that is the composition of the attaching map of $e_\alpha^n$ with the quotient map collapsing $X^{n-1} - e_\beta^{n-1}$ to a point.

Here we are identifying the cells $e_\alpha^n$ and $e_\beta^{n-1}$ with generators of the corresponding summands of the cellular chain groups. The summation in the formula contains only finitely many terms since the attaching map of $e_\alpha^n$ has compact image, so this image meets only finitely many cells $e_\beta^{n-1}$.

To derive the cellular boundary formula, consider the commutative diagram

%% --- END page-0149.tex ---
%% --- BEGIN page-0150.tex ---
\[
\begin{array}{ccccccc}
H_n(D_\alpha^n, \partial D_\alpha^n) & \xrightarrow{\partial \approx} & \tilde{H}_{n-1}(\partial D_\alpha^n) & \xrightarrow{\Delta_{\alpha\beta*}} & \tilde{H}_{n-1}(S^{n-1}) \\
\big\downarrow \Phi_{\alpha*} & & \big\downarrow \varphi_{\alpha*} & & \big\downarrow q_{\beta*} \\
H_n(X^n, X^{n-1}) & \xrightarrow{\partial_n \approx} & \tilde{H}_{n-1}(X^{n-1}) & \xrightarrow{a_*} & \tilde{H}_{n-1}(X^{n-1}/X^{n-2}) \\
& \searrow d_n & \big\downarrow j_{n-1} & & \big\downarrow \approx \\
& & H_{n-1}(X^{n-1}, X^{n-2}) & \xrightarrow{\approx} & H_{n-1}(X^{n-1}/X^{n-2}, X^{n-2}/X^{n-2})
\end{array}
\]

where:
\begin{itemize}
\item $\Phi_\alpha$ is the characteristic map of the cell $e_\alpha^n$ and $\varphi_\alpha$ is its attaching map.
\item $q: X^{n-1} \to X^{n-1}/X^{n-2}$ is the quotient map.
\item $q_\beta: X^{n-1}/X^{n-2} \to S_\beta^{n-1}$ collapses the complement of the cell $e_\beta^{n-1}$ to a point, the resulting quotient sphere being identified with $S_\beta^{n-1} = D_\beta^{n-1}/\partial D_\beta^{n-1}$ via the characteristic map $\Phi_\beta$.
\item $\Delta_{\alpha\beta}: \partial D_\alpha^n \to S_\beta^{n-1}$ is the composition $q_\beta q\varphi_\alpha$, in other words, the attaching map of $e_\alpha^n$ followed by the quotient map $X^{n-1} \to S_\beta^{n-1}$ collapsing the complement of $e_\beta^{n-1}$ in $X^{n-1}$ to a point.
\end{itemize}

The map $\Phi_{\alpha*}$ takes a chosen generator $[D_\alpha^n] \in H_n(D_\alpha^n, \partial D_\alpha^n)$ to a generator of the $\mathbb{Z}$ summand of $H_n(X^n, X^{n-1})$ corresponding to $e_\alpha^n$. Letting $e_\alpha^n$ denote this generator, commutativity of the left half of the diagram then gives $d_n(e_\alpha^n) = j_{n-1} \varphi_{\alpha*} \partial [D_\alpha^n]$. In terms of the basis for $H_{n-1}(X^{n-1}, X^{n-2})$ corresponding to the cells $e_\beta^{n-1}$, the map $q_{\beta*}$ is the projection of $\tilde{H}_{n-1}(X^{n-1}/X^{n-2})$ onto its $\mathbb{Z}$ summand corresponding to $e_\beta^{n-1}$. Commutativity of the diagram then yields the formula for $d_n$ given above.

\begin{example}[Cellular homology of spheres]
\label{ex:cellular-homology-spheres}
\dcref{ch2:2.36}
\group{Cellular Homology}
\uses{thm:cellular-singular-homology-isomorphic,ex:sphere-two-cells}
Let $M_g$ be the closed orientable surface of genus $g$ with its usual CW structure consisting of one 0-cell, $2g$ 1-cells, and one 2-cell attached by the product of commutators $[a_1, b_1] \cdots [a_g, b_g]$. The associated cellular chain complex is
\[
0 \to \mathbb{Z} \xrightarrow{d_2} \mathbb{Z}^{2g} \xrightarrow{d_1} \mathbb{Z} \to 0
\]
\end{example}

As observed above, $d_1$ must be $0$ since there is only one 0-cell. Also, $d_2$ is $0$ because each $a_i$ or $b_i$ appears with its inverse in $[a_1, b_1] \cdots [a_g, b_g]$, so the maps $\Delta_{\alpha\beta}$ are homotopic to constant maps. Since $d_1$ and $d_2$ are both zero, the homology groups of $M_g$ are the same as the cellular chain groups, namely, $\mathbb{Z}$ in dimensions 0 and 2, and $\mathbb{Z}^{2g}$ in dimension 1.

\begin{example}[Cellular homology of surfaces]
\label{ex:cellular-homology-surfaces}
\dcref{ch2:2.37}
\group{Cellular Homology}
\uses{thm:cellular-singular-homology-isomorphic}
The closed nonorientable surface $N_g$ of genus $g$ has a cell structure with one 0-cell, $g$ 1-cells, and one 2-cell attached by the word $a_1^2 a_2^2 \cdots a_g^2$. Again, $d_1 = 0$, and $d_2: \mathbb{Z} \to \mathbb{Z}^g$ is specified by the equation $d_2(1) = (2, \ldots, 2)$ since each $a_i$ appears in the attaching word of the 2-cell with total exponent 2, which means that each $\Delta_{\alpha\beta}$ is homotopic to the map $z \mapsto z^2$, of degree 2. Since $d_2(1) = (2, \ldots, 2)$, we have $d_2$ injective and hence $H_2(N_g) = 0$. If we change the basis for $\mathbb{Z}^g$ by replacing the last standard basis element $(0, \ldots, 0, 1)$ by $(1, \ldots, 1)$, we see that $H_1(N_g) \approx \mathbb{Z}^{g-1} \oplus \mathbb{Z}_2$.
\end{example}

%% --- END page-0150.tex ---
%% --- BEGIN page-0151.tex ---
These two examples illustrate the general fact that the orientability of a closed connected manifold $M$ of dimension $n$ is detected by $H_n(M)$, which is $\mathbb{Z}$ if $M$ is orientable and $0$ otherwise. This is shown in Theorem 3.26.

\begin{example}[An Acyclic Space]
\label{ex:acyclic-space}
\dcref{ch2:2.38}
\group{Cellular Homology}
\uses{thm:cellular-singular-homology-isomorphic}
Let $X$ be obtained from $S^1 \vee S^1$ by attaching two 2-cells by the words $a^5 b^{-3}$ and $b^3(ab)^{-2}$. Then $d_2: \mathbb{Z}^2 \to \mathbb{Z}^2$ has matrix $\begin{pmatrix} 5 & -2 \\ -3 & 1 \end{pmatrix}$, with the two columns coming from abelianizing $a^5 b^{-3}$ and $b^3(ab)^{-2}$ to $5a - 3b$ and $-2a + b$, in additive notation. The matrix has determinant $-1$, so $d_z$ is an isomorphism and $\tilde{H}_i(X) = 0$ for all $i$. Such a space $X$ is called acyclic.
\end{example}

We can see that this acyclic space is not contractible by considering $\pi_1(X)$, which has the presentation $\langle a, b \mid a^5 b^{-3}, b^3(ab)^{-2} \rangle$. There is a nontrivial homomorphism from this group to the group $G$ of rotational symmetries of a regular dodecahedron, sending $a$ to the rotation $\rho_a$ through angle $2\pi/5$ about the axis through the center of a pentagonal face, and $b$ to the rotation $\rho_b$ through angle $2\pi/3$ about the axis through a vertex of this face. The composition $\rho_a \rho_b$ is a rotation through angle $\pi$ about the axis through the midpoint of an edge abutting this vertex. Thus the relations $a^5 = b^3 = (ab)^2$ defining $\pi_1(X)$ become $\rho_a^5 = \rho_b^3 = (\rho_a \rho_b)^2 = 1$ in $G$, which means there is a well-defined homomorphism $\rho: \pi_1(X) \to G$ sending $a$ to $\rho_a$ and $b$ to $\rho_b$. It is not hard to see that $G$ is generated by $\rho_a$ and $\rho_b$, so $\rho$ is surjective. With more work one can compute that the kernel of $\rho$ is $\mathbb{Z}_*$, generated by the element $a^5 = b^3 = (ab)^2$, and this $\mathbb{Z}_*$ is in fact the center of $\pi_1(X)$. In particular, $\pi_1(X)$ has order 120 since $G$ has order 60.

After these 2-dimensional examples, let us now move up to three dimensions, where we have the additional task of computing the cellular boundary map $d_3$.

\begin{example}[A 3-dimensional torus]
\label{ex:three-torus}
\dcref{ch2:2.39}
\group{Cellular Homology}
\uses{thm:cellular-singular-homology-isomorphic}
$T^3 = S^1 \times S^1 \times S^1$ can be constructed from a cube by identifying each pair of opposite square faces as in the first of the two figures. The second figure shows a slightly different pattern of identifications of opposite faces, with the front and back faces now identified via a rotation of the cube around a horizontal left-right axis. The space produced by these identifications is the product $K \times S^1$ of a Klein bottle and a circle. For both $T^3$ and $K \times S^1$ we have a CW structure with one 3-cell, three 2-cells, three 1-cells, and one 0-cell. The cellular chain complexes thus have the form
\end{example}



\[ 0 \to \mathbb{Z} \xrightarrow{d_3} \mathbb{Z}^3 \xrightarrow{d_2} \mathbb{Z}^3 \xrightarrow{d_1} \mathbb{Z} \to 0 \]

In the case of the 3-torus $T^3$ the cellular boundary map $d_2$ is zero by the same calculation as for the 2-dimensional torus. We claim that $d_3$ is zero as well. This amounts to saying that the three maps $\Delta_{\alpha \beta}: S^2 \to S^2$ corresponding to the three 2-cells

%% --- END page-0151.tex ---
%% --- BEGIN page-0152.tex ---
have degree zero. Each $\Delta_{\alpha\beta}$ maps the interiors of two opposite faces of the cube homeomorphically onto the complement of a point in the target $S^2$ and sends the remaining four faces to this point. Computing local degrees at the center points of the two opposite faces, we see that the local degree is $+1$ at one of these points and $-1$ at the other, since the restrictions of $\Delta_{\alpha\beta}$ to these two faces differ by a reflection of the boundary of the cube across the plane midway between them, and a reflection has degree $-1$. Since the cellular boundary maps are all zero, we deduce that $H_i(T^3)$ is $\mathbb{Z}$ for $i=0,3$, $\mathbb{Z}^3$ for $i=1,2$, and $0$ for $i > 3$.

For $K \times S^1$, when we compute local degrees for the front and back faces we find that the degrees now have the same rather than opposite signs since the map $\Delta_{\alpha\beta}$ on these two faces differs not by a reflection but by a rotation of the boundary of the cube. The local degrees for the other faces are the same as before. Using the letters $A, B, C$ to denote the 2-cells given by the faces orthogonal to the edges $a, b, c$, respectively, we have the boundary formulas $d_3c^3 = 2C$, $d_2A = 2b$, $d_2B = 0$, and $d_2C = 0$. It follows that $H_3(K \times S^1) = 0$, $H_2(K \times S^1) = \mathbb{Z} \oplus \mathbb{Z}_2$, and $H_1(K \times S^1) = \mathbb{Z} \oplus \mathbb{Z} \oplus \mathbb{Z}_2$.

Many more examples of a similar nature, quotients of a cube or other polyhedron with faces identified in some pattern, could be worked out in similar fashion. But let us instead turn to some higher-dimensional examples.

\begin{example}[Moore Spaces]
\label{ex:moore-spaces}
\dcref{ch2:2.40}
\group{Cellular Homology}
\uses{thm:cellular-singular-homology-isomorphic}
Given an abelian group $G$ and an integer $n \geq 1$, we will construct a CW complex $X$ such that $H_n(X) = G$ and $\tilde{H}_i(X) = 0$ for $i \neq n$. Such a space is called a \textbf{Moore space}, commonly written $M(G, n)$ to indicate the dependence on $G$ and $n$. It is probably best for the definition of a Moore space to include the condition that $M(G, n)$ be simply-connected if $n > 1$. The spaces we construct will have this property.
\end{example}

As an easy special case, when $G = \mathbb{Z}_m$ we can take $X$ to be $S^n$ with a cell $e^{n+1}$ attached by a map $S^n \to S^n$ of degree $m$. More generally, any finitely generated $G$ can be realized by taking wedge sums of examples of this type for finite cyclic summands of $G$, together with copies of $S^n$ for infinite cyclic summands of $G$.

In the general nonfinitely generated case let $F \to G$ be a homomorphism of a free abelian group $F$ onto $G$, sending a basis for $F$ onto some set of generators of $G$. The kernel $K$ of this homomorphism is a subgroup of a free abelian group, hence is itself free abelian. Choose bases $\{x_\alpha\}$ for $F$ and $\{y_\beta\}$ for $K$, and write $y_\beta = \sum_\alpha a_{\beta\alpha} x_\alpha$. Let $X^n = \bigvee_\alpha S^n_\alpha$, so $H_n(X^n) \approx F$ via Corollary 2.25. We will construct $X$ from $X^n$ by attaching cells $e^{n+1}_\beta$ via maps $f_\beta: S^n \to X^n$ such that the composition of $f_\beta$ with the projection onto the summand $S^n_\alpha$ has degree $d_{\beta\alpha}$. Then the cellular boundary map $d_{n+1}$ will be the inclusion $K \to F$, hence $X$ will have the desired homology groups.

The construction of $f_\beta$ generalizes the construction in Example 2.31 of a map $S^n \to S^n$ of given degree. Namely, we can let $f_\beta$ map the complement of $\sum_\alpha |d_{\beta\alpha}|$

%% --- END page-0152.tex ---
%% --- BEGIN page-0153.tex ---
disjoint balls in $S^n$ to the 0-cell of $X^n$ while sending $|d_{\beta\alpha}|$ of the balls onto the summand $S^n_\alpha$ by maps of degree $+1$ if $d_{\beta\alpha} > 0$, or degree $-1$ if $d_{\beta\alpha} < 0$.

\begin{example}[Cellular homology of CPn]
\label{ex:cellular-homology-cpn}
\dcref{ch2:2.41}
\group{Cellular Homology}
\uses{thm:cellular-singular-homology-isomorphic,ex:complex-projective-space-cw}
By taking a wedge sum of the Moore spaces constructed in the preceding example for varying $n$ we obtain a connected CW complex with any prescribed sequence of homology groups in dimensions $1, 2, 3, \cdots$.
\end{example}

\begin{example}[Cellular homology of $\mathbb{RP}^n$]
\label{ex:cellular-homology-rpn}
\dcref{ch2:2.42}
\group{Cellular Homology}
\uses{thm:cellular-singular-homology-isomorphic,ex:real-projective-space-cw}
As we saw in Example 0.4, $\mathbb{RP}^n$ has a CW structure with one cell $e^k$ in each dimension $k \le n$, and the attaching map for $e^k$ is the 2-sheeted covering projection $\varphi: S^{k-1} \to \mathbb{RP}^{k-1}$. To compute the boundary map $d_k$, we compute the degree of the composition $S^{k-1} \xrightarrow{\varphi} \mathbb{RP}^{k-1} \xrightarrow{q} \mathbb{RP}^{k-1}/\mathbb{RP}^{k-2} = S^{k-1}$, with $q$ the quotient map. The map $q\varphi$ restricts to a homeomorphism from each component of $S^{k-1} - S^{k-2}$ onto $\mathbb{RP}^{k-1} - \mathbb{RP}^{k-2}$, and these two homeomorphisms are obtained from each other by precomposing with the antipodal map of $S^{k-1}$, which has degree $(-1)^k$. Hence $\deg q\varphi = \deg \mathbb{1} + \deg(-\mathbb{1}) = 1 + (-1)^k$, and so $d_k$ is either 0 or multiplication by 2 according to whether $k$ is odd or even. Thus the cellular chain complex for $\mathbb{RP}^n$ is
\end{example}

\[
0 \to \mathbb{Z} \xrightarrow{2} \mathbb{Z} \xrightarrow{0} \cdots \xrightarrow{2} \mathbb{Z} \xrightarrow{0} \mathbb{Z} \xrightarrow{2} \mathbb{Z} \xrightarrow{0} \mathbb{Z} \to 0 \quad \text{if } n \text{ is even}
\]

\[
0 \to \mathbb{Z} \xrightarrow{0} \mathbb{Z} \xrightarrow{2} \cdots \xrightarrow{2} \mathbb{Z} \xrightarrow{0} \mathbb{Z} \xrightarrow{2} \mathbb{Z} \xrightarrow{0} \mathbb{Z} \to 0 \quad \text{if } n \text{ is odd}
\]

From this it follows that

\[
H_k(\mathbb{RP}^n) = \begin{cases}
\mathbb{Z} & \text{for } k = 0 \text{ and for } k = n \text{ odd} \\
\mathbb{Z}_2 & \text{for } k \text{ odd}, 0 < k < n \\
0 & \text{otherwise}
\end{cases}
\]

\begin{example}[Lens Spaces]
\label{ex:lens-spaces}
\dcref{ch2:2.43}
\group{Cellular Homology}
\uses{thm:cellular-singular-homology-isomorphic}
This example is somewhat more complicated. Given an integer $m > 1$ and integers $\ell_1, \ldots, \ell_n$ relatively prime to $m$, define the \textbf{lens space} $L = L_m(\ell_1, \ldots, \ell_n)$ to be the orbit space $S^{2n-1}/\mathbb{Z}_m$ of the unit sphere $S^{2n-1} \subset \mathbb{C}^n$ with the action of $\mathbb{Z}_m$ generated by the rotation $\rho(z_1, \ldots, z_n) = (e^{2\pi i\ell_1/m} z_1, \ldots, e^{2\pi i\ell_n/m} z_n)$, rotating the $j^{\text{th}}$ $\mathbb{C}$ factor of $\mathbb{C}^n$ by the angle $2\pi \ell_j/m$. In particular, when $m = 2$, $\rho$ is the antipodal map, so $L = \mathbb{RP}^{2n-1}$ in this case. In the general case, the projection $S^{2n-1} \to L$ is a covering space since the action of $\mathbb{Z}_m$ on $S^{2n-1}$ is free; only the identity element fixes any point of $S^{2n-1}$ since each point of $S^{2n-1}$ has some coordinate $z_j$ nonzero and then $e^{2\pi ik\ell_j/m} z_j \ne z_j$ for $0 < k < m$, as a result of the assumption that $\ell_j$ is relatively prime to $m$.
\end{example}

We shall construct a CW structure on $L$ with one cell $e^k$ for each $k \le 2n - 1$ and show that the resulting cellular chain complex is

\[
0 \to \mathbb{Z} \xrightarrow{0} \mathbb{Z} \xrightarrow{m} \mathbb{Z} \xrightarrow{0} \cdots \xrightarrow{0} \mathbb{Z} \xrightarrow{m} \mathbb{Z} \xrightarrow{0} \mathbb{Z} \to 0
\]

with boundary maps alternately 0 and multiplication by $m$. Hence

\[
H_k(L_m(\ell_1, \ldots, \ell_n)) = \begin{cases}
\mathbb{Z} & \text{for } k = 0, 2n - 1 \\
\mathbb{Z}_m & \text{for } k \text{ odd}, 0 < k < 2n - 1 \\
0 & \text{otherwise}
\end{cases}
\]

%% --- END page-0153.tex ---
%% --- BEGIN page-0154.tex ---
To obtain the CW structure, first subdivide the unit circle $C$ in the $n^{th}$ $\mathbb{C}$ factor of $\mathbb{C}^n$ by taking the points $e^{2\pi ij/m} \in C$ as vertices, $j = 1, \ldots, m$. Joining the $j^{th}$ vertex of $C$ to the unit sphere $S^{2n-3} \subset \mathbb{C}^{n-1}$ by arcs of great circles in $S^{2n-1}$ yields a $(2n-2)$-dimensional ball $B_j^{2n-2}$ bounded by $S^{2n-3}$. Specifically, $B_j^{2n-2}$ consists of the points $\cos\theta(0, \cdots, 0, e^{2\pi ij/m}) + \sin\theta(z_1, \cdots, z_{n-1}, 0)$ for $0 \le \theta \le \pi/2$. Similarly, joining the $j^{th}$ edge of $C$ to $S^{2n-3}$ gives a ball $B_j^{2n-1}$ bounded by $B_j^{2n-2}$ and $B_{j+1}^{2n-2}$, subscripts being taken mod $m$. The rotation $\rho$ carries $S^{2n-3}$ to itself and rotates $C$ by the angle $2\pi\ell_n/m$, hence $\rho$ permutes the $B_j^{2n-2}$'s and the $B_j^{2n-1}$'s. A suitable power of $\rho$, namely $\rho^r$ where $r\ell_n \equiv 1 \pmod{m}$, takes each $B_j^{2n-2}$ and $B_j^{2n-1}$ to the next one. Since $\rho^r$ has order $m$, it is also a generator of the rotation group $\mathbb{Z}_m$, and hence we may obtain $L$ as the quotient of one $B_j^{2n-1}$ by identifying its two faces $B_j^{2n-2}$ and $B_{j+1}^{2n-2}$ together via $\rho^r$.

In particular, when $n = 2$, $B_j^{2n-1}$ is a lens-shaped 3-ball and $L$ is obtained from this ball by identifying its two curved disk faces via $\rho^r$, which may be described as the composition of the reflection across the plane containing the rim of the lens, taking one face of the lens to the other, followed by a rotation of this face through the angle $2\pi\ell/m$ where $\ell = r\ell_1$. The figure illustrates the case $(m, \ell) = (7, 2)$, with the two dots indicating a typical pair of identified points in the upper and lower faces of the lens. Since the lens space $L$ is determined by the rotation angle $2\pi\ell/m$, it is conveniently written $L_{\ell/m}$. Clearly only the mod $m$ value of $\ell$ matters. It is a classical theorem of Reidemeister from the 1930s that $L_{\ell/m}$ is homeomorphic to $L_{\ell'/m'}$ iff $m' = m$ and $\ell'' = \pm\ell^{\pm 1} \pmod{m}$. For example, when $m = 7$ there are only two distinct lens spaces $L_{1/7}$ and $L_{2/7}$. The `if' part of this theorem is easy: Reflecting the lens through a mirror shows that $L_{\ell/m} \cong L_{-\ell/m}$, and interchanging the roles of the two $\mathbb{C}$ factors of $\mathbb{C}^2$ one obtains $L_{\ell/m} \cong L_{\ell-1/m}$. In the converse direction, $L_{\ell/m} \cong L_{\ell'/m'}$ clearly implies $m = m'$ since $\pi_1(L_{\ell/m}) \cong \mathbb{Z}_m$. The rest of the theorem takes considerably more work, involving either special 3-dimensional techniques or more algebraic methods that generalize to classify the higher-dimensional lens spaces as well. The latter approach is explained in [Cohen 1973].

Returning to the construction of a CW structure on $L_m(\ell_1, \cdots, \ell_n)$, observe that the $(2n-3)$-dimensional lens space $L_m(\ell_1, \cdots, \ell_{n-1})$ sits in $L_m(\ell_1, \cdots, \ell_n)$ as the quotient of $S^{2n-3}$, and $L_m(\ell_1, \cdots, \ell_n)$ is obtained from this subspace by attaching two cells, of dimensions $2n-2$ and $2n-1$, coming from the interiors of $B_j^{2n-1}$ and its two identified faces $B_j^{2n-2}$ and $B_{j+1}^{2n-2}$. Inductively this gives a CW structure on $L_m(\ell_1, \cdots, \ell_n)$ with one cell $e^k$ in each dimension $k \le 2n-1$.

The boundary maps in the associated cellular chain complex are computed as follows. The first one, $d_{2n-1}$, is zero since the identification of the two faces of $B_j^{2n-1}$ is via a reflection (degree $-1$) across $B_j^{2n-2}$ fixing $S^{2n-3}$, followed by a rotation-



%% --- END page-0154.tex ---
%% --- BEGIN page-0155.tex ---
tion (degree $+1$), so $d_{2n-1}(e^{2n-1}) = e^{2n-2} - e^{2n-2} = 0$. The next boundary map
$d_{2n-2}$ takes $e^{2n-2}$ to $me^{2n-3}$ since the attaching map for $e^{2n-2}$ is the quotient map
$S^{2n-3} \to L_m(\ell_1, \ldots, \ell_{n-1})$ and the balls $B_j^{2n-3}$ in $S^{2n-3}$ which project down onto $e^{2n-3}$
are permuted cyclically by the rotation $\rho$ of degree $+1$. Inductively, the subsequent
boundary maps $d_k$ then alternate between $0$ and multiplication by $m$.

Also of interest are the infinite-dimensional lens spaces $L_m(\ell_1, \ell_2, \ldots) = S^\infty/\mathbb{Z}_m$
defined in the same way as in the finite-dimensional case, starting from a sequence of
integers $\ell_1, \ell_2, \ldots$ relatively prime to $m$. The space $L_m(\ell_1, \ell_2, \ldots)$ is the union of the
increasing sequence of finite-dimensional lens spaces $L_m(\ell_1, \ldots, \ell_n)$ for $n = 1, 2, \ldots$,
each of which is a subcomplex of the next in the cell structure we have just constructed, so $L_m(\ell_1, \ell_2, \ldots)$ is also a CW complex. Its cellular chain complex consists
of a $\mathbb{Z}$ in each dimension with boundary maps alternately $0$ and $m$, so its reduced
homology consists of a $\mathbb{Z}_m$ in each odd dimension.

In the terminology of \S1.B, the infinite-dimensional lens space $L_m(\ell_1, \ell_2, \ldots)$ is
an Eilenberg-MacLane space $K(\mathbb{Z}_m, 1)$ since its universal cover $S^\infty$ is contractible, as
we showed there. By Theorem 1B.8 the homotopy type of $L_m(\ell_1, \ell_2, \ldots)$ depends
only on $m$, and not on the $\ell_i$'s. This is not true in the finite-dimensional case, when
two lens spaces $L_m(\ell_1, \ldots, \ell_n)$ and $L_m(\ell_1', \ldots, \ell_n')$ have the same homotopy type
iff $\ell_1 \cdots \ell_n \equiv \pm k^n \ell_1' \cdots \ell_n'$ mod $m$ for some integer $k$. A proof of this is outlined in
Exercise 2 in \S3.E and Exercise 29 in \S4.2. For example, the 3-dimensional lens spaces
$L_{1/5}$ and $L_{2/5}$ are not homotopy equivalent, though they have the same fundamental
group and the same homology groups. On the other hand, $L_{1/7}$ and $L_{2/7}$ are homotopy
equivalent but not homeomorphic.

\textbf{Euler Characteristic}

For a finite CW complex $X$, the \textbf{Euler characteristic} $\chi(X)$ is defined to be the
alternating sum $\sum_n (-1)^n c_n$ where $c_n$ is the number of $n$-cells of $X$, generalizing
the familiar formula \textit{vertices} $-$ \textit{edges} $+$ \textit{faces} for 2-dimensional complexes. The
following result shows that $\chi(X)$ can be defined purely in terms of homology, and
hence depends only on the homotopy type of $X$. In particular, $\chi(X)$ is independent
of the choice of CW structure on $X$.

\begin{theorem}[Euler characteristic in terms of homology]
\label{thm:euler-characteristic-homology}
\dcref{ch2:2.44}
\group{Cellular Homology}
\uses{def:dimension-cw-complex}
$\chi(X) = \sum_n (-1)^n \text{rank} H_n(X)$.
\end{theorem}

Here the rank of a finitely generated abelian group is the number of $\mathbb{Z}$ summands
when the group is expressed as a direct sum of cyclic groups. We shall need the
following fact, whose proof we leave as an exercise: If $0 \to A \to B \to C \to 0$ is a short
exact sequence of finitely generated abelian groups, then rank $B$ = rank $A$ + rank $C$.
\begin{proof}[2.44]
 This is purely algebraic. Let

\[ 0 \to C_k \xrightarrow{d_k} C_{k-1} \to \cdots \to C_1 \xrightarrow{d_1} C_0 \to 0 \]

%% --- END page-0155.tex ---
%% --- BEGIN page-0156.tex ---
be a chain complex of finitely generated abelian groups, with cycles $Z_n = \ker d_n$,
boundaries $B_n = \operatorname{Im} d_{n+1}$, and homology $H_n = Z_n/B_n$. Thus we have short exact
sequences $0 \to Z_n \to C_n \to B_{n-1} \to 0$ and $0 \to B_n \to Z_n \to H_n \to 0$, hence

\[
\text{rank } C_n = \text{rank } Z_n + \text{rank } B_{n-1}
\]

\[
\text{rank } Z_n = \text{rank } B_n + \text{rank } H_n
\]

Now substitute the second equation into the first, multiply the resulting equation by
$(-1)^n$, and sum over $n$ to get $\sum_n (-1)^n \text{rank } C_n = \sum_n (-1)^n \text{rank } H_n$. Applying this
with $C_n = H_n(X^n, X^{n-1})$ then gives the theorem.
\end{proof}
For example, the surfaces $M_g$ and $N_g$ have Euler characteristics $\chi(M_g) = 2 - 2g$
and $\chi(N_g) = 2 - g$. Thus all the orientable surfaces $M_g$ are distinguished from each
other by their Euler characteristics, as are the nonorientable surfaces $N_g$, and there
are only the relations $\chi(M_g) = \chi(N_{2g})$.

\textbf{Split Exact Sequences}

Suppose one has a retraction $r: X \to A$, so $ri = \mathbb{1}$ where $i: A \to X$ is the inclusion.
The induced map $i_*: H_n(A) \to H_n(X)$ is then injective since $r_* i_* = \mathbb{1}$. From this it
follows that the boundary maps in the long exact sequence for $(X, A)$ are zero, so the
long exact sequence breaks up into short exact sequences

\[
0 \to H_n(A) \xrightarrow{i_*} H_n(X) \xrightarrow{j_*} H_n(X, A) \to 0
\]

The relation $r_* i_* = \mathbb{1}$ actually gives more information than this, by the following
piece of elementary algebra:

\begin{lemma}[Splitting Lemma]
\label{lem:splitting-lemma}
\group{Singular Homology}
\uses{}
For a short exact sequence $0 \to A \xrightarrow{i} B \xrightarrow{j} C \to 0$ of abelian groups the following statements are equivalent:
\begin{enumerate}
\item[(a)] There is a homomorphism $p: B \to A$ such that $pi = \mathbb{1}: A \to A$.
\item[(b)] There is a homomorphism $s: C \to B$ such that $js = \mathbb{1}: C \to C$.
\item[(c)] There is an isomorphism $B \cong A \oplus C$ making a commutative diagram as at the right, where
\[
\begin{array}{ccccccccc}
0 & \to & A & \xrightarrow{i} & B & \xrightarrow{j} & C & \to & 0 \\
  &     &   &                & \cong &               &   &     &  \\
0 & \to & A & \xrightarrow{a \mapsto (a,0)} & A \oplus C & \xrightarrow{(a,c) \mapsto c} & C & \to & 0
\end{array}
\]
the maps in the lower row are the obvious ones, $a \mapsto (a, 0)$ and $(a, c) \mapsto c$.
\end{enumerate}

If these conditions are satisfied, the exact sequence is said to \textbf{split}. Note that (c) is symmetric: There is no essential difference between the roles of $A$ and $C$.
\end{lemma}

\begin{proof}[Sketch of proof]
For the implication (a) $\Rightarrow$ (c) one checks that the map $B \to A \oplus C$, $b \mapsto (p(b), j(b))$, is an isomorphism with the desired properties. For (b) $\Rightarrow$ (c) one uses instead the map $A \oplus C \to B$, $(a, c) \mapsto i(a) + s(c)$. The opposite implications (c) $\Rightarrow$ (a) and (c) $\Rightarrow$ (b) are fairly obvious. If one wants to show (b) $\Rightarrow$ (a) directly, one can define $p(b) = i^{-1}(b - sj(b))$. Further details are left to the reader.
\end{proof}

%% --- END page-0156.tex ---
%% --- BEGIN page-0157.tex ---
Except for the implications (b) $\Rightarrow$ (a) and (b) $\Rightarrow$ (c), the proof works equally well for nonabelian groups. In the nonabelian case, (b) is definitely weaker than (a) and (c), and short exact sequences satisfying (b) only determine $B$ as a semidirect product of $A$ and $C$. The difficulty is that $s(C)$ might not be a normal subgroup of $B$. In the nonabelian case one defines `splitting' to mean that (b) is satisfied.

In both the abelian and nonabelian contexts, if $C$ is free then every exact sequence $0 \to A \to B \to C \to 0$ splits, since one can define $s : C \to B$ by choosing a basis $\{c_\alpha\}$ for $C$ and letting $s(c_\alpha)$ be any element $b_\alpha \in B$ such that $j(b_\alpha) = c_\alpha$. The converse is also true: If every short exact sequence ending in $C$ splits, then $C$ is free. This is because for every $C$ there is a short exact sequence $0 \to A \to B \to C \to 0$ with $B$ free --- choose generators for $C$ and let $B$ have a basis in one-to-one correspondence with these generators, then let $B \to C$ send each basis element to the corresponding generator --- so if this sequence $0 \to A \to B \to C \to 0$ splits, $C$ is isomorphic to a subgroup of a free group, hence is free.

From the Splitting Lemma and the remarks preceding it we deduce that a retraction $r : X \to A$ gives a splitting $H_n(X) \cong H_n(A) \oplus H_n(X,A)$. This can be used to show the nonexistence of such a retraction in some cases, for example in the situation of the Brouwer fixed point theorem, where a retraction $D^n \to S^{n-1}$ would give an impossible splitting $H_{n-1}(D^n) \cong H_{n-1}(S^{n-1}) \oplus H_{n-1}(D^n, S^{n-1})$. For a somewhat more subtle example, consider the mapping cylinder $M_f$ of a degree $m$ map $f : S^n \to S^n$ with $m > 1$. If $M_f$ retracted onto the $S^n \subset M_f$ corresponding to the domain of $f$, we would have a split short exact sequence



But this sequence does not split since $\mathbb{Z}$ is not isomorphic to $\mathbb{Z} \oplus \mathbb{Z}_m$ if $m > 1$, so the retraction cannot exist. In the simplest case of the degree 2 map $S^1 \to S^1$, $z \mapsto z^2$, this says that the M\"obius band does not retract onto its boundary circle.

\textbf{Homology of Groups}

In \S1.B we constructed for each group $G$ a CW complex $K(G,1)$ having a contractible universal cover, and we showed that the homotopy type of such a space $K(G,1)$ is uniquely determined by $G$. The homology groups $H_n(K(G,1))$ therefore depend only on $G$, and are usually denoted simply $H_n(G)$. The calculations for lens spaces in Example 2.43 show that $H_n(\mathbb{Z}_m)$ is $\mathbb{Z}_m$ for odd $n$ and $0$ for even $n > 0$. Since $S^1$ is a $K(\mathbb{Z},1)$ and the torus is a $K(\mathbb{Z} \times \mathbb{Z},1)$, we also know the homology of these two groups. More generally, the homology of finitely generated abelian groups can be computed from these examples using the K\"unneth formula in \S3.B and the fact that a product $K(G,1) \times K(H,1)$ is a $K(G \times H,1)$.

Here is an application of the calculation of $H_n(\mathbb{Z}_m)$:

%% --- END page-0157.tex ---
%% --- BEGIN page-0158.tex ---
\begin{proposition}[Fundamental group of finite-dimensional K(G,1) is torsionfree]
\label{prop:fundamental-group-kg1-torsionfree}
\dcref{ch2:2.45}
\group{Cellular Homology}
\uses{def:k-g-1-space,ex:lens-spaces}
If a finite-dimensional CW complex $X$ is a $K(G,1)$, then the group $G = \pi_1(X)$ must be torsionfree.
\end{proposition}

This applies to quite a few manifolds, for example closed surfaces other than $S^2$ and $\mathbb{RP}^2$, and also many 3-dimensional manifolds such as complements of knots in $S^3$.
\begin{proof}
 If $G$ had torsion, it would have a finite cyclic subgroup $\mathbb{Z}_m$ for some $m > 1$, and the covering space of $X$ corresponding to this subgroup of $G = \pi_1(X)$ would be a $K(\mathbb{Z}_m, 1)$. Since $X$ is a finite-dimensional CW complex, the same would be true of its covering space $K(\mathbb{Z}_m, 1)$, and hence the homology of the $K(\mathbb{Z}_m, 1)$ would be nonzero in only finitely many dimensions. But this contradicts the fact that $H_n(\mathbb{Z}_m)$ is nonzero for infinitely many values of $n$.
\end{proof}
Reflecting the richness of group theory, the homology of groups has been studied quite extensively. A good starting place for those wishing to learn more is the textbook [Brown 1982]. At a more advanced level the books [Adem \& Milgram 1994] and [Benson 1992] treat the subject from a mostly topological viewpoint.

\textbf{Mayer-Vietoris Sequences}

In addition to the long exact sequence of homology groups for a pair $(X,A)$, there is another sort of long exact sequence, known as a Mayer-Vietoris sequence, which is equally powerful but is sometimes more convenient to use. For a pair of subspaces $A, B \subset X$ such that $X$ is the union of the interiors of $A$ and $B$, this exact sequence has the form

\[ \cdots \to H_n(A \cap B) \xrightarrow{\Phi} H_n(A) \oplus H_n(B) \xrightarrow{\psi} H_n(X) \xrightarrow{\delta} H_{n-1}(A \cap B) \to \cdots \]
\[ \cdots \to H_0(X) \to 0 \]

In addition to its usefulness for calculations, the Mayer-Vietoris sequence is also applied frequently in induction arguments, where one might know that a certain statement is true for $A$, $B$, and $A \cap B$ by induction and then deduce that it is true for $A \cup B$ by the exact sequence.

The Mayer-Vietoris sequence is easy to derive from the machinery of \S2.1. Let $C_n(A+B)$ be the subgroup of $C_n(X)$ consisting of chains that are sums of chains in $A$ and chains in $B$. The usual boundary map $\partial : C_n(X) \to C_{n-1}(X)$ takes $C_n(A+B)$ to $C_{n-1}(A+B)$, so the $C_n(A+B)$'s form a chain complex. According to Proposition 2.21, the inclusions $C_n(A+B) \hookrightarrow C_n(X)$ induce isomorphisms on homology groups. The Mayer-Vietoris sequence is then the long exact sequence of homology groups associated to the short exact sequence of chain complexes formed by the short exact sequences

\[ 0 \to C_n(A \cap B) \xrightarrow{\Phi} C_n(A) \oplus C_n(B) \xrightarrow{\psi} C_n(A+B) \to 0 \]

%% --- END page-0158.tex ---
%% --- BEGIN page-0159.tex ---
where $\varphi(x) = (x,-x)$ and $\psi(x,y) = x + y$. The exactness of this short exact sequence can be checked as follows. First, $\ker \varphi = 0$ since a chain in $A \cap B$ that is zero as a chain in $A$ (or in $B$) must be the zero chain. Next, $\operatorname{Im} \varphi \subset \ker \psi$ since $\psi\varphi = 0$. Also, $\ker \psi \subset \operatorname{Im} \varphi$ since for a pair $(x,y) \in C_n(A) \oplus C_n(B)$ the condition $x + y = 0$ implies $x = -y$, so $x$ is a chain in both $A$ and $B$, that is, $x \in C_n(A \cap B)$, and $(x,y) = (x,-x) \in \operatorname{Im} \varphi$. Finally, exactness at $C_n(A + B)$ is immediate from the definition of $C_n(A + B)$.

The boundary map $\partial: H_n(X) \to H_{n-1}(A \cap B)$ can easily be made explicit. A class $\alpha \in H_n(X)$ is represented by a cycle $z$, and by barycentric subdivision or some other method we can choose $z$ to be a sum $x + y$ of chains in $A$ and $B$, respectively. It need not be true that $x$ and $y$ are cycles individually, but $\partial x = -\partial y$ since $\partial(x + y) = 0$, and the element $\partial \alpha \in H_{n-1}(A \cap B)$ is represented by the cycle $\partial x = -\partial y$, as is clear from the definition of the boundary map in the long exact sequence of homology groups associated to a short exact sequence of chain complexes.

There is also a formally identical Mayer-Vietoris sequence for reduced homology groups, obtained by augmenting the previous short exact sequence of chain complexes in the obvious way:

\[
\begin{array}{ccccccccc}
0 & \to & C_0(A \cap B) & \xrightarrow{\varphi} & C_0(A) \oplus C_0(B) & \xrightarrow{\psi} & C_0(A + B) & \to & 0 \\
& & \downarrow \varepsilon & & \downarrow \varepsilon \oplus \varepsilon & & \downarrow \varepsilon & & \\
0 & \to & \mathbb{Z} & \xrightarrow{\varphi} & \mathbb{Z} \oplus \mathbb{Z} & \xrightarrow{\psi} & \mathbb{Z} & \to & 0
\end{array}
\]

Mayer-Vietoris sequences can be viewed as analogs of the van Kampen theorem since if $A \cap B$ is path-connected, the $H_1$ terms of the reduced Mayer-Vietoris sequence yield an isomorphism $H_1(X) \cong (H_1(A) \oplus H_1(B)) / \operatorname{Im} \Phi$. This is exactly the abelianized statement of the van Kampen theorem, and $H_1$ is the abelianization of $\pi_1$ for path-connected spaces, as we show in \S 2.A.

There are also Mayer-Vietoris sequences for decompositions $X = A \cup B$ such that $A$ and $B$ are deformation retracts of neighborhoods $U$ and $V$ with $U \cap V$ deformation retracting onto $A \cap B$. Under these assumptions the five-lemma implies that the maps $C_n(A + B) \to C_n(U \cup V)$ induce isomorphisms on homology, and hence so do the maps $C_n(A + B) \to C_n(X)$, which was all that we needed to obtain a Mayer-Vietoris sequence. For example, if $X$ is a CW complex and $A$ and $B$ are subcomplexes, then we can choose for $U$ and $V$ neighborhoods of the form $N_\varepsilon(A)$ and $N_\varepsilon(B)$ constructed in the Appendix, which have the property that $N_\varepsilon(A) \cap N_\varepsilon(B) = N_\varepsilon(A \cap B)$.

\begin{example}[Mayer-Vietoris sequence of sphere split into hemispheres]
\label{ex:mayer-vietoris-sphere}
\dcref{ch2:2.46}
\group{Singular Homology}
\uses{thm:excision-general}
Take $X = S^n$ with $A$ and $B$ the northern and southern hemispheres, so that $A \cap B = S^{n-1}$. Then in the reduced Mayer-Vietoris sequence the terms $\tilde{H}_i(A)$ and $\tilde{H}_i(B)$ are zero, so we obtain isomorphisms $\tilde{H}_i(S^n) \cong \tilde{H}_{i-1}(S^{n-1})$. This gives another way of calculating the homology groups of $S^n$ by induction.
\end{example}

\begin{example}[Mayer-Vietoris sequence of Klein bottle]
\label{ex:mayer-vietoris-klein-bottle}
\dcref{ch2:2.47}
\group{Singular Homology}
\uses{thm:excision-general}
We can decompose the Klein bottle $K$ as the union of two Möbius bands $A$ and $B$ glued together by a homeomorphism between their boundary circles.
\end{example}

%% --- END page-0159.tex ---
%% --- BEGIN page-0160.tex ---
Then $A$, $B$, and $A \cap B$ are homotopy equivalent to circles, so the interesting part of the reduced Mayer-Vietoris sequence for the decomposition $K = A \cup B$ is the segment

\[0 \to H_2(K) \to H_1(A \cap B) \xrightarrow{\Phi} H_1(A) \oplus H_1(B) \to H_1(K) \to 0\]

The map $\Phi$ is $\mathbb{Z} \to \mathbb{Z} \oplus \mathbb{Z}$, $1 \mapsto (2, -2)$, since the boundary circle of a M\"{o}bius band wraps twice around the core circle. Since $\Phi$ is injective we obtain $H_2(K) = 0$. Furthermore, we have $H_1(K) \cong \mathbb{Z} \oplus \mathbb{Z}$, since we can choose $(1, 0)$ and $(1, -1)$ as a basis for $\mathbb{Z} \oplus \mathbb{Z}$. All the higher homology groups of $K$ are zero from the earlier part of the Mayer-Vietoris sequence.

\begin{example}[Homology of mapping torus]
\label{ex:homology-mapping-torus}
\dcref{ch2:2.48}
\group{Singular Homology}
\uses{thm:excision-general}
Let us describe an exact sequence which is somewhat similar to the Mayer-Vietoris sequence and which in some cases generalizes it. If we are given two maps $f, g: X \to Y$ then we can form a quotient space $Z$ of the disjoint union of $X \times I$ and $Y$ via the identifications $(x, 0) \sim f(x)$ and $(x, 1) \sim g(x)$, thus attaching one end of $X \times I$ to $Y$ by $f$ and the other end by $g$. For example, if $f$ and $g$ are each the identity map $X \to X$ then $Z = X \times S^1$. If only one of $f$ and $g$, say $f$, is the identity map, then $Z$ is homeomorphic to what is called the mapping torus of $g$, the quotient space of $X \times I$ under the identifications $(x, 0) \sim (g(x), 1)$. The Klein bottle is an example, with $g$ a reflection $S^1 \to S^1$.
\end{example}

The exact sequence we want has the form

\[(*) \quad \cdots \to H_n(X) \xrightarrow{f_* - g_*} H_n(Y) \xrightarrow{i_*} H_n(Z) \to H_{n-1}(X) \xrightarrow{f_* - g_*} H_{n-1}(Y) \to \cdots\]

where $i$ is the evident inclusion $Y \hookrightarrow Z$. To derive this exact sequence, consider the map $q: (X \times I, X \times \partial I) \to (Z, Y)$ that is the restriction to $X \times I$ of the quotient map $X \times I \sqcup Y \to Z$. The map $q$ induces a map of long exact sequences:

\[\begin{array}{ccccccccc}
\cdots & \to & H_{n+1}(X \times I, X \times \partial I) & \xrightarrow{\partial} & H_n(X \times \partial I) & \xrightarrow{i_*} & H_n(X \times I) & \xrightarrow{0} & \cdots \\
& & \downarrow q_* & & \downarrow q_* & & \downarrow q_* & & \\
\cdots & \to & H_{n+1}(Z, Y) & \xrightarrow{\partial} & H_n(Y) & \xrightarrow{i_*} & H_n(Z) & \to & \cdots
\end{array}\]

In the upper row the middle term is the direct sum of two copies of $H_n(X)$, and the map $i_*$ is surjective since $X \times I$ deformation retracts onto $X \times \{0\}$ and $X \times \{1\}$. Surjectivity of the maps $i_*$ in the upper row implies that the next maps are $0$, which in turn implies that the maps $\partial$ are injective. Thus the map $\partial$ in the upper row gives an isomorphism of $H_{n+1}(X \times I, X \times \partial I)$ onto the kernel of $i_*$, which consists of the pairs $(\alpha, -\alpha)$ for $\alpha \in H_n(X)$. This kernel is a copy of $H_n(X)$, and the middle vertical map $q_*$ takes $(\alpha, -\alpha)$ to $f_*(\alpha) - g_*(\alpha)$. The left-hand $q_*$ is an isomorphism since these are good pairs and $q$ induces a homeomorphism of quotient spaces $(X \times I)/(X \times \partial I) = Z/Y$. Hence if we replace $H_{n+1}(Z, Y)$ in the lower exact sequence by the isomorphic group $H_n(X) \cong \ker i_*$ we obtain the long exact sequence we want.

In the case of the mapping torus of a reflection $g: S^1 \to S^1$, with $Z$ a Klein bottle, the interesting portion of the exact sequence $(*)$ is

%% --- END page-0160.tex ---
%% --- BEGIN page-0161.tex ---
\[
0 \to H_2(Z) \to H_1(S^1) \xrightarrow{1-g_*} H_1(S^1) \to H_1(Z) \to H_0(S^1) \xrightarrow{1-g_*} H_0(S^1)
\]

\[
\mathbb{Z} \xrightarrow{2} \mathbb{Z} \quad\quad \mathbb{Z} \xrightarrow{0} \mathbb{Z}
\]

Thus $H_2(Z) = 0$ and we have a short exact sequence $0 \to \mathbb{Z}_2 \to H_1(Z) \to \mathbb{Z} \to 0$. This splits since $\mathbb{Z}$ is free, so $H_1(Z) \cong \mathbb{Z}_2 \oplus \mathbb{Z}$. Other examples are given in the Exercises.

If $Y$ is the disjoint union of spaces $Y_1$ and $Y_2$, with $f: X \to Y_1$ and $g: X \to Y_2$, then $Z$ consists of the mapping cylinders of these two maps with their domain ends identified. For example, suppose we have a CW complex decomposed as the union of two subcomplexes $A$ and $B$ and we take $f$ and $g$ to be the inclusions $A \cap B \hookrightarrow A$ and $A \cap B \hookrightarrow B$. Then the double mapping cylinder $Z$ is homotopy equivalent to $A \cup B$ since we can view $Z$ as $(A \cap B) \times I$ with $A$ and $B$ attached at the two ends, and then slide the attaching of $A$ down to the $B$ end to produce $A \cup B$ with $(A \cap B) \times I$ attached at one of its ends. By Proposition 0.18 the sliding operation preserves homotopy type, so we obtain a homotopy equivalence $Z \simeq A \cup B$. The exact sequence $(\ast)$ in this case is the Mayer-Vietoris sequence.

A relative form of the Mayer-Vietoris sequence is sometimes useful. If one has a pair of spaces $(X, Y) = (A \cup B, C \cup D)$ with $C \subset A$ and $D \subset B$, such that $X$ is the union of the interiors of $A$ and $B$, and $Y$ is the union of the interiors of $C$ and $D$, then there is a relative Mayer-Vietoris sequence

\[
\cdots \to H_n(A \cap B, C \cap D) \xrightarrow{\phi} H_n(A, C) \oplus H_n(B, D) \xrightarrow{\psi} H_n(X, Y) \xrightarrow{\partial} \cdots
\]

To derive this, consider the commutative diagram



where $C_n(A + B, C + D)$ is the quotient of the subgroup $C_n(A + B) \subset C_n(X)$ by its subgroup $C_n(C + D) \subset C_n(Y)$. Thus the three columns of the diagram are exact. We have seen that the first two rows are exact, and we claim that the third row is exact also, with the maps $\varphi$ and $\psi$ induced from the $\varphi$ and $\psi$ in the second row. Since $\psi \varphi = 0$ in the second row, this holds also in the third row, so the third row is at least a chain complex. Viewing the three rows as chain complexes, the diagram then represents a short exact sequence of chain complexes. The associated long exact sequence of homology groups has two out of every three terms zero since the first two rows of the diagram are exact. Hence the remaining homology groups are zero and the third row is exact.

%% --- END page-0161.tex ---
%% --- BEGIN page-0162.tex ---
The third column maps to $0 \to C_n(Y) \to C_n(X) \to C_n(X,Y) \to 0$, inducing maps of homology groups that are isomorphisms for the $X$ and $Y$ terms as we have seen above. So by the five-lemma the maps $C_n(A+B, C+D) \to C_n(X,Y)$ also induce isomorphisms on homology. The relative Mayer-Vietoris sequence is then the long exact sequence of homology groups associated to the short exact sequence of chain complexes given by the third row of the diagram.

\textbf{Homology with Coefficients}

There is an easy generalization of the homology theory we have considered so far that behaves in a very similar fashion and sometimes offers technical advantages. The generalization consists of using chains of the form $\sum_i n_i \sigma_i$ where each $\sigma_i$ is a singular $n$-simplex in $X$ as before, but now the coefficients $n_i$ are taken to lie in a fixed abelian group $G$ rather than $\mathbb{Z}$. Such $n$-chains form an abelian group $C_n(X; G)$, and there is the expected relative version $C_n(X, A; G) = C_n(X; G)/C_n(A; G)$. The old formula for the boundary maps $\partial$ can still be used for arbitrary $G$, namely $\partial(\sum_i n_i \sigma_i) = \sum_i (-1)^j n_i \sigma_i|[v_0, \cdots, \hat{v}_j, \cdots, v_n]$. Just as before, a calculation shows that $\partial^2 = 0$, so the groups $C_n(X; G)$ and $C_n(X, A; G)$ form chain complexes. The resulting homology groups $H_n(X; G)$ and $H_n(X, A; G)$ are called \textbf{homology groups with coefficients in $G$}. Reduced groups $\tilde{H}_n(X; G)$ are defined via the augmented chain complex $\cdots \to C_0(X; G) \xrightarrow{\varepsilon} G \to 0$ with $\varepsilon$ again defined by summing coefficients.

The case $G = \mathbb{Z}_2$ is particularly simple since one is just considering sums of singular simplices with coefficients $0$ or $1$, so by discarding terms with coefficient $0$ one can think of chains as just finite `unions' of singular simplices. The boundary formulas also simplify since one no longer has to worry about signs. Since signs are an algebraic representation of orientation considerations, one can also ignore orientations. This means that homology with $\mathbb{Z}_2$ coefficients is often the most natural tool in the absence of orientability.

All the theory we developed in \S2.1 for $\mathbb{Z}$ coefficients carries over directly to general coefficient groups $G$ with no change in the proofs. The same is true for Mayer-Vietoris sequences. Differences between $H_n(X; G)$ and $H_n(X)$ begin to appear only when one starts making calculations. When $X$ is a point, the method used to compute $H_n(X)$ shows that $H_n(X; G)$ is $G$ for $n = 0$ and $0$ for $n > 0$. From this it follows just as for $G = \mathbb{Z}$ that $\tilde{H}_n(S^k; G)$ is $G$ for $n = k$ and $0$ otherwise.

Cellular homology also generalizes to homology with coefficients, with the cellular chain group $H_n(X^n, X^{n-1})$ replaced by $H_n(X^n, X^{n-1}; G)$, which is a direct sum of $G$'s, one for each $n$-cell. The proof that the cellular homology groups $H_n^{\mathrm{CW}}(X)$ agree with singular homology $H_n(X)$ extends immediately to give $H_n^{\mathrm{CW}}(X; G) \approx H_n(X; G)$. The cellular boundary maps are given by the same formula as for $\mathbb{Z}$ coefficients, $d_n(\sum_\alpha n_\alpha e_\alpha^n) = \sum_{\alpha,\beta} d_{\alpha\beta} n_\alpha e_\beta^{n-1}$. The old proof applies, but the following result is needed to know that the coefficients $d_{\alpha\beta}$ are the same as before:

%% --- END page-0162.tex ---
%% --- BEGIN page-0163.tex ---
\begin{lemma}[Homology map of degree m map with G coefficients]
\label{lem:homology-map-degree-m-coefficients}
\dcref{ch2:2.49}
\group{Singular Homology}
\uses{cor:homology-spheres}
If $f: S^k \to S^k$ has degree $m$, then $f_*: H_k(S^k; G) \to H_k(S^k; G)$ is multiplication by $m$.
\end{lemma}
\begin{proof}
 As a preliminary observation, note that a homomorphism $\varphi: G_1 \to G_2$ induces maps $\varphi_*: C_*(X, A; G_1) \to C_*(X, A; G_2)$ commuting with boundary maps, so there are induced homomorphisms $\varphi_*: H_*(X, A; G_1) \to H_*(X, A; G_2)$. These have various naturality properties. For example, they give a commutative diagram mapping the long exact sequence of homology for the pair $(X, A)$ with $G_1$ coefficients to the corresponding sequence with $G_2$ coefficients. Also, the maps $\varphi_*$ commute with homomorphisms $f_*$ induced by maps $f: (X, A) \to (Y, B)$.

Now let $f: S^k \to S^k$ have degree $m$ and let $\varphi: \mathbb{Z} \to G$ take 1 to a given element $g \in G$. Then we have a commutative


diagram as at the right, where commutativity of the outer two squares comes from the inductive calculation of these homology groups, reducing to the case $k=0$ when the commutativity is obvious.

Since the diagram commutes, the assumption that the map across the top takes 1 to $m$ implies that the map across the bottom takes $g$ to $mg$.
\end{proof}
\begin{example}[Homology of RPn with field coefficients]
\label{ex:homology-rpn-field-coefficients}
\dcref{ch2:2.50}
\group{Singular Homology}
\uses{ex:cellular-homology-rpn}
It is instructive to see what happens to the homology of $\mathbb{R}P^n$ when the coefficient group $G$ is chosen to be a field $F$. The cellular chain complex is
\end{example}

\[\cdots \xrightarrow{0} F \xrightarrow{2} F \xrightarrow{0} F \xrightarrow{2} F \xrightarrow{0} F \to 0\]

Hence if $F$ has characteristic 2, for example if $F = \mathbb{Z}_2$, then $H_k(\mathbb{R}P^n; F) \cong F$ for $0 \le k \le n$, a more uniform answer than with $\mathbb{Z}$ coefficients. On the other hand, if $F$ has characteristic different from 2 then the boundary maps $F \xrightarrow{2} F$ are isomorphisms, hence $H_k(\mathbb{R}P^n; F)$ is $F$ for $k = 0$ and for $k = n$ odd, and is zero otherwise.

In \S3.A we will see that there is a general algebraic formula expressing homology with arbitrary coefficients in terms of homology with $\mathbb{Z}$ coefficients. Some easy special cases that give much of the flavor of the general result are included in the Exercises.

In spite of the fact that homology with $\mathbb{Z}$ coefficients determines homology with other coefficient groups, there are many situations where homology with a suitably chosen coefficient group can provide more information than homology with $\mathbb{Z}$ coefficients. A good example of this is the proof of the Borsuk-Ulam theorem using $\mathbb{Z}_2$ coefficients in \S2.B.

As another illustration, we will now give an example of a map $f: X \to Y$ with the property that the induced maps $f_*$ are trivial for homology with $\mathbb{Z}$ coefficients but not for homology with $\mathbb{Z}_m$ coefficients for suitably chosen $m$. Thus homology with $\mathbb{Z}_m$ coefficients tells us that $f$ is not homotopic to a constant map, which we would not know using only $\mathbb{Z}$ coefficients.

%% --- END page-0163.tex ---
%% --- BEGIN page-0164.tex ---
\begin{example}[Homology of Moore space quotient]
\label{ex:homology-moore-space-quotient}
\dcref{ch2:2.51}
\group{Singular Homology}
\uses{ex:moore-spaces}
Let $X$ be a Moore space $M(\mathbb{Z}_m, n)$ obtained from $S^n$ by attaching a cell $e^{n+1}$ by a map of degree $m$. The quotient map $f:X \to X/S^n = S^{n+1}$ induces trivial homomorphisms on reduced homology with $\mathbb{Z}$ coefficients since the nonzero reduced homology groups of $X$ and $S^{n+1}$ occur in different dimensions. But with $\mathbb{Z}_m$ coefficients the story is different, as we can see by considering the long exact sequence of the pair $(X, S^n)$, which contains the segment
\[
0 = \widetilde{H}_{n+1}(S^n;\mathbb{Z}_m) \to \widetilde{H}_{n+1}(X;\mathbb{Z}_m) \xrightarrow{f_*} \widetilde{H}_{n+1}(X/S^n;\mathbb{Z}_m)
\]
Exactness says that $f_*$ is injective, hence nonzero since $\widetilde{H}_{n+1}(X;\mathbb{Z}_m)$ is $\mathbb{Z}_m$, the cellular boundary map $H_{n+1}(X^{n+1}, X^n;\mathbb{Z}_m) \to H_n(X^n, X^{n-1};\mathbb{Z}_m)$ being $\mathbb{Z}_m \xrightarrow{m} \mathbb{Z}_m$.
\end{example}

\section*{Exercises}

%ORIGINAL-NUMBER: 1
\begin{exercise}
Prove the Brouwer fixed point theorem for maps $f:D^n \to D^n$ by applying degree theory to the map $S^n \to S^n$ that sends both the northern and southern hemispheres of $S^n$ to the southern hemisphere via $f$. [This was Brouwer's original proof.]
\end{exercise}

%ORIGINAL-NUMBER: 2
\begin{exercise}
Given a map $f:S^{2n} \to S^{2n}$, show that there is some point $x \in S^{2n}$ with either $f(x) = x$ or $f(x) = -x$. Deduce that every map $\mathbb{RP}^{2n} \to \mathbb{RP}^{2n}$ has a fixed point. Construct maps $\mathbb{RP}^{2n-1} \to \mathbb{RP}^{2n-1}$ without fixed points from linear transformations $\mathbb{R}^{2n} \to \mathbb{R}^{2n}$ without eigenvectors.
\end{exercise}

%ORIGINAL-NUMBER: 3
\begin{exercise}
Let $f:S^n \to S^n$ be a map of degree zero. Show that there exist points $x,y \in S^n$ with $f(x) = x$ and $f(y) = -y$. Use this to show that if $F$ is a continuous vector field defined on the unit ball $D^n$ in $\mathbb{R}^n$ such that $F(x) \neq 0$ for all $x$, then there exists a point on $\partial D^n$ where $F$ points radially outward and another point on $\partial D^n$ where $F$ points radially inward.
\end{exercise}

%ORIGINAL-NUMBER: 4
\begin{exercise}
Construct a surjective map $S^n \to S^n$ of degree zero, for each $n \geq 1$.
\end{exercise}

%ORIGINAL-NUMBER: 5
\begin{exercise}
Show that any two reflections of $S^n$ across different $n$-dimensional hyperplanes are homotopic, in fact homotopic through reflections. [The linear algebra formula for a reflection in terms of inner products may be helpful.]
\end{exercise}

%ORIGINAL-NUMBER: 6
\begin{exercise}
Show that every map $S^n \to S^n$ can be homotoped to have a fixed point if $n > 0$.
\end{exercise}

%ORIGINAL-NUMBER: 7
\begin{exercise}
For an invertible linear transformation $f:\mathbb{R}^n \to \mathbb{R}^n$ show that the induced map on $H_n(\mathbb{R}^n, \mathbb{R}^n - \{0\}) \cong \widetilde{H}_{n-1}(\mathbb{R}^n - \{0\}) \cong \mathbb{Z}$ is $\pm 1$ according to whether the determinant of $f$ is positive or negative. [Use Gaussian elimination to show that the matrix of $f$ can be joined by a path of invertible matrices to a diagonal matrix with $\pm 1$'s on the diagonal.]
\end{exercise}

%ORIGINAL-NUMBER: 8
\begin{exercise}
A polynomial $f(z)$ with complex coefficients, viewed as a map $\mathbb{C} \to \mathbb{C}$, can always be extended to a continuous map of one-point compactifications $\bar{f}:S^2 \to S^2$. Show that the degree of $\bar{f}$ equals the degree of $f$ as a polynomial. Show also that the local degree of $\bar{f}$ at a root of $f$ is the multiplicity of the root.
\end{exercise}

%% --- END page-0164.tex ---
%% --- BEGIN page-0165.tex ---
%ORIGINAL-NUMBER: 9
\begin{exercise}
Compute the homology groups of the following 2-complexes:
\begin{enumerate}
\item[(a)] The quotient of $S^2$ obtained by identifying north and south poles to a point.
\item[(b)] $S^1 \times (S^1 \vee S^1)$.
\item[(c)] The space obtained from $D^2$ by first deleting the interiors of two disjoint subdisks in the interior of $D^2$ and then identifying all three resulting boundary circles together via homeomorphisms preserving clockwise orientations of these circles.
\item[(d)] The quotient space of $S^1 \times S^1$ obtained by identifying points in the circle $S^1 \times \{x_0\}$ that differ by $2\pi/m$ rotation and identifying points in the circle $\{x_0\} \times S^1$ that differ by $2\pi/n$ rotation.
\end{enumerate}
\end{exercise}

%ORIGINAL-NUMBER: 10
\begin{exercise}
Let $X$ be the quotient space of $S^2$ under the identifications $x \sim -x$ for $x$ in the equator $S^1$. Compute the homology groups $H_r(X)$. Do the same for $S^3$ with antipodal points of the equatorial $S^2 \subset S^3$ identified.
\end{exercise}

%ORIGINAL-NUMBER: 11
\begin{exercise}
In an exercise for \S 1.2 we described a 3-dimensional CW complex obtained from the cube $I^3$ by identifying opposite faces via a one-quarter twist. Compute the homology groups of this complex.
\end{exercise}

%ORIGINAL-NUMBER: 12
\begin{exercise}
Show that the quotient map $S^1 \times S^1 \to S^2$ collapsing the subspace $S^1 \vee S^1$ to a point is not nullhomotopic by showing that it induces an isomorphism on $H_2$. On the other hand, show via covering spaces that any map $S^2 \to S^1 \times S^1$ is nullhomotopic.
\end{exercise}

%ORIGINAL-NUMBER: 13
\begin{exercise}
Let $X$ be the 2-complex obtained from $S^1$ with its usual cell structure by attaching two 2-cells by maps of degrees 2 and 3, respectively.
\begin{enumerate}
\item[(a)] Compute the homology groups of all the subcomplexes $A \subset X$ and the corresponding quotient complexes $X/A$.
\item[(b)] Show that $X \simeq S^2$ and that the only subcomplex $A \subset X$ for which the quotient map $X \to X/A$ is a homotopy equivalence is the trivial subcomplex, the 0-cell.
\end{enumerate}
\end{exercise}

%ORIGINAL-NUMBER: 14
\begin{exercise}
A map $f:S^n \to S^n$ satisfying $f(x) = f(-x)$ for all $x$ is called an even map. Show that an even map $S^n \to S^n$ must have even degree, and that the degree must in fact be zero when $n$ is even. When $n$ is odd, show there exist even maps of any given even degree. [Hints: If $f$ is even, it factors as a composition $S^n \to \mathbb{RP}^n \to S^n$. Using the calculation of $H_n(\mathbb{RP}^n)$ in the text, show that the induced map $H_n(S^n) \to H_n(\mathbb{RP}^n)$ sends a generator to twice a generator when $n$ is odd. It may be helpful to show that the quotient map $\mathbb{RP}^n \to \mathbb{RP}^n/\mathbb{RP}^{n-1}$ induces an isomorphism on $H_n$ when $n$ is odd.]
\end{exercise}

%ORIGINAL-NUMBER: 15
\begin{exercise}
Show that if $X$ is a CW complex then $H_n(X^n)$ is free by identifying it with the kernel of the cellular boundary map $H_n(X^n, X^{n-1}) \to H_{n-1}(X^{n-1}, X^{n-2})$.
\end{exercise}

%ORIGINAL-NUMBER: 16
\begin{exercise}
Let $\Delta^n = [v_0, \cdots, v_n]$ have its natural $\Delta$-complex structure with $k$-simplices $[v_{i_0}, \cdots, v_{i_k}]$ for $i_0 < \cdots < i_k$. Compute the ranks of the simplicial (or cellular) chain groups $\Delta_c(\Delta^n)$ and the subgroups of cycles and boundaries. [Hint: Pascal's triangle.] Apply this to show that the $k$-skeleton of $\Delta^n$ has homology groups $\widetilde{H}_i((\Delta^n)^k)$ equal to $0$ for $i < k$, and free of rank $\binom{n}{k+1}$ for $i = k$.
\end{exercise}

%% --- END page-0165.tex ---
%% --- BEGIN page-0166.tex ---
%ORIGINAL-NUMBER: 17
\begin{exercise}
Show the isomorphism between cellular and singular homology is natural in the following sense: A map $f:X \to Y$ that is \textit{cellular} --- satisfying $f(X^n) \subset Y^n$ for all $n$ --- induces a chain map $f_*$ between the cellular chain complexes of $X$ and $Y$, and the map $f_*: H_n^{CW}(X) \to H_n^{CW}(Y)$ induced by this chain map corresponds to $f_*: H_n(X) \to H_n(Y)$ under the isomorphism $H_n^{CW} \cong H_n$.
\end{exercise}

%ORIGINAL-NUMBER: 18
\begin{exercise}
For a CW pair $(X,A)$ show there is a relative cellular chain complex formed by the groups $H_i(X^i, X^{i-1} \cup A^i)$, having homology groups isomorphic to $H_n(X,A)$.
\end{exercise}

%ORIGINAL-NUMBER: 19
\begin{exercise}
Compute $H_i(\mathbb{R}P^n / \mathbb{R}P^m)$ for $m < n$ by cellular homology, using the standard CW structure on $\mathbb{R}P^n$ with $\mathbb{R}P^m$ as its $m$-skeleton.
\end{exercise}

%ORIGINAL-NUMBER: 20
\begin{exercise}
For finite CW complexes $X$ and $Y$, show that $\chi(X \times Y) = \chi(X) \chi(Y)$.
\end{exercise}

%ORIGINAL-NUMBER: 21
\begin{exercise}
If a finite CW complex $X$ is the union of subcomplexes $A$ and $B$, show that $\chi(X) = \chi(A) + \chi(B) - \chi(A \cap B)$.
\end{exercise}

%ORIGINAL-NUMBER: 22
\begin{exercise}
For $X$ a finite CW complex and $p: \tilde{X} \to X$ an $n$-sheeted covering space, show that $\chi(\tilde{X}) = n \chi(X)$.
\end{exercise}

%ORIGINAL-NUMBER: 23
\begin{exercise}
Show that if the closed orientable surface $M_g$ of genus $g$ is a covering space of $M_h$, then $g = n(h-1) + 1$ for some $n$, namely, $n$ is the number of sheets in the covering. [Conversely, if $g = n(h-1) + 1$ then there is an $n$-sheeted covering $M_g \to M_h$, as we saw in Example 1.41.]
\end{exercise}

%ORIGINAL-NUMBER: 24
\begin{exercise}
Suppose we build $S^2$ from a finite collection of polygons by identifying edges in pairs. Show that in the resulting CW structure on $S^2$ the 1-skeleton cannot be either of the two graphs shown, with five and six vertices. [This is one step in a proof that neither of these graphs embeds in $\mathbb{R}^2$.]


\end{exercise}

%ORIGINAL-NUMBER: 25
\begin{exercise}
Show that for each $n \in \mathbb{Z}$ there is a unique function $\varphi$ assigning an integer to each finite CW complex, such that (a) $\varphi(X) = \varphi(Y)$ if $X$ and $Y$ are homeomorphic, (b) $\varphi(X) = \varphi(A) + \varphi(X/A)$ if $A$ is a subcomplex of $X$, and (c) $\varphi(S^0) = n$. For such a function $\varphi$, show that $\varphi(X) = \varphi(Y)$ if $X = Y$.
\end{exercise}

%ORIGINAL-NUMBER: 26
\begin{exercise}
For a pair $(X,A)$, let $X \cup CA$ be $X$ with a cone on $A$ attached.

(a) Show that $X$ is a retract of $X \cup CA$ iff $A$ is \textit{contractible in} $X$: There is a homotopy $f_t: A \to X$ with $f_0$ the inclusion $A \hookrightarrow X$ and $f_1$ a constant map.

(b) Show that if $A$ is contractible in $X$ then $H_n(X,A) \cong \tilde{H}_n(X) \oplus \tilde{H}_{n-1}(A)$, using the fact that $(X \cup CA)/X$ is the suspension $SA$ of $A$.
\end{exercise}

%ORIGINAL-NUMBER: 27
\begin{exercise}
The short exact sequences $0 \to C_n(A) \to C_n(X) \to C_n(X,A) \to 0$ always split, but why does this not always yield splittings $H_n(X) \cong H_n(A) \oplus H_n(X,A)$?
\end{exercise}

%ORIGINAL-NUMBER: 28
\begin{exercise}
(a) Use the Mayer-Vietoris sequence to compute the homology groups of the space obtained from a torus $S^1 \times S^1$ by attaching a Möbius band via a homeomorphism from the boundary circle of the Möbius band to the circle $S^1 \times \{x_0\}$ in the torus.

(b) Do the same for the space obtained by attaching a Möbius band to $\mathbb{R}P^2$ via a homeomorphism of its boundary circle to the standard $\mathbb{R}P^1 \subset \mathbb{R}P^2$.
\end{exercise}

%% --- END page-0166.tex ---
%% --- BEGIN page-0167.tex ---
%ORIGINAL-NUMBER: 29
\begin{exercise}
The surface $M_g$ of genus $g$, embedded in $\mathbb{R}^3$ in the standard way, bounds a compact region $R$. Two copies of $R$, glued together by the identity map between their boundary surfaces $M_g$, form a closed 3-manifold $X$. Compute the homology groups of $X$ via the Mayer-Vietoris sequence for this decomposition of $X$ into two copies of $R$. Also compute the relative groups $H_i(R, M_g)$.
\end{exercise}

%ORIGINAL-NUMBER: 30
\begin{exercise}
For the mapping torus $T_f$ of a map $f: X \to X$, we constructed in Example 2.48 a long exact sequence $\cdots \to H_n(X) \xrightarrow{1-f_*} H_n(X) \to H_n(T_f) \to H_{n-1}(X) \to \cdots$. Use this to compute the homology of the mapping tori of the following maps:

(a) A reflection $S^2 \to S^2$.

(b) A map $S^2 \to S^2$ of degree 2.

(c) The map $S^1 \times S^1 \to S^1 \times S^1$ that is the identity on one factor and a reflection on the other.

(d) The map $S^1 \times S^1 \to S^1 \times S^1$ that is a reflection on each factor.

(e) The map $S^1 \times S^1 \to S^1 \times S^1$ that interchanges the two factors and then reflects one of the factors.
\end{exercise}

%ORIGINAL-NUMBER: 31
\begin{exercise}
Use the Mayer-Vietoris sequence to show there are isomorphisms $\tilde{H}_i(X \vee Y) \cong \tilde{H}_i(X) \oplus \tilde{H}_i(Y)$ if the basepoints of $X$ and $Y$ that are identified in $X \vee Y$ are deformation retracts of neighborhoods $U \subset X$ and $V \subset Y$.
\end{exercise}

%ORIGINAL-NUMBER: 32
\begin{exercise}
For $SX$ the suspension of $X$, show by a Mayer-Vietoris sequence that there are isomorphisms $\tilde{H}_n(SX) \cong \tilde{H}_{n-1}(X)$ for all $n$.
\end{exercise}

%ORIGINAL-NUMBER: 33
\begin{exercise}
Suppose the space $X$ is the union of open sets $A_1, \cdots, A_n$ such that each intersection $A_{i_1} \cap \cdots \cap A_{i_k}$ is either empty or has trivial reduced homology groups. Show that $\tilde{H}_i(X) = 0$ for $i \geq n - 1$, and give an example showing this inequality is best possible, for each $n$.
\end{exercise}

%ORIGINAL-NUMBER: 34
\begin{exercise}
[Deleted --- see the errata for comments.]
\end{exercise}

%ORIGINAL-NUMBER: 35
\begin{exercise}
Use the Mayer-Vietoris sequence to show that a nonorientable closed surface, or more generally a finite simplicial complex $X$ for which $H_1(X)$ contains torsion, cannot be embedded as a subspace of $\mathbb{R}^3$ in such a way as to have a neighborhood homeomorphic to the mapping cylinder of some map from a closed orientable surface to $X$. [This assumption on a neighborhood is in fact not needed if one deduces the result from Alexander duality in $\S 3.3$.]
\end{exercise}

%ORIGINAL-NUMBER: 36
\begin{exercise}
Show that $H_i(X \times S^n) \cong H_i(X) \oplus H_{i-n}(X)$ for all $i$ and $n$, where $H_i = 0$ for $i < 0$ by definition. Namely, show $H_i(X \times S^n) \cong H_i(X) \oplus H_i(X \times S^n, X \times \{x_0\})$ and $H_j(X \times S^n, X \times \{x_0\}) \cong H_{j-1}(X \times S^{n-1}, X \times \{x_0\})$. [For the latter isomorphism the relative Mayer-Vietoris sequence yields an easy proof.]
\end{exercise}

%ORIGINAL-NUMBER: 37
\begin{exercise}
Give an elementary derivation for the Mayer-Vietoris sequence in simplicial homology for a $\Delta$-complex $X$ decomposed as the union of subcomplexes $A$ and $B$.
\end{exercise}

%% --- END page-0167.tex ---
%% --- BEGIN page-0168.tex ---
%ORIGINAL-NUMBER: 38
\begin{exercise}
Show that a commutative diagram
\[
\begin{array}{ccccccccc}
\cdots \to & C_{n+1} & \to & B_n & \to & C_n & \to & B_{n-1} & \to \cdots \\
& \downarrow & \searrow & A_n & \downarrow & \downarrow & \searrow & A_{n-1} & \\
\cdots \to & E_{n+1} & \to & D_n & \to & E_n & \to & D_{n-1} & \to \cdots
\end{array}
\]
with the two sequences across the top and bottom exact, gives rise to an exact sequence $\cdots \to E_{n+1} \to B_n \to C_n \oplus D_n \to E_n \to B_{n-1} \to \cdots$ where the maps are obtained from those in the previous diagram in the obvious way, except that $B_n \to C_n \oplus D_n$ has a minus sign in one coordinate.
\end{exercise}

%ORIGINAL-NUMBER: 39
\begin{exercise}
Use the preceding exercise to derive relative Mayer-Vietoris sequences for CW pairs $(X, Y) = (A \cup B, C \cup D)$ with $A = B$ or $C = D$.
\end{exercise}

%ORIGINAL-NUMBER: 40
\begin{exercise}
From the long exact sequence of homology groups associated to the short exact sequence of chain complexes $0 \to C_i(X) \xrightarrow{n} C_i(X) \to C_i(X; \mathbb{Z}_n) \to 0$ deduce immediately that there are short exact sequences
\[
0 \to H_i(X) / nH_i(X) \to H_i(X; \mathbb{Z}_n) \to n\text{-Torsion}(H_{i-1}(X)) \to 0
\]
where $n\text{-Torsion}(G)$ is the kernel of the map $G \xrightarrow{n} G$, $g \mapsto ng$. Use this to show that $\tilde{H}_i(X; \mathbb{Z}_p) = 0$ for all $i$ and all primes $p$ iff $\tilde{H}_i(X)$ is a vector space over $\mathbb{Q}$ for all $i$.
\end{exercise}

%ORIGINAL-NUMBER: 41
\begin{exercise}
For $X$ a finite CW complex and $F$ a field, show that the Euler characteristic $\chi(X)$ can also be computed by the formula $\chi(X) = \sum_n (-1)^n \dim H_n(X; F)$, the alternating sum of the dimensions of the vector spaces $H_n(X; F)$.
\end{exercise}

%ORIGINAL-NUMBER: 42
\begin{exercise}
Let $X$ be a finite connected graph having no vertex that is the endpoint of just one edge, and suppose that $H_1(X; \mathbb{Z})$ is free abelian of rank $n > 1$, so the group of automorphisms of $H_1(X; \mathbb{Z})$ is $GL_n(\mathbb{Z})$, the group of invertible $n \times n$ matrices with integer entries whose inverse matrix also has integer entries. Show that if $G$ is a finite group of homeomorphisms of $X$, then the homomorphism $G \to GL_n(\mathbb{Z})$ assigning to $g: X \to X$ the induced homomorphism $g_*: H_1(X; \mathbb{Z}) \to H_1(X; \mathbb{Z})$ is injective. Show the same result holds if the coefficient group $\mathbb{Z}$ is replaced by $\mathbb{Z}_m$ with $m > 2$. What goes wrong when $m = 2$?
\end{exercise}

%ORIGINAL-NUMBER: 43
\begin{exercise}
(a) Show that a chain complex of free abelian groups $C_n$ splits as a direct sum of subcomplexes $0 \to L_{n+1} \to K_n \to 0$ with at most two nonzero terms. [Show the short exact sequence $0 \to \text{Ker}\partial \to C_n \to \text{Im}\partial \to 0$ splits and take $K_n = \text{Ker}\partial$.]

(b) In case the groups $C_n$ are finitely generated, show there is a further splitting into summands $0 \to \mathbb{Z} \to 0$ and $0 \to \mathbb{Z} \xrightarrow{m} \mathbb{Z} \to 0$. [Reduce the matrix of the boundary map $L_{n+1} \to K_n$ to echelon form by elementary row and column operations.]

(c) Deduce that if $X$ is a CW complex with finitely many cells in each dimension, then $H_n(X; G)$ is the direct sum of the following groups:
\begin{itemize}
\item a copy of $G$ for each $\mathbb{Z}$ summand of $H_n(X)$
\item a copy of $G / mG$ for each $\mathbb{Z}_m$ summand of $H_n(X)$
\item a copy of the kernel of $G \xrightarrow{m} G$ for each $\mathbb{Z}_m$ summand of $H_{n-1}(X)$
\end{itemize}
\end{exercise}

%% --- END page-0168.tex ---
%% --- BEGIN page-0169.tex ---
\section{The Formal Viewpoint}\label{sec:formal-viewpoint}

Sometimes it is good to step back from the forest of details and look for general patterns. In this rather brief section we will first describe the general pattern of homology by axioms, then we will look at some common formal features shared by many of the constructions we have made, using the language of categories and functors which has become common in much of modern mathematics.

\textbf{Axioms for Homology}

For simplicity let us restrict attention to CW complexes and focus on reduced homology to avoid mentioning relative homology. A (reduced) \textbf{homology theory} assigns to each nonempty CW complex $X$ a sequence of abelian groups $\tilde{h}_n(X)$ and to each map $f : X \to Y$ between CW complexes a sequence of homomorphisms $f_* : \tilde{h}_n(X) \to \tilde{h}_n(Y)$ such that $(fg)_* = f_* g_*$ and $\mathbb{1}_* = \mathbb{1}$, and so that the following three axioms are satisfied.

\begin{enumerate}
\item If $f = g : X \to Y$, then $f_* = g_* : \tilde{h}_n(X) \to \tilde{h}_n(Y)$.

\item There are boundary homomorphisms $\partial : \tilde{h}_n(X/A) \to \tilde{h}_{n-1}(A)$ defined for each CW pair $(X,A)$, fitting into an exact sequence
\[ \cdots \xrightarrow{\partial} \tilde{h}_n(A) \xrightarrow{i_*} \tilde{h}_n(X) \xrightarrow{q_*} \tilde{h}_n(X/A) \xrightarrow{\partial} \tilde{h}_{n-1}(A) \xrightarrow{i_*} \cdots \]
where $i$ is the inclusion and $q$ is the quotient map. Furthermore the boundary maps are natural: For $f : (X,A) \to (Y,B)$ inducing a quotient map $\bar{f} : X/A \to Y/B$, there are commutative diagrams
\[
\begin{array}{ccc}
\tilde{h}_n(X/A) & \xrightarrow{\partial} & \tilde{h}_{n-1}(A) \\
\downarrow \bar{f}_* & & \downarrow f_* \\
\tilde{h}_n(Y/B) & \xrightarrow{\partial} & \tilde{h}_{n-1}(B)
\end{array}
\]

\item For a wedge sum $X = \bigvee_\alpha X_\alpha$ with inclusions $i_\alpha : X_\alpha \hookrightarrow X$, the direct sum map $\bigoplus_\alpha i_{\alpha *} : \bigoplus_\alpha \tilde{h}_n(X_\alpha) \to \tilde{h}_n(X)$ is an isomorphism for each $n$.
\end{enumerate}

Negative values for the subscripts $n$ are permitted. Ordinary singular homology is zero in negative dimensions by definition, but interesting homology theories with nontrivial groups in negative dimensions do exist.

The third axiom may seem less substantial than the first two, and indeed for finite wedge sums it can be deduced from the first two axioms, though not in general for infinite wedge sums, as an example in the Exercises shows.

It is also possible, and not much more difficult, to give axioms for unreduced homology theories. One supposes one has relative groups $h_n(X,A)$ defined, specializing to absolute groups by setting $h_n(X) = h_n(X, \emptyset)$. Axiom (1) is replaced by its

%% --- END page-0169.tex ---
%% --- BEGIN page-0170.tex ---
obvious relative form, and axiom (2) is broken into two parts, the first hypothesizing a long exact sequence involving these relative groups, with natural boundary maps, the second stating some version of excision, for example $h_n(X, A) \approx h_n(X/A, A/A)$ if one is dealing with CW pairs. In axiom (3) the wedge sum is replaced by disjoint union.

These axioms for unreduced homology are essentially the same as those originally laid out in the highly influential book [Eilenberg \& Steenrod 1952], except that axiom (3) was omitted since the focus there was on finite complexes, and there was another axiom specifying that the groups $h_n(\text{point})$ are zero for $n \neq 0$, as is true for singular homology. This axiom was called the `dimension axiom', presumably because it specifies that a point has nontrivial homology only in dimension zero. It can be regarded as a normalization axiom, since one can trivially define a homology theory where it fails by setting $h_n(X, A) = H_{n+k}(X, A)$ for a fixed nonzero integer $k$. At the time there were no interesting homology theories known for which the dimension axiom did not hold, but soon thereafter topologists began studying a homology theory called `bordism' having the property that the bordism groups of a point are nonzero in infinitely many dimensions. Axiom (3) seems to have appeared first in [Milnor 1962].

Reduced and unreduced homology theories are essentially equivalent. From an unreduced theory $h$ one gets a reduced theory $\tilde{h}$ by setting $\tilde{h}_n(X)$ equal to the kernel of the canonical map $h_n(X) \to h_n(\text{point})$. In the other direction, one sets $h_n(X) = \tilde{h}_n(X_+)$ where $X_+$ is the disjoint union of $X$ with a point. We leave it as an exercise to show that these two transformations between reduced and unreduced homology are inverses of each other. Just as with ordinary homology, one has $h_n(X) \approx \tilde{h}_n(X) \oplus h_n(x_0)$ for any point $x_0 \in X$, since the long exact sequence of the pair $(X, x_0)$ splits via the retraction of $X$ onto $x_0$. Note that $\tilde{h}_n(x_0) = 0$ for all $n$, as can be seen by looking at the long exact sequence of reduced homology groups of the pair $(x_0, x_0)$.

The groups $h_n(x_0) \approx \tilde{h}_n(S^0)$ are called the \textbf{coefficients} of the homology theories $h$ and $\tilde{h}$, by analogy with the case of singular homology with coefficients. One can trivially realize any sequence of abelian groups $G_i$ as the coefficient groups of a homology theory by setting $h_n(X, A) = \bigoplus_i H_{n-i}(X, A; G_i)$.

In general, homology theories are not uniquely determined by their coefficient groups, but this is true for singular homology: If $h$ is a homology theory defined for CW pairs, whose coefficient groups $h_n(x_0)$ are zero for $n \neq 0$, then there are natural isomorphisms $h_n(X, A) \approx H_n(X, A; G)$ for all CW pairs $(X, A)$ and all $n$, where $G = h_0(x_0)$. This will be proved in Theorem 4.59.

We have seen how Mayer-Vietoris sequences can be quite useful for singular homology, and in fact every homology theory has Mayer-Vietoris sequences, at least for CW complexes. These can be obtained directly from the axioms in the follow-

%% --- END page-0170.tex ---
%% --- BEGIN page-0171.tex ---
ing way. For a CW complex $X = A \cup B$ with $A$ and $B$ subcomplexes, the inclusion
$(B, A \cap B) \to (X, A)$ induces a commutative diagram of exact sequences

\[
\cdots \to h_{n+1}(B, A \cap B) \to h_n(A \cap B) \to h_n(B) \to h_n(B, A \cap B) \to \cdots
\]
\[
\Big\downarrow \cong \quad \Big\downarrow \quad \Big\downarrow \quad \Big\downarrow \cong
\]
\[
\cdots \to h_{n+1}(X, A) \to h_n(A) \to h_n(X) \to h_n(X, A) \to \cdots
\]

The vertical maps between relative groups are isomorphisms since $B/(A \cap B) = X/A$.
Then it is a purely algebraic fact, whose proof is Exercise 38 at the end of the previous
section, that a diagram such as this with every third vertical map an isomorphism
gives rise to a long exact sequence involving the remaining nonisomorphic terms. In
the present case this takes the form of a Mayer-Vietoris sequence

\[
\cdots \to h_n(A \cap B) \xrightarrow{\varphi} h_n(A) \oplus h_n(B) \xrightarrow{\psi} h_n(X) \xrightarrow{\partial} h_{n-1}(A \cap B) \to \cdots
\]

\textbf{Categories and Functors}

Formally, singular homology can be regarded as a sequence of functions $H_n$ that
assign to each space $X$ an abelian group $H_n(X)$ and to each map $f : X \to Y$ a homomorphism $H_n(f) = f_* : H_n(X) \to H_n(Y)$, and similarly for relative homology groups.
This sort of situation arises quite often, and not just in algebraic topology, so it is
useful to introduce some general terminology for it. Roughly speaking, `functions'
like $H_n$ are called `functors', and the domains and ranges of these functors are called
`categories'. Thus for $H_n$ the domain category consists of topological spaces and continuous maps, or in the relative case, pairs of spaces and continuous maps of pairs,
and the range category consists of abelian groups and homomorphisms. A key point
is that one is interested not only in the objects in the category, for example spaces or
groups, but also in the maps, or `morphisms', between these objects.

Now for the precise definitions. A \textbf{category} $\mathcal{C}$ consists of three things:

\begin{enumerate}
\item A collection Ob$(\mathcal{C})$ of \textbf{objects}.
\item Sets Mor$(X, Y)$ of \textbf{morphisms} for each pair $X, Y \in \text{Ob}(\mathcal{C})$, including a distinguished `identity' morphism $\mathbb{1} = \mathbb{1}_X \in \text{Mor}(X, X)$ for each $X$.
\item A `composition of morphisms' function $\circ : \text{Mor}(X, Y) \times \text{Mor}(Y, Z) \to \text{Mor}(X, Z)$ for each triple $X, Y, Z \in \text{Ob}(\mathcal{C})$, satisfying $f \circ \mathbb{1} = f$, $\mathbb{1} \circ f = f$, and $(f \circ g) \circ h = f \circ (g \circ h)$.
\end{enumerate}

There are plenty of obvious examples, such as:

\begin{itemize}
\item The category of topological spaces, with continuous maps as the morphisms. Or
we could restrict to special classes of spaces such as CW complexes, keeping
continuous maps as the morphisms. We could also restrict the morphisms, for
example to homeomorphisms.
\item The category of groups, with homomorphisms as morphisms. Or the subcategory
of abelian groups, again with homomorphisms as the morphisms. Generalizing

%% --- END page-0171.tex ---
%% --- BEGIN page-0172.tex ---
this is the category of modules over a fixed ring, with morphisms the module homomorphisms.


\item The category of sets, with arbitrary functions as the morphisms. Or the morphisms could be restricted to injections, surjections, or bijections.
\end{itemize}

There are also many categories where the morphisms are not simply functions, for example:

\begin{itemize}
\item Any group $G$ can be viewed as a category with only one object and with $G$ as the morphisms of this object, so that condition (3) reduces to two of the three axioms for a group. If we require only these two axioms, associativity and a left and right identity, we have a `group without inverses', usually called a \emph{monoid} since it is the same thing as a category with one object.

\item A partially ordered set $(X, \le)$ can be considered a category where the objects are the elements of $X$ and there is a unique morphism from $x$ to $y$ whenever $x \le y$. The relation $x \le x$ gives the morphism $\mathrm{id}$ and transitivity gives the composition $\mathrm{Mor}(x,y) \times \mathrm{Mor}(y,z) \to \mathrm{Mor}(x,z)$. The condition that $x \le y$ and $y \le x$ implies $x = y$ says that there is at most one morphism between any two objects.

\item There is a `homotopy category' whose objects are topological spaces and whose morphisms are homotopy classes of maps, rather than actual maps. This uses the fact that composition is well-defined on homotopy classes: $f_0 g_0 = f_1 g_1$ if $f_0 \simeq f_1$ and $g_0 \simeq g_1$.

\item Chain complexes are the objects of a category, with chain maps as morphisms. This category has various interesting subcategories, obtained by restricting the objects. For example, we could take chain complexes whose groups are zero in negative dimensions, or zero outside a finite range. Or we could restrict to exact sequences, or short exact sequences. In each case we take morphisms to be chain maps, which are commutative diagrams. Going a step further, there is a category whose objects are short exact sequences of chain complexes and whose morphisms are commutative diagrams of maps between such short exact sequences.
\end{itemize}

A functor $F$ from a category $\mathcal{C}$ to a category $\mathcal{D}$ assigns to each object $X$ in $\mathcal{C}$ an object $F(X)$ in $\mathcal{D}$ and to each morphism $f \in \mathrm{Mor}(X,Y)$ in $\mathcal{C}$ a morphism $F(f) \in \mathrm{Mor}(F(X), F(Y))$ in $\mathcal{D}$, such that $F(\mathrm{id}) = \mathrm{id}$ and $F(f \circ g) = F(f) \circ F(g)$. In the case of the singular homology functor $H_n$, the latter two conditions are the familiar properties $\mathrm{id}_* = \mathrm{id}$ and $(fg)_* = f_* g_*$ of induced maps. Strictly speaking, what we have just defined is a \textbf{covariant functor}. A \textbf{contravariant functor} would differ from this by assigning to $f \in \mathrm{Mor}(X,Y)$ a `backwards' morphism $F(f) \in \mathrm{Mor}(F(Y), F(X))$ with $F(\mathrm{id}) = \mathrm{id}$ and $F(f \circ g) = F(g) \circ F(f)$. A classical example of this is the dual vector space functor, which assigns to a vector space $V$ over a fixed scalar field $K$ the dual vector space $F(V) = V^*$ of linear maps $V \to K$, and to each linear transformation
%% --- END page-0172.tex ---
%% --- BEGIN page-0173.tex ---
$f:V \to W$ the dual map $F(f) = f^*: W^* \to V^*$, going in the reverse direction. In the next chapter we will study the contravariant version of homology, called cohomology.

A number of the constructions we have studied in this chapter are functors:

\begin{itemize}
\item The singular chain complex functor assigns to a space $X$ the chain complex of singular chains in $X$ and to a map $f:X \to Y$ the induced chain map.  This is a functor from the category of spaces and continuous maps to the category of chain complexes and chain maps.

\item The algebraic homology functor assigns to a chain complex its sequence of homology groups and to a chain map the induced homomorphisms on homology. This is a functor from the category of chain complexes and chain maps to the category whose objects are sequences of abelian groups and whose morphisms are sequences of homomorphisms.

\item The composition of the two preceding functors is the functor assigning to a space its singular homology groups.

\item The first example above, the singular chain complex functor, can itself be regarded as the composition of two functors. The first functor assigns to a space $X$ its singular complex $S(X)$, a $\Delta$-complex, and the second functor assigns to a $\Delta$-complex its simplicial chain complex. This is what the two functors do on objects, and what they do on morphisms can be described in the following way. A map of spaces $f:X \to Y$ induces a map $f_*: S(X) \to S(Y)$ by composing singular simplices $\Delta^n \to X$ with $f$. The map $f_*$ is a map between $\Delta$-complexes taking the distinguished characteristic maps in the domain $\Delta$-complex to the distinguished characteristic maps in the target $\Delta$-complex. Call such maps $\mathbf{\Delta}$\textbf{-maps} and let them be the morphisms in the category of $\Delta$-complexes. Note that a $\Delta$-map induces a chain map between simplicial chain complexes, taking basis elements to basis elements, so we have a simplicial chain complex functor taking the category of $\Delta$-complexes and $\Delta$-maps to the category of chain complexes and chain maps.

\item There is a functor assigning to a pair of spaces $(X,A)$ the associated long exact sequence of homology groups. Morphisms in the domain category are maps of pairs, and in the target category morphisms are maps between exact sequences forming commutative diagrams. This functor is the composition of two functors, the first assigning to $(X,A)$ a short exact sequence of chain complexes, the second assigning to such a short exact sequence the associated long exact sequence of homology groups. Morphisms in the intermediate category are the evident commutative diagrams.
\end{itemize}

Another sort of process we have encountered is the transformation of one functor into another, for example:

\begin{itemize}
\item Boundary maps $H_n(X,A) \to H_{n-1}(A)$ in singular homology, or indeed in any homology theory.
\end{itemize}

%% --- END page-0173.tex ---
%% --- BEGIN page-0174.tex ---
\begin{itemize}
\item Change-of-coefficient homomorphisms $H_n(X;G_1) \to H_n(X;G_2)$ induced by a homomorphism $G_1 \to G_2$, as in the proof of Lemma 2.49.
\end{itemize}

In general, if one has two functors $F, G: \mathcal{C} \to \mathcal{D}$ then a \textbf{natural transformation} $T$ from $F$ to $G$ assigns a morphism $T_X: F(X) \to G(X)$ to each object $X \in \mathcal{C}$, in such a way that for each morphism $f: X \to Y$ in $\mathcal{C}$ the square at the right commutes. The case that $F$ and $G$ are contravariant rather than covariant is similar.



We have been describing the passage from topology to the abstract world of categories and functors, but there is also a nice path in the opposite direction:

\begin{itemize}
\item To each category $\mathcal{C}$ there is associated a $\Delta$-complex $B\mathcal{C}$ called the \textbf{classifying space} of $\mathcal{C}$, whose $n$-simplices are the strings $X_0 \to X_1 \to \cdots \to X_n$ of morphisms in $\mathcal{C}$. The faces of this simplex are obtained by deleting an $X_i$, and then composing the two adjacent morphisms if $i \neq 0, n$. Thus when $n = 2$ the three faces of $X_0 \to X_1 \to X_2$ are $X_0 \to X_1$, $X_1 \to X_2$, and the composed morphism $X_0 \to X_2$. In case $\mathcal{C}$ has a single object and the morphisms of $\mathcal{C}$ form a group $G$, then $B\mathcal{C}$ is the same as the $\Delta$-complex $BG$ constructed in Example 1B.7, a $K(G, 1)$. In general, the space $B\mathcal{C}$ need not be a $K(G, 1)$, however. For example, if we start with a $\Delta$-complex $X$ and regard its set of simplices as a partially ordered set $\mathcal{C}(X)$ under the relation of inclusion of faces, then $B\mathcal{C}(X)$ is the barycentric subdivision of $X$.

\item A functor $F: \mathcal{C} \to \mathcal{D}$ induces a map $B\mathcal{C} \to B\mathcal{D}$. This is the $\Delta$-map that sends an $n$-simplex $X_0 \to X_1 \to \cdots \to X_n$ to the $n$-simplex $F(X_0) \to F(X_1) \to \cdots \to F(X_n)$.

\item A natural transformation from a functor $F$ to a functor $G$ induces a homotopy between the induced maps of classifying spaces. We leave this for the reader to make explicit, using the subdivision of $\Delta^n \times I$ into $(n + 1)$-simplices described earlier in the chapter.
\end{itemize}

\section*{Exercises}

%ORIGINAL-NUMBER: 1
\begin{exercise}
If $T_n(X, A)$ denotes the torsion subgroup of $H_n(X, A; \mathbb{Z})$, show that the functors $(X, A) \mapsto T_n(X, A)$, with the obvious induced homomorphisms $T_n(X, A) \to T_n(Y, B)$ and boundary maps $T_n(X, A) \to T_{n-1}(A)$, do not define a homology theory. Do the same for the `mod torsion' functor $MT_n(X, A) = H_n(X, A; \mathbb{Z})/T_n(X, A)$.
\end{exercise}

%ORIGINAL-NUMBER: 2
\begin{exercise}
Define a candidate for a reduced homology theory on CW complexes by $\tilde{h}_n(X) = \prod_i \tilde{H}_i(X) / \bigoplus_i \tilde{H}_i(X)$. Thus $\tilde{h}_n(X)$ is independent of $n$ and is zero if $X$ is finite-dimensional, but is not identically zero, for example for $X = \bigvee_i S^1$. Show that the axioms for a homology theory are satisfied except that the wedge axiom fails.
\end{exercise}

%ORIGINAL-NUMBER: 3
\begin{exercise}
Show that if $\tilde{h}$ is a reduced homology theory, then $\tilde{h}_n(\text{point}) = 0$ for all $n$. Deduce that there are suspension isomorphisms $\tilde{h}_n(X) \cong \tilde{h}_{n+1}(SX)$ for all $n$.
\end{exercise}

%ORIGINAL-NUMBER: 4
\begin{exercise}
Show that the wedge axiom for homology theories follows from the other axioms in the case of finite wedge sums.
\end{exercise}

%% --- END page-0174.tex ---
%% --- BEGIN page-0175.tex ---
\section*{Additional Topics}

\section{2.A Homology and Fundamental Group}

There is a close connection between $H_1(X)$ and $\pi_1(X)$, arising from the fact that a map $f:I \to X$ can be viewed as either a path or a singular 1-simplex. If $f$ is a loop, with $f(0) = f(1)$, this singular 1-simplex is a cycle since $\partial f = f(1) - f(0)$.

\begin{theorem}[Hurewicz theorem in dimension 1]
\label{thm:hurewicz-dimension-1}
\dcref{ch2:2A.1}
\group{Singular Homology}
\uses{def:loop-and-pi1-set,prop:h0-path-connected}
By regarding loops as singular 1-cycles, we obtain a homomorphism $h: \pi_1(X, x_0) \to H_1(X)$. If $X$ is path-connected, then $h$ is surjective and has kernel the commutator subgroup of $\pi_1(X)$, so $h$ induces an isomorphism from the abelianization of $\pi_1(X)$ onto $H_1(X)$.
\end{theorem}
\begin{proof}
 Recall the notation $f \simeq g$ for the relation of homotopy, fixing endpoints, between paths $f$ and $g$. Regarding $f$ and $g$ as chains, the notation $f \sim g$ will mean that $f$ is homologous to $g$, that is, $f - g$ is the boundary of some 2-chain. Here are some facts about this relation.

(i) If $f$ is a constant path, then $f \sim 0$. Namely, $f$ is a cycle since it is a loop, and since $H_1(\text{point}) = 0$, $f$ must then be a boundary. Explicitly, $f$ is the boundary of the constant singular 2-simplex $\sigma$ having the same image as $f$ since

\[
\partial \sigma = \sigma|[v_1, v_2] - \sigma|[v_0, v_2] + \sigma|[v_0, v_1] = f - f + f = f
\]

(ii) If $f \simeq g$ then $f \sim g$. To see this, consider a homotopy $F: I \times I \to X$ from $f$ to $g$. This yields a pair of singular 2-simplices $\sigma_1$ and $\sigma_2$ in $X$ by subdividing the square $I \times I$ into two triangles $[v_0, v_1, v_2]$ and $[v_0, v_2, v_3]$ as shown in the figure. When one computes $\partial(\sigma_1 - \sigma_2)$, the two restrictions of $F$ to the diagonal of the square cancel, leaving $f - g$ together with two constant singular 1-simplices from the left and right edges of the square. By (i) these are boundaries, so $f - g$ is also a boundary.

(iii) $f \cdot g \sim f + g$, where $f \cdot g$ denotes the product of the paths $f$ and $g$. For if $\sigma: \Delta^2 \to X$ is the composition of orthogonal projection of $\Delta^2 = [v_0, v_1, v_2]$ onto the edge $[v_0, v_2]$ followed by $f \cdot g: [v_0, v_2] \to X$, then $\partial \sigma = g - f \cdot g + f$.

(iv) $\bar{f} \sim -f$, where $\bar{f}$ is the inverse path of $f$. This follows from the preceding three observations, which give $f + \bar{f} \sim f \cdot \bar{f} \sim 0$.

Applying (ii) and (iii) to loops, it follows that we have a well-defined homomorphism $h: \pi_1(X, x_0) \to H_1(X)$ sending the homotopy class of a loop $f$ to the homology class of the 1-cycle $f$.



%% --- END page-0175.tex ---
%% --- BEGIN page-0176.tex ---
To show $h$ is surjective when $X$ is path-connected, let $\sum_i n_i\sigma_i$ be a 1-cycle representing a given element of $H_1(X)$. After relabeling the $\sigma_i$'s we may assume each $n_i$ is $\pm 1$. By (iv) we may in fact take each $n_i$ to be $+1$, so our 1-cycle is $\sum_i \sigma_i$. If some $\sigma_i$ is not a loop, then the fact that $\partial(\sum_i \sigma_i) = 0$ means there must be another $\sigma_j$ such that the composed path $\sigma_i \circ \sigma_j$ is defined. By (iii) we may then combine the terms $\sigma_i$ and $\sigma_j$ into a single term $\sigma_i \circ \sigma_j$. Iterating this, we reduce to the case that each $\sigma_i$ is a loop. Since $X$ is path-connected, we may choose a path $y_i$ from $x_0$ to the basepoint of $\sigma_i$. We have $y_i \circ \sigma_i \circ \overline{y_i} \sim \sigma_i$ by (iii) and (iv), so we may assume all $\sigma_i$'s are loops at $x_0$. Then we can combine all the $\sigma_i$'s into a single $\sigma$ by (iii). This says the given element of $H_1(X)$ is in the image of $h$.

The commutator subgroup of $\pi_1(X)$ is contained in the kernel of $h$ since $H_1(X)$ is abelian. To obtain the reverse inclusion we will show that every class $[f]$ in the kernel of $h$ is trivial in the abelianization $\pi_1(X)_{ab}$ of $\pi_1(X)$.

If an element $[f] \in \pi_1(X)$ is in the kernel of $h$, then $f$, as a 1-cycle, is the boundary of a 2-chain $\sum_i n_i \sigma_i$. Again we may assume each $n_i$ is $\pm 1$. As in the discussion preceding Proposition 2.6, we can associate to the chain $\sum_i n_i\sigma_i$ a 2-dimensional $\Delta$-complex $K$ by taking a 2-simplex $\Delta_i^2$ for each $\sigma_i$ and identifying certain pairs of edges of these 2-simplices. Namely, if we apply the usual boundary formula to write $\partial \sigma_i = \tau_{i0} - \tau_{i1} + \tau_{i2}$ for singular 1-simplices $\tau_{ij}$, then the formula



\[f = \partial\left(\sum_i n_i\sigma_i\right) = \sum_i n_i \partial \sigma_i = \sum_{i,j}(-1)^j n_i \tau_{ij}\]

implies that we can group all but one of the $\tau_{ij}$'s into pairs for which the two coefficients $(-1)^j n_i$ in each pair are $+1$ and $-1$. The one remaining $\tau_{ij}$ is equal to $f$. We then identify edges of the $\Delta_i^2$'s corresponding to the paired $\tau_{ij}$'s, preserving orientations of these edges so that we obtain a $\Delta$-complex $K$.

The maps $\sigma_i$ fit together to give a map $\sigma : K \to X$. We can deform $\sigma$, staying fixed on the edge corresponding to $f$, so that each vertex maps to the basepoint $x_0$, in the following way. Paths from the images of these vertices to $x_0$ define such a homotopy on the union of the 0-skeleton of $K$ with the edge corresponding to $f$, and then we can appeal to the homotopy extension property in Proposition 0.16 to extend this homotopy to all of $K$. Alternatively, it is not hard to construct such an extension by hand. Restricting the new $\sigma$ to the simplices $\Delta_i^2$, we obtain a new chain $\sum_i n_i\sigma_i$ with boundary equal to $f$ and with all $\tau_{ij}$'s loops at $x_0$.

Using additive notation in the abelian group $\pi_1(X)_{ab}$, we have the formula $[f] = \sum_{i,j}(-1)^j n_i[\tau_{ij}]$ because of the cancelling pairs of $\tau_{ij}$'s. We can rewrite the summation $\sum_{i,j}(-1)^j n_i[\tau_{ij}]$ as $\sum_i n_i[\partial \sigma_i]$ where $[\partial \sigma_i] = [\tau_{i0}] - [\tau_{i1}] + [\tau_{i2}]$. Since $\sigma_i$ gives a nullhomotopy of the composed loop $\tau_{i0} - \tau_{i1} + \tau_{i2}$, we conclude that $[f] = 0$ in $\pi_1(X)_{ab}$.
\end{proof}
%% --- END page-0176.tex ---
%% --- BEGIN page-0177.tex ---
The end of this proof can be illuminated by looking more closely at the geometry. The complex $K$ is in fact a compact surface with boundary consisting of a single circle formed by the edge corresponding to $f$. This is because any pattern of identifications of pairs of edges of a finite collection of disjoint 2-simplices produces a compact surface with boundary. We leave it as an exercise for the reader to check that the algebraic formula $f = \partial(\sum_i n_i \sigma_i)$ with each $n_i = \pm 1$ implies that $K$ is an orientable surface. The component of $K$ containing the boundary circle is a standard closed orientable surface of some genus $g$ with an open disk removed, by the basic structure theorem for compact orientable surfaces. Giving this surface the cell structure indicated in the figure, it then becomes obvious that $f$ is homotopic to a product of $g$ commutators in $\pi_1(X)$.



The map $h:\pi_1(X,x_0) \to H_1(X)$ can also be defined by $h([f]) = f_* (\alpha)$ where $f:S^1 \to X$ represents a given element of $\pi_1(X,x_0)$, $f_*$ is the induced map on $H_1$, and $\alpha$ is the generator of $H_1(S^1) = \mathbb{Z}$ represented by the standard map $\sigma:I \to S^1$, $\sigma(s) = e^{2\pi i s}$. This is because both $[f] \in \pi_1(X,x_0)$ and $f_*(\alpha) \in H_1(X)$ are represented by the loop $f \circ \sigma: I \to X$. A consequence of this definition is that $h([f]) = h([g])$ if $f$ and $g$ are homotopic maps $S^1 \to X$, since $f_* = g_*$ by Theorem 2.10.

\begin{example}[H1 of closed orientable surface]
\label{ex:h1-closed-orientable-surface}
\dcref{ch2:2A.2}
\group{Singular Homology}
\uses{thm:hurewicz-dimension-1}
For the closed orientable surface $M$ of genus $g$, the abelianization of $\pi_1(M)$ is $\mathbb{Z}^{2g}$, the product of $2g$ copies of $\mathbb{Z}$, and a basis for $H_1(M)$ consists of the 1-cycles represented by the 1-cells of $M$ in its standard CW structure. We can also represent a basis by the loops $\alpha_i$ and $\beta_i$ shown in the figure below since these
\end{example}



loops are homotopic to the loops represented by the 1-cells, as one can see in the picture of the cell structure in Chapter 0. The loops $y_i$, on the other hand, are trivial in homology since the portion of $M$ on one side of $y_i$ is a compact surface bounded by $y_i$, so $y_i$ is homotopic to a loop that is a product of commutators, as we saw a couple paragraphs earlier. The loop $\alpha_i'$ represents the same homology class as $\alpha_i$ since the region between $y_i$ and $\alpha_i \cup \alpha_i'$ provides a homotopy between $y_i$ and a product of two loops homotopic to $\alpha_i$ and the inverse of $\alpha_i'$, so $\alpha_i - \alpha_i' \sim y_i \sim 0$, hence $\alpha_i \sim \alpha_i'$.



%% --- END page-0177.tex ---
%% --- BEGIN page-0178.tex ---
\section{Classical Applications}

In this section we use homology theory to prove several interesting results in topology and algebra whose statements give no hint that algebraic topology might be involved.

To begin, we calculate the homology of complements of embedded spheres and disks in a sphere. Recall that an embedding is a map that is a homeomorphism onto its image.

\begin{proposition}[Homology of sphere complement of disk/sphere]
\label{prop:homology-sphere-complement-disk-sphere}
\dcref{ch2:2B.1}
\group{Singular Homology}
\uses{thm:excision-general,ex:mayer-vietoris-sphere}
(a) For an embedding $h:D^k \to S^n$, $\tilde{H}_i(S^n - h(D^k)) = 0$ for all $i$.
(b) For an embedding $h:S^k \to S^n$ with $k < n$, $\tilde{H}_i(S^n - h(S^k))$ is $\mathbb{Z}$ for $i = n - k - 1$ and $0$ otherwise.
\end{proposition}

As a special case of (b) we have the Jordan curve theorem: A subspace of $S^2$ homeomorphic to $S^1$ separates $S^2$ into two complementary components, or equivalently, path-components since open subsets of $S^n$ are locally path-connected. One could just as well use $\mathbb{R}^2$ in place of $S^2$ here since deleting a point from an open set in $S^2$ does not affect its connectedness. More generally, (b) says that a subspace of $S^n$ homeomorphic to $S^{n-1}$ separates it into two components, and these components have the same homology groups as a point. Somewhat surprisingly, there are embeddings where these complementary components are not simply-connected as they are for the standard embedding. An example is the Alexander horned sphere in $S^3$ which we describe in detail following the proof of the proposition. These complications involving embedded $S^{n-1}$'s in $S^n$ are all local in nature since it is known that any locally nicely embedded $S^{n-1}$ in $S^n$ is equivalent to the standard $S^{n-1} \subset S^n$, equivalent in the sense that there is a homeomorphism of $S^n$ taking the given embedded $S^{n-1}$ onto the standard $S^{n-1}$. In particular, both complementary regions are homeomorphic to open balls. See [Brown 1960] for a precise statement and proof. When $n = 2$ it is a classical theorem of Schoenflies that all embeddings $S^1 \hookrightarrow S^2$ are equivalent.

By contrast, when we come to embeddings of $S^{n-2}$ in $S^n$, even locally nice embeddings need not be equivalent to the standard one. This is the subject of knot theory, including the classical case of knotted embeddings of $S^1$ in $S^3$ or $\mathbb{R}^3$. For embeddings of $S^{n-2}$ in $S^n$ the complement always has the same homology as $S^1$, according to the theorem, but the fundamental group can be quite different. In spite of the fact that the homology of a knot complement does not detect knottedness, it is still possible to use homology to distinguish different knots by looking at the homology of covering spaces of their complements.
\begin{proof}
 We prove (a) by induction on $k$. When $k = 0$, $S^n - h(D^0)$ is homeomorphic to $\mathbb{R}^n$, so this case is trivial. For the induction step it will be convenient to replace the domain disk $D^k$ of $h$ by the cube $I^k$. Let $A = S^n - h(I^{k-1} \times [0, 1/2])$ and let

%% --- END page-0178.tex ---
%% --- BEGIN page-0179.tex ---
$B = S^n - h(I^{k-1}\times[1/2, 1])$, so $A \cap B = S^n - h(I^k)$ and $A \cup B = S^n - h(I^{k-1} \times \{1/2\})$. By
induction $\tilde{H}_i(A \cup B) = 0$ for all $i$, so the Mayer-Vietoris sequence gives isomorphisms
$\Phi: \tilde{H}_i(S^n - h(I^k)) \to \tilde{H}_i(A) \oplus \tilde{H}_i(B)$ for all $i$. Modulo signs, the two components of $\Phi$
are induced by the inclusions $S^n - h(I^k) \hookrightarrow A$ and $S^n - h(I^k) \hookrightarrow B$, so if there exists
an $i$-dimensional cycle $\alpha$ in $S^n - h(I^k)$ that is not a boundary in $S^n - h(I^k)$, then
$\alpha$ is also not a boundary in at least one of $A$ and $B$. (When $i = 0$ the word `cycle'
here is to be interpreted in the sense of augmented chain complexes since we are
dealing with reduced homology.) By iteration we can then produce a nested sequence
of closed intervals $I_1 \supset I_2 \supset \cdots$ in the last coordinate of $I^k$ shrinking down to a
point $p \in I$, such that $\alpha$ is not a boundary in $S^n - h(I^{k-1}\times I_m)$ for any $m$. On
the other hand, by induction on $k$ we know that $\alpha$ is the boundary of a chain $\beta$ in
$S^n - h(I^{k-1}\times\{p\})$. This $\beta$ is a finite linear combination of singular simplices with
compact image in $S^n - h(I^{k-1}\times\{p\})$. The union of these images is covered by the
nested sequence of open sets $S^n - h(I^{k-1}\times I_m)$, so by compactness $\beta$ must actually
be a chain in $S^n - h(I^{k-1}\times I_m)$ for some $m$. This contradiction shows that $\alpha$ must
be a boundary in $S^n - h(I^k)$, finishing the induction step.

Part (b) is also proved by induction on $k$, starting with the trivial case $k = 0$ when
$S^n - h(S^0)$ is homeomorphic to $S^{n-1}\times\mathbb{R}$. For the induction step, write $S^k$ as the
union of hemispheres $D_+^k$ and $D_-^k$ intersecting in $S^{k-1}$. The Mayer-Vietoris sequence
for $A = S^n - h(D_+^k)$ and $B = S^n - h(D_-^k)$, both of which have trivial reduced homology
by part (a), then gives isomorphisms $\tilde{H}_i(S^n - h(S^k)) \approx \tilde{H}_{i+1}(S^n - h(S^{k-1}))$.
\end{proof}
If we apply the last part of this proof to an embedding $h:S^n \to S^n$, the Mayer-Vietoris sequence ends with the terms $\tilde{H}_0(A) \oplus \tilde{H}_0(B) \to \tilde{H}_0(S^n - h(S^{n-1})) \to 0$. Both
$\tilde{H}_0(A)$ and $\tilde{H}_0(B)$ are zero, so exactness would imply that $\tilde{H}_0(S^n - h(S^{n-1})) = 0$
which appears to contradict the fact that $S^n - h(S^{n-1})$ has two path-components.
The only way out of this dilemma is for $h$ to be surjective, so that $A \cap B$ is empty and
the $0$ at the end of the Mayer-Vietoris sequence is $\tilde{H}_{-1}(\emptyset)$ which is $\mathbb{Z}$ rather than $0$.

In particular, this shows that $S^n$ cannot be embedded in $\mathbb{R}^n$ since this would
yield a nonsurjective embedding in $S^n$. A consequence is that there is no embedding
$\mathbb{R}^m \to \mathbb{R}^n$ for $m > n$ since this would restrict to an embedding of $S^n \subset \mathbb{R}^m$ into $\mathbb{R}^n$.
More generally there is no continuous injection $\mathbb{R}^m \to \mathbb{R}^n$ for $m > n$ since this too
would give an embedding $S^n \to \mathbb{R}^n$.

\begin{example}[The Alexander Horned Sphere]
\label{ex:alexander-horned-sphere}
\dcref{ch2:2B.2}
\group{Singular Homology}
\uses{prop:homology-sphere-complement-disk-sphere}
This is a subspace $S \subset \mathbb{R}^3$ homeomorphic to $S^2$ such that the unbounded component of $\mathbb{R}^3 - S$ is not simply-connected
as it is for the standard $S^2 \subset \mathbb{R}^3$. We will construct $S$ by defining a sequence of compact subspaces $X_0 \supset X_1 \supset \cdots$ of $\mathbb{R}^3$ whose intersection is homeomorphic to a ball,
and then $S$ will be the boundary sphere of this ball.
\end{example}

We begin with $X_0$ a solid torus $S^1 \times D^2$ obtained from a ball $B_0$ by attaching
a handle $I \times D^2$ along $\partial I \times D^2$. In the figure this handle is shown as the union of

%% --- END page-0179.tex ---
%% --- BEGIN page-0180.tex ---
two `horns' attached to the ball, together with a shorter handle drawn as dashed lines. To form the space $X_1 \subset X_0$ we delete part of the short handle, so that what remains is a pair of linked handles attached to the ball $B_1$ that is the union of $B_0$ with the two horns. To form $X_2$ the process is repeated: Decompose each of the second stage handles as a pair of horns and a short handle, then delete a part of the short handle. In the same way $X_n$ is constructed inductively from $X_{n-1}$. Thus $X_n$ is a ball $B_n$ with $2^n$ handles attached, and $B_n$ is obtained from $B_{n-1}$ by attaching $2^n$ horns. There are homeomorphisms $h_n : B_{n-1} \to B_n$ that are the identity outside a small neighborhood of $B_n - B_{n-1}$. As $n$ goes to infinity, the composition $h_n \cdots h_1$ approaches a map $f : B_0 \to \mathbb{R}^3$ which is continuous since the convergence is uniform. The set of points in $B_0$ where $f$ is not equal to $h_n \cdots h_1$ for large $n$ is a Cantor set, whose image under $f$ is the intersection of all the handles. It is not hard to see that $f$ is one-to-one. By compactness it follows that $f$ is a homeomorphism onto its image, a ball $B \subset \mathbb{R}^3$ whose boundary sphere $f(\partial B_0)$ is $S$, the Alexander horned sphere.



Now we compute $\pi_1(\mathbb{R}^3 - B)$. Note that $B$ is the intersection of the $X_n$'s, so $\mathbb{R}^3 - B$ is the union of the complements $Y_n$ of the $X_n$'s, which form an increasing sequence $Y_0 \subset Y_1 \subset \cdots$. We will show that the groups $\pi_1(Y_n)$ also form an increasing sequence of successively larger groups, whose union is $\pi_1(\mathbb{R}^3 - B)$. To begin we have $\pi_1(Y_0) \cong \mathbb{Z}$ since $X_0$ is a solid torus embedded in $\mathbb{R}^3$ in a standard way. To compute $\pi_1(Y_1)$, let $\overline{Y_0}$ be the closure of $Y_0$ in $Y_1$, so $\overline{Y_0} - Y_0$ is an open annulus $A$ and $\pi_1(\overline{Y_0})$ is also $\mathbb{Z}$. We obtain $Y_1$ from $\overline{Y_0}$ by attaching the space $Z = Y_1 - Y_0$ along $A$. The group $\pi_1(Z)$ is the free group $F_2$ on two generators $\alpha_1$ and $\alpha_2$ represented by loops linking the two handles, since $Z - A$ is homeomorphic to an open ball with two straight tubes deleted. A loop $\alpha$ generating $\pi_1(A)$ represents the commutator $[\alpha_1, \alpha_2]$, as one can see by noting that the closure of $Z$ is obtained from $Z$ by adjoining two disjoint surfaces, each homeomorphic to a torus with an open disk removed; the boundary of this disk is homotopic to $\alpha$ and is also homotopic to the commutator of meridian and longitude circles in the torus, which correspond to $\alpha_1$ and $\alpha_2$. Van Kampen's theorem now implies that the inclusion $Y_0 \hookrightarrow Y_1$ induces an injection of $\pi_1(Y_0)$ into $\pi_1(Y_1)$ as the infinite cyclic subgroup generated by $[\alpha_1, \alpha_2]$.

In a similar way we can regard $Y_{n+1}$ as being obtained from $Y_n$ by adjoining $2^n$ copies of $Z$. Assuming inductively that $\pi_1(Y_n)$ is the free group $F_{2^n}$ with generators represented by loops linking the $2^n$ smallest handles of $X_n$, then each copy of $Z$ ad-

%% --- END page-0180.tex ---
%% --- BEGIN page-0181.tex ---
joined to $Y_n$ changes $\pi_1(Y_n)$ by making one of the generators into the commutator of
two new generators. Note that adjoining a copy of $Z$ induces an injection on $\pi_1$ since
the induced homomorphism is the free product of the injection $\pi_1(A) \to \pi_1(Z)$ with
the identity map on the complementary free factor. Thus the map $\pi_1(Y_n) \to \pi_1(Y_{n+1})$
is an injection $F_{2n} \to F_{2n+1}$. The group $\pi_1(\mathbb{R}^3 - B)$ is isomorphic to the union of this
increasing sequence of groups by a compactness argument: Each loop in $\mathbb{R}^3 - B$ has
compact image and hence must lie in some $Y_n$, and similarly for homotopies of loops.

In particular we see explicitly why $\pi_1(\mathbb{R}^3 - B)$ has trivial abelianization, because
each of its generators is exactly equal to the commutator of two other generators.
This inductive construction in which each generator of a free group is decreed to
be the commutator of two new generators is perhaps the simplest way of building a
nontrivial group with trivial abelianization, and for the construction to have such a
nice geometric interpretation is something to marvel at. From a naive viewpoint it may
seem a little odd that a highly nonfree group can be built as a union of an increasing
sequence of free groups, but this can also easily happen for abelian groups, as $\mathbb{Q}$ for
example is the union of an increasing sequence of infinite cyclic subgroups.

The next theorem says that for subspaces of $\mathbb{R}^n$, the property of being open is
a topological invariant. This result is known classically as Invariance of Domain, the
word `domain' being an older designation for an open set in $\mathbb{R}^n$.

\begin{theorem}[Invariance of domain]
\label{thm:invariance-of-domain-general}
\dcref{ch2:2B.3}
\group{Singular Homology}
\uses{prop:homology-sphere-complement-disk-sphere}
If $U$ is an open set in $\mathbb{R}^n$ and $h:U \to \mathbb{R}^n$ is an embedding, or more
generally just a continuous injection, then the image $h(U)$ is an open set in $\mathbb{R}^n$.
\end{theorem}
\begin{proof}
 Viewing $S^n$ as the one-point compactification of $\mathbb{R}^n$, an equivalent statement
is that $h(U)$ is open in $S^n$, and this is what we will prove. Each $x \in U$ is the center
point of a disk $D^n \subset U$. It will suffice to prove that $h(D^n - \partial D^n)$ is open in $S^n$. The
hypothesis on $h$ implies that its restrictions to $D^n$ and $\partial D^n$ are embeddings. By the
previous proposition $S^n - h(\partial D^n)$ has two path-components. These path-components
are $h(D^n - \partial D^n)$ and $S^n - h(D^n)$ since these two subspaces are disjoint and path-
connected, the first since it is homeomorphic to $D^n - \partial D^n$ and the second by the
proposition. Since $S^n - h(\partial D^n)$ is open in $S^n$, its path-components are the same as
its components. The components of a space with finitely many components are open,
so $h(D^n - \partial D^n)$ is open in $S^n - h(\partial D^n)$ and hence also in $S^n$.
\end{proof}
Here is an application involving the notion of an $n$-manifold, which is a Hausdorff
space locally homeomorphic to $\mathbb{R}^n$:

\begin{corollary}[Embedding of compact manifold in connected manifold]
\label{cor:embedding-compact-manifold-connected}
\dcref{ch2:2B.4}
\group{Singular Homology}
\uses{thm:invariance-of-domain-general}
If $M$ is a compact $n$-manifold and $N$ is a connected $n$-manifold,
then an embedding $h:M \to N$ must be surjective, hence a homeomorphism.
\end{corollary}
\begin{proof}
 $h(M)$ is closed in $N$ since it is compact and $N$ is Hausdorff. Since $N$ is
connected it suffices to show $h(M)$ is also open in $N$, and this is immediate from the
theorem.
\end{proof}
%% --- END page-0181.tex ---
%% --- BEGIN page-0182.tex ---
The Invariance of Domain and the $n$-dimensional generalization of the Jordan curve theorem were first proved by Brouwer around 1910, at a very early stage in the development of algebraic topology.

\textbf{Division Algebras}

Here is an algebraic application of homology theory due to H. Hopf:

\begin{theorem}[Real and complex division algebras]
\label{thm:real-complex-division-algebras}
\dcref{ch2:2B.5}
\group{Cellular Homology}
\uses{prop:z2-acts-freely-even-spheres}
$\mathbb{R}$ and $\mathbb{C}$ are the only finite-dimensional division algebras over $\mathbb{R}$ which are commutative and have an identity.
\end{theorem}

By definition, an \textbf{algebra} structure on $\mathbb{R}^n$ is simply a bilinear multiplication map $\mathbb{R}^n \times \mathbb{R}^n \to \mathbb{R}^n$, $(a,b) \mapsto ab$. Thus the product satisfies left and right distributivity, $a(b+c) = ab + ac$ and $(a+b)c = ac + bc$, and scalar associativity, $\alpha(ab) = (\alpha a)b = a(\alpha b)$ for $\alpha \in \mathbb{R}$. Commutativity, full associativity, and an identity element are not assumed. An algebra is a \textbf{division algebra} if the equations $ax = b$ and $xa = b$ are always solvable whenever $a \neq 0$. In other words, the linear transformations $x \mapsto ax$ and $x\mapsto xa$ are surjective when $a \neq 0$. These are linear maps $\mathbb{R}^n \to \mathbb{R}^n$, so surjectivity is equivalent to having trivial kernel, which means there are no zero-divisors.

The four classical examples are $\mathbb{R}$, $\mathbb{C}$, the quaternion algebra $\mathbb{H}$, and the octonion algebra $\mathbb{O}$. Frobenius proved in 1877 that $\mathbb{R}$, $\mathbb{C}$, and $\mathbb{H}$ are the only finite-dimensional associative division algebras over $\mathbb{R}$ with an identity element. If the product satisfies $|ab| = |a||b|$ as in the classical examples, then Hurwitz showed in 1898 that the dimension of the algebra must be 1, 2, 4, or 8, and others subsequently showed that the only examples with an identity element are the classical ones. A full discussion of all this, including some examples showing the necessity of the hypothesis of an identity element, can be found in [Ebbinghaus 1991]. As one would expect, the proofs of these results are algebraic, but if one drops the condition that $|ab| = |a||b|$ it seems that more topological proofs are required. We will show in Theorem 3.21 that a finite-dimensional division algebra over $\mathbb{R}$ must have dimension a power of 2. The fact that the dimension can be at most 8 is a famous theorem of [Bott \& Milnor 1958] and [Kervaire 1958]. See \S 4.B for a few more comments on this.
\begin{proof}
 Suppose first that $\mathbb{R}^n$ has a commutative division algebra structure. Define a map $f: S^{n-1} \to S^{n-1}$ by $f(x) = x^2/|x|^2$. This is well-defined since $x \neq 0$ implies $x^2 \neq 0$ in a division algebra. The map $f$ is continuous since the multiplication map $\mathbb{R}^n \times \mathbb{R}^n \to \mathbb{R}^n$ is bilinear, hence continuous. Since $f(-x) = f(x)$ for all $x$, $f$ induces a quotient map $\bar{f}: \mathbb{RP}^{n-1} \to S^{n-1}$. The following argument shows that $\bar{f}$ is injective. An equality $f(x) = f(y)$ implies $x^2 = \alpha^2 y^2$ for $\alpha = (|x|^2/|y|^2)^{1/2} > 0$. Thus we have $x^2 - \alpha^2 y^2 = 0$, which factors as $(x + \alpha y)(x - \alpha y) = 0$ using commutativity and the fact that $\alpha$ is a real scalar. Since there are no divisors of zero, we deduce that $x = \pm \alpha y$. Since $x$ and $y$ are unit vectors and $\alpha$ is real, this yields $x = \pm y$, so $x$ and $y$ determine the same point of $\mathbb{RP}^{n-1}$, which means that $\bar{f}$ is injective.

%% --- END page-0182.tex ---
%% --- BEGIN page-0183.tex ---
Since $\bar{f}$ is an injective map of compact Hausdorff spaces, it must be a homeomorphism onto its image. By Corollary 2B.4, $\bar{f}$ must in fact be surjective if we are not in the trivial case $n = 1$. Thus we have a homeomorphism $\mathbb{RP}^{n-1} \simeq S^{n-1}$. This implies $n = 2$ since if $n > 2$ the spaces $\mathbb{RP}^{n-1}$ and $S^{n-1}$ have different homology groups (or different fundamental groups).

It remains to show that a 2-dimensional commutative division algebra $A$ with identity is isomorphic to $\mathbb{C}$. This is elementary algebra: If $j \in A$ is not a real scalar multiple of the identity element $1 \in A$ and we write $j^2 = a + bj$ for $a, b \in \mathbb{R}$, then $(j - b/2)^2 = a + b^2/4$ so by rechoosing $j$ we may assume that $j^2 = a \in \mathbb{R}$. If $a = 0$, say $a = c^2$, then $j^2 = c^2$ implies $(j-c)(j-c) = 0$, so $j = \pm c$, but this contradicts the choice of $j$. So $j^2 = -c^2$ and by rescaling $j$ we may assume $j^2 = -1$, hence $A$ is isomorphic to $\mathbb{C}$.
\end{proof}
Leaving out the last paragraph, the proof shows that a finite-dimensional commutative division algebra, not necessarily with an identity, must have dimension at most 2. Oddly enough, there do exist 2-dimensional commutative division algebras without identity elements, for example $\mathbb{C}$ with the modified multiplication $z \cdot w = z\bar{w}$, the bar denoting complex conjugation.

\textbf{The Borsuk-Ulam Theorem}

In Theorem 1.10 we proved the 2-dimensional case of the Borsuk-Ulam theorem, and now we will give a proof for all dimensions, using the following theorem of Borsuk:

\begin{proposition}[Odd map of Sn has odd degree]
\label{prop:odd-map-odd-degree}
\dcref{ch2:2B.6}
\group{Cellular Homology}
\uses{cor:homology-spheres}
An odd map $f : S^n \to S^n$, satisfying $f(-x) = -f(x)$ for all $x$, must have odd degree.
\end{proposition}

The corresponding result that even maps have even degree is easier, and was an exercise for $\S 2.2$.

The proof will show that using homology with a coefficient group other than $\mathbb{Z}$ can sometimes be a distinct advantage. The main ingredient will be a certain exact sequence associated to a two-sheeted covering space $p : \tilde{X} \to X$,

\[ \cdots \to H_n(X; \mathbb{Z}_2) \xrightarrow{\tau} H_n(\tilde{X}; \mathbb{Z}_2) \xrightarrow{p_*} H_n(X; \mathbb{Z}_2) \to H_{n-1}(X; \mathbb{Z}_2) \to \cdots \]

This is the long exact sequence of homology groups associated to a short exact sequence of chain complexes consisting of short exact sequences of chain groups

\[ 0 \to C_n(X; \mathbb{Z}_2) \xrightarrow{\tau} C_n(\tilde{X}; \mathbb{Z}_2) \xrightarrow{p_\#} C_n(X; \mathbb{Z}_2) \to 0 \]

The map $p_\#$ is surjective since singular simplices $\sigma : \Delta^n \to X$ always lift to $\tilde{X}$, as $\Delta^n$ is simply-connected. Each $\sigma$ has in fact precisely two lifts $\tilde{\sigma}_1$ and $\tilde{\sigma}_2$. Because we are using $\mathbb{Z}_2$ coefficients, the kernel of $p_\#$ is generated by the sums $\tilde{\sigma}_1 + \tilde{\sigma}_2$. So if we define $\tau$ to send each $\sigma : \Delta^n \to X$ to the sum of its two lifts to $\tilde{\Delta}^n$, then the image of $\tau$ is the kernel of $p_\#$. Obviously $\tau$ is injective, so we have the short exact sequence

%% --- END page-0183.tex ---
%% --- BEGIN page-0184.tex ---
indicated. Since $\tau$ and $p_*$ commute with boundary maps, we have a short exact sequence of chain complexes, yielding the long exact sequence of homology groups.

The map $\tau_*$ is a special case of more general \textit{transfer homomorphisms} considered in \S3.G, so we will refer to the long exact sequence involving the maps $\tau_*$ as the \textit{transfer sequence}. This sequence can also be viewed as a special case of the Gysin sequences discussed in \S4.D. There is a generalization of the transfer sequence to homology with other coefficients, but this uses a more elaborate form of homology called homology with local coefficients, as we show in \S3.H.
\begin{proof}[2B.6]
 The proof will involve the transfer sequence for the covering space $p : S^n \to \mathbb{RP}^n$. This has the following form, where to simplify notation we abbreviate $\mathbb{RP}^n$ to $P^n$ and we let the coefficient group $\mathbb{Z}_2$ be implicit:

\[
0 \longrightarrow H_n(P^n) \xrightarrow{\tau_*} H_n(S^n) \xrightarrow{p_*} H_n(P^n) \longrightarrow H_{n-1}(P^n) \longrightarrow 0 \longrightarrow \cdots
\]
\[
\cdots \longrightarrow 0 \longrightarrow H_i(P^n) \xrightarrow{\cong} H_{i-1}(P^n) \longrightarrow 0 \longrightarrow \cdots
\]
\[
\cdots \longrightarrow 0 \longrightarrow H_1(P^n) \xrightarrow{\cong} H_0(P^n) \xrightarrow{p_*} H_0(S^n) \xrightarrow{p_*} H_0(P^n) \longrightarrow 0
\]

The initial $0$ is $H_{n+1}(P^n; \mathbb{Z}_2)$, which vanishes since $P^n$ is an $n$-dimensional CW complex. The other terms that are zero are $H_i(S^n)$ for $0 < i < n$. We assume $n > 1$, leaving the minor modifications needed for the case $n = 1$ to the reader. All the terms that are not zero are $\mathbb{Z}_2$, by cellular homology. Alternatively, this exact sequence can be used to compute the homology groups $H_i(\mathbb{RP}^n; \mathbb{Z}_2)$ if one does not already know them. Since all the nonzero groups in the sequence are $\mathbb{Z}_2$, exactness forces the maps to be isomorphisms or zero as indicated.

An odd map $f : S^n \to S^n$ induces a quotient map $\overline{f} : \mathbb{RP}^n \to \mathbb{RP}^n$. These two maps induce a map from the transfer sequence to itself, and we will need to know that the squares in the resulting diagram commute. This follows from the naturality of the long exact sequence of homology associated to a short exact sequence of chain complexes, once we verify commutativity of the diagram

\[
\begin{array}{ccccc}
0 \longrightarrow C_i(P^n) \xrightarrow{\tau} C_i(S^n) \xrightarrow{p_*} C_i(P^n) \longrightarrow 0 \\
\downarrow \overline{f}_* & & \downarrow f_* & & \downarrow \overline{f}_* \\
0 \longrightarrow C_i(P^n) \xrightarrow{\tau} C_i(S^n) \xrightarrow{p_*} C_i(P^n) \longrightarrow 0
\end{array}
\]

Here the right-hand square commutes since $p f = \overline{f} p$. The left-hand square commutes since for a singular $i$-simplex $\sigma : \Delta^i \to P^n$ with lifts $\overline{\sigma}_1$ and $\overline{\sigma}_2$, the two lifts of $\overline{f} \sigma$ are $f \overline{\sigma}_1$ and $f \overline{\sigma}_2$ since $f$ takes antipodal points to antipodal points.

Now we can see that all the maps $f_*$ and $\overline{f}_*$ in the commutative diagram of transfer sequences are isomorphisms by induction on dimension, using the evident fact that if three maps in a commutative square are isomorphisms, so is the fourth. The induction starts with the trivial fact that $f_*$ and $\overline{f}_*$ are isomorphisms in dimension zero.

%% --- END page-0184.tex ---
%% --- BEGIN page-0185.tex ---
In particular we deduce that the map $f_*: H_n(S^n; \mathbb{Z}_2) \to H_n(S^n; \mathbb{Z}_2)$ is an isomorphism. By Lemma 2.49 this map is multiplication by the degree of $f$ mod 2, so the degree of $f$ must be odd.
\end{proof}
The fact that odd maps have odd degree easily implies the Borsuk--Ulam theorem:

\begin{corollary}[Borsuk-Ulam theorem]
\label{cor:borsuk-ulam}
\dcref{ch2:2B.7}
\group{Cellular Homology}
\uses{prop:odd-map-odd-degree}
For every map $g : S^n \to \mathbb{R}^n$ there exists a point $x \in S^n$ with $g(x) = g(-x)$.
\end{corollary}
\begin{proof}
 Let $f(x) = g(x) - g(-x)$, so $f$ is odd. We need to show that $f(x) = 0$ for some $x$. If this is not the case, we can replace $f(x)$ by $f(x)/|f(x)|$ to get a new map $f : S^n \to S^{n-1}$ which is still odd. The restriction of this $f$ to the equator $S^{n-1}$ then has odd degree by the proposition. But this restriction is nullhomotopic via the restriction of $f$ to one of the hemispheres bounded by $S^{n-1}$.
\end{proof}
\section*{Exercises}

%ORIGINAL-NUMBER: 1
\begin{exercise}
Compute $H_i(S^n - X)$ when $X$ is a subspace of $S^n$ homeomorphic to $S^k \vee S^\ell$ or to $S^k \sqcup S^\ell$.
\end{exercise}

%ORIGINAL-NUMBER: 2
\begin{exercise}
Show that $\tilde{H}_i(S^n - X) \cong \tilde{H}_{i-1}(X)$ when $X$ is homeomorphic to a finite connected graph. [First do the case that the graph is a tree.]
\end{exercise}

%ORIGINAL-NUMBER: 3
\begin{exercise}
Let $(D, S) \subset (D^n, S^{n-1})$ be a pair of subspaces homeomorphic to $(D^k, S^{k-1})$, with $D \cap S^{n-1} = S$. Show the inclusion $S^{n-1} - S \hookrightarrow D^n - D$ induces an isomorphism on homology. [Glue two copies of $(D^n, D)$ to the two ends of $(S^{n-1} \times I, S \times I)$ to produce a $k$-sphere in $S^n$ and look at a Mayer--Vietoris sequence for the complement of this $k$-sphere.]
\end{exercise}

%ORIGINAL-NUMBER: 4
\begin{exercise}
In the unit sphere $S^{p+q-1} \subset \mathbb{R}^{p+q}$ let $S^{p-1}$ and $S^{q-1}$ be the subspheres consisting of points whose last $q$ and first $p$ coordinates are zero, respectively.

(a) Show that $S^{p+q-1} - S^{p-1}$ deformation retracts onto $S^{q-1}$, and is in fact homeomorphic to $S^{q-1} \times \mathbb{R}^p$.

(b) Show that $S^{p-1}$ and $S^{q-1}$ are not the boundaries of any pair of disjointly embedded disks $D^p$ and $D^q$ in $D^{p+q}$. [The preceding exercise may be useful.]
\end{exercise}

%ORIGINAL-NUMBER: 5
\begin{exercise}
Let $S$ be an embedded $k$-sphere in $S^n$ for which there exists a disk $D^n \subset S^n$ intersecting $S$ in the disk $D^k \subset D^n$ defined by the first $k$ coordinates of $D^n$. Let $D^{n-k} \subset D^n$ be the disk defined by the last $n-k$ coordinates, with boundary sphere $S^{n-k-1}$. Show that the inclusion $S^{n-k-1} \hookrightarrow S^n - S$ induces an isomorphism on homology groups.
\end{exercise}

%ORIGINAL-NUMBER: 6
\begin{exercise}
Modify the construction of the Alexander horned sphere to produce an embedding $S^2 \hookrightarrow \mathbb{R}^3$ for which neither component of $\mathbb{R}^3 - S^2$ is simply-connected.
\end{exercise}

%% --- END page-0185.tex ---
%% --- BEGIN page-0186.tex ---
\begin{exercise}
%ORIGINAL-NUMBER: 7
Analyze what happens when the number of handles in the basic building block for the Alexander horned sphere is doubled, as in the figure at the right.
\end{exercise}

\begin{exercise}
%ORIGINAL-NUMBER: 8
Show that $\mathbb{R}^{2n+1}$ is not a division algebra over $\mathbb{R}$ if $n > 0$ by considering how the determinant of the linear map $x \mapsto ax$ given by the multiplication in a division algebra structure would vary as $a$ moves along a path in $\mathbb{R}^{2n+1} - \{0\}$ joining two antipodal points.
\end{exercise}

\begin{exercise}
%ORIGINAL-NUMBER: 9
Make the transfer sequence explicit in the case of a trivial covering $\tilde{X} \to X$, where $\tilde{X} = X \times S^0$.
\end{exercise}

\begin{exercise}
%ORIGINAL-NUMBER: 10
Use the transfer sequence for the covering $S^\infty \to \mathbb{RP}^\infty$ to compute $H_n(\mathbb{RP}^\infty; \mathbb{Z}_2)$.
\end{exercise}

\begin{exercise}
%ORIGINAL-NUMBER: 11
Use the transfer sequence for the covering $X \times S^\infty \to X \times \mathbb{RP}^\infty$ to produce isomorphisms $H_n(X \times \mathbb{RP}^\infty; \mathbb{Z}_2) \cong \bigoplus_{i \leq n} H_i(X; \mathbb{Z}_2)$ for all $n$.
\end{exercise}

\section{Simplicial Approximation}

Many spaces of interest in algebraic topology can be given the structure of simplicial complexes, and early in the history of the subject this structure was exploited as one of the main technical tools. Later, CW complexes largely superseded simplicial complexes in this role, but there are still some occasions when the extra structure of simplicial complexes can be quite useful. This will be illustrated nicely by the proof of the classical Lefschetz fixed point theorem in this section.

One of the good features of simplicial complexes is that arbitrary continuous maps between them can always be deformed to maps that are linear on the simplices of some subdivision of the domain complex. This is the idea of `simplicial approximation', developed by Brouwer and Alexander before 1920. Here is the relevant definition: If $K$ and $L$ are simplicial complexes, then a map $f : K \to L$ is \textbf{simplicial} if it sends each simplex of $K$ to a simplex of $L$ by a linear map taking vertices to vertices. In barycentric coordinates, a linear map of a simplex $[v_0, \cdots, v_n]$ has the form $\sum_i t_i v_i \mapsto \sum_i t_i f(v_i)$. Since a linear map from a simplex to a simplex is uniquely determined by its values on vertices, this means that a simplicial map is uniquely determined by its values on vertices. It is easy to see that a map from the vertices of $K$ to the vertices of $L$ extends to a simplicial map iff it sends the vertices of each simplex of $K$ to the vertices of some simplex of $L$.

Here is the most basic form of the Simplicial Approximation Theorem:

\begin{theorem}[Simplicial Approximation Theorem]
\label{thm:simplicial-approximation}
\dcref{ch2:2C.1}
\group{Singular Homology}
\uses{lem:star-intersections}
If $K$ is a finite simplicial complex and $L$ is an arbitrary simplicial complex, then any map $f : K \to L$ is homotopic to a map that is simplicial with respect to some iterated barycentric subdivision of $K$.
\end{theorem}

%% --- END page-0186.tex ---
%% --- BEGIN page-0187.tex ---
To see that subdivision of $K$ is essential, consider the case of maps $S^n \to S^n$. With fixed simplicial structures on the domain and range spheres there are only finitely many simplicial maps since there are only finitely many ways to map vertices to vertices. Hence only finitely many degrees are realized by maps that are simplicial with respect to fixed simplicial structures in both the domain and range spheres. This remains true even if the simplicial structure on the range sphere is allowed to vary, since if the range sphere has more vertices than the domain sphere then the map cannot be surjective, hence must have degree zero.

Before proving the simplicial approximation theorem we need some terminology and a lemma. The \textit{star} $\text{St}\,\sigma$ of a simplex $\sigma$ in a simplicial complex $X$ is defined to be the subcomplex consisting of all the simplices of $X$ that contain $\sigma$. Closely related to this is the \textit{open star} st $\sigma$, which is the union of the interiors of all simplices containing $\sigma$, where the interior of a simplex $\tau$ is by definition $\tau - \partial \tau$. Thus st $\sigma$ is an open set in $X$ whose closure is $\text{St}\,\sigma$.

\begin{lemma}[Intersections of star neighborhoods]
\label{lem:star-intersections}
\dcref{ch2:2C.2}
\group{Singular Homology}
\uses{def:simplex}
For vertices $v_1, \ldots, v_n$ of a simplicial complex $X$, the intersection st $v_1 \cap \cdots \cap$ st $v_n$ is empty unless $v_1, \ldots, v_n$ are the vertices of a simplex $\sigma$ of $X$, in which case st $v_1 \cap \cdots \cap$ st $v_n = $ st $\sigma$.
\end{lemma}
\begin{proof}
 The intersection st $v_1 \cap \cdots \cap$ st $v_n$ consists of the interiors of all simplices $\tau$ whose vertex set contains $\{v_1, \ldots, v_n\}$. If st $v_1 \cap \cdots \cap$ st $v_n$ is nonempty, such a $\tau$ exists and contains the simplex $\sigma = [v_1, \cdots, v_n] \subset X$. The simplices $\tau$ containing $\{v_1, \ldots, v_n\}$ are just the simplices containing $\sigma$, so st $v_1 \cap \cdots \cap$ st $v_n = $ st $\sigma$.
\end{proof}
\begin{proof}[2C.1]
 Choose a metric on $K$ that restricts to the standard Euclidean metric on each simplex of $K$. For example, $K$ can be viewed as a subcomplex of a simplex $\Delta^N$ whose vertices are all the vertices of $K$, and we can restrict a standard metric on $\Delta^N$ to give a metric on $K$. Let $\varepsilon$ be a Lebesgue number for the open cover $\{f^{-1}(\text{st}\,w) \mid w \text{ is a vertex of } L\}$ of $K$. After iterated barycentric subdivision of $K$ we may assume that each simplex has diameter less than $\varepsilon/2$. The closed star of each vertex $v$ of $K$ then has diameter less than $\varepsilon$, hence this closed star maps by $f$ to the open star of some vertex $g(v)$ of $L$. The resulting map $g: K^0 \to L^0$ thus satisfies $f(\text{St}\,v) \subset \text{st}\,g(v)$ for all vertices $v$ of $K$.

To see that $g$ extends to a simplicial map $g: K \to L$, consider the problem of extending $g$ over a simplex $[v_1, \cdots, v_n]$ of $K$. An interior point $x$ of this simplex lies in st $v_i$ for each $i$, so $f(x)$ lies in st $g(v_i)$ for each $i$, since $f(\text{St}\,v_i) \subset \text{st}\,g(v_i)$ by the definition of $g(v_i)$. Thus st $g(v_1) \cap \cdots \cap$ st $g(v_n) \neq \emptyset$, so $[g(v_1), \cdots, g(v_n)]$ is a simplex of $L$ by the lemma, and we can extend $g$ linearly over $[v_1, \cdots, v_n]$. Both $f(x)$ and $g(x)$ lie in a single simplex of $L$ since $g(x)$ lies in $[g(v_1), \cdots, g(v_n)]$ and $f(x)$ lies in the star of this simplex. So taking the linear path $(1-t)f(x) + tg(x)$, $0 \le t \le 1$, in the simplex containing $f(x)$ and $g(x)$ defines a homotopy from $f$ to $g$. To check continuity of this homotopy it suffices to restrict to the simplex $[v_1, \cdots, v_n]$, where

%% --- END page-0187.tex ---
%% --- BEGIN page-0188.tex ---
continuity is clear since $f(x)$ varies continuously in the star of $[g(v_1), \ldots, g(v_n)]$ and $g(x)$ varies continuously in $[g(v_1), \ldots, g(v_n)]$.
\end{proof}
Notice that if $f$ already sends some vertices of $K$ to vertices of $L$ then we may choose $g$ to equal $f$ on these vertices, and hence the homotopy from $f$ to $g$ will be stationary on these vertices. This is convenient if one is in a situation where one wants maps and homotopies to preserve basepoints.

The proof makes it clear that the simplicial approximation $g$ can be chosen not just homotopic to $f$ but also close to $f$ if we allow subdivisions of $L$ as well as $K$.

\textbf{The Lefschetz Fixed Point Theorem}

This very classical application of homology is a considerable generalization of the Brouwer fixed point theorem. It is also related to the Euler characteristic formula.

For a homomorphism $\varphi: \mathbb{Z}^n \to \mathbb{Z}^n$ with matrix $[a_{ij}]$, the trace tr $\varphi$ is defined to be $\sum_i a_{ii}$, the sum of the diagonal elements of $[a_{ij}]$. Since tr$([a_{ij}][b_{ij}]) = $ tr$([b_{ij}][a_{ij}])$, conjugate matrices have the same trace, and it follows that tr $\varphi$ is independent of the choice of basis for $\mathbb{Z}^n$. For a homomorphism $\varphi: A \to A$ of a finitely generated abelian group $A$ we can then define tr $\varphi$ to be the trace of the induced homomorphism $\bar{\varphi}: A/\text{Torsion} \to A/\text{Torsion}$.

For a map $f: X \to X$ of a finite CW complex $X$, or more generally any space whose homology groups are finitely generated and vanish in high dimensions, the \textbf{Lefschetz number} $\tau(f)$ is defined to be $\sum_n (-1)^n \text{tr}(f_*: H_n(X) \to H_n(X))$. In particular, if $f$ is the identity, or is homotopic to the identity, then $\tau(f)$ is the Euler characteristic $\chi(X)$ since the trace of the $n \times n$ identity matrix is $n$.

Here is the Lefschetz fixed point theorem:

\begin{theorem}[Lefschetz Fixed Point Theorem]
\label{thm:lefschetz-fixed-point}
\dcref{ch2:2C.3}
\group{Singular Homology}
\uses{thm:cellular-singular-homology-isomorphic}
If $X$ is a finite simplicial complex, or more generally a retract of a finite simplicial complex, and $f: X \to X$ is a map with $\tau(f) \neq 0$, then $f$ has a fixed point.
\end{theorem}

As we show in Theorem A.7 in the Appendix, every compact, locally contractible space that can be embedded in $\mathbb{R}^n$ for some $n$ is a retract of a finite simplicial complex. This includes compact manifolds and finite CW complexes, for example. The compactness hypothesis is essential, since a translation of $\mathbb{R}$ has $\tau = 1$ but no fixed points. For an example showing that local properties are also significant, let $X$ be the compact subspace of $\mathbb{R}^2$ consisting of two concentric circles together with a copy of $\mathbb{R}$ between them whose two ends spiral in to the two circles, wrapping around them infinitely often, and let $f: X \to X$ be a homeomorphism translating the copy of $\mathbb{R}$ along itself and rotating the circles, with no fixed points. Since $f$ is homotopic to the identity, we have $\tau(f) = \chi(X)$, which equals 1 since the three path components of $X$ are two circles and a line.



%% --- END page-0188.tex ---
%% --- BEGIN page-0189.tex ---
If $X$ has the same homology groups as a point, at least modulo torsion, then
the theorem says that every map $X \to X$ has a fixed point. This holds for exam-
ple for $\mathbb{R}P^n$ if $n$ is even. The case of projective spaces is interesting because of
its connection with linear algebra. An invertible linear transformation $f: \mathbb{R}^n \to \mathbb{R}^n$
takes lines through $0$ to lines through $0$, hence induces a map $\bar{f}: \mathbb{R}P^{n-1} \to \mathbb{R}P^{n-1}$.
Fixed points of $\bar{f}$ are equivalent to eigenvectors of $f$. The characteristic polyno-
mial of $f$ has odd degree if $n$ is odd, hence has a real root, so an eigenvector ex-
ists in this case. This is in agreement with the observation above that every map
$\mathbb{R}P^{2k} \to \mathbb{R}P^{2k}$ has a fixed point. On the other hand the rotation of $\mathbb{R}^{2k}$ defined by
$f(x_1, \cdots, x_{2k}) = (x_2, -x_1, x_4, -x_3, \cdots, x_{2k}, -x_{2k-1})$ has no eigenvectors and its pro-
jectivization $\bar{f}: \mathbb{R}P^{2k-1} \to \mathbb{R}P^{2k-1}$ has no fixed points.

Similarly, in the complex case an invertible linear transformation $f: \mathbb{C}^n \to \mathbb{C}^n$ in-
duces $\bar{f}: \mathbb{C}P^{n-1} \to \mathbb{C}P^{n-1}$, and this always has a fixed point since the characteristic
polynomial always has a complex root. Nevertheless, as in the real case there is a
map $\mathbb{C}P^{2k-1} \to \mathbb{C}P^{2k-1}$ without fixed points. Namely, consider $f: \mathbb{C}^{2k} \to \mathbb{C}^{2k}$ defined
by $f(z_1, \cdots, z_{2k}) = (\bar{z}_2, -z_1, \bar{z}_4, -z_3, \cdots, \bar{z}_{2k}, -\bar{z}_{2k-1})$. This map is only `conjugate-
linear' over $\mathbb{C}$, but this is still good enough to imply that $f$ induces a well-defined
map $\bar{f}$ on $\mathbb{C}P^{2k-1}$, and it is easy to check that $\bar{f}$ has no fixed points. The similarity
between the real and complex cases persists in the fact that every map $\mathbb{C}P^{2k} \to \mathbb{C}P^{2k}$
has a fixed point, though to deduce this from the Lefschetz fixed point theorem re-
quires more structure than homology has, so this will be left as an exercise for \S 3.2,
using cup products in cohomology.

One could go further and consider the quaternionic case. The antipodal map of
$S^4 = \mathbb{H}P^1$ has no fixed points, but every map $\mathbb{H}P^n \to \mathbb{H}P^n$ with $n > 1$ does have a
fixed point. This is shown in Example 4L.4 using considerably heavier machinery.
\begin{proof}[2C.3]
 The general case easily reduces to the case of finite simplicial com-
plexes, for suppose $r: K \to X$ is a retraction of a finite simplicial complex $K$ onto $X$.
For a map $f: X \to X$, the composition $fr: K \to X \subset K$ then has exactly the same fixed
points as $f$. Since $r_*: H_n(K) \to H_n(X)$ is projection onto a direct summand, we have
$\text{tr}(f_* r_*) = \text{tr}(f_*)$ and hence $\tau(fr) = \tau(f)$.

For $X$ a finite simplicial complex, suppose that $f: X \to X$ has no fixed points. We
claim there is a subdivision $L$ of $X$, a further subdivision $K$ of $L$, and a simplicial map
$g: K \to L$ homotopic to $f$ such that $g(\sigma) \cap \sigma = \emptyset$ for each simplex $\sigma$ of $K$. To see this,
first choose a metric $d$ on $X$ as in the proof of the simplicial approximation theorem.
Since $f$ has no fixed points, $d(x, f(x)) > 0$ for all $x \in X$, so by the compactness
of $X$ there is an $\varepsilon > 0$ such that $d(x, f(x)) > \varepsilon$ for all $x$. Choose a subdivision
$L$ of $X$ so that the stars of all simplices have diameter less than $\varepsilon/2$. Applying the
simplicial approximation theorem, there is a subdivision $K$ of $L$ and a simplicial map
$g: K \to L$ homotopic to $f$. By construction, $g$ has the property that for each simplex
$\sigma$ of $K$, $f(\sigma)$ is contained in the star of the simplex $g(\sigma)$. Then $g(\sigma) \cap \sigma = \emptyset$

%% --- END page-0189.tex ---
%% --- BEGIN page-0190.tex ---
for each simplex $\sigma$ of $K$ since for any choice of $x \in \sigma$ we have $d(x, f(x)) > \varepsilon$,
while $g(\sigma)$ lies within distance $\varepsilon/2$ of $f(x)$ and $\sigma$ lies within distance $\varepsilon/2$ of $x$, as
a consequence of the fact that $\sigma$ is contained in a simplex of $L$, $K$ being a subdivision
of $L$.

The Lefschetz numbers $\tau(f)$ and $\tau(g)$ are equal since $f$ and $g$ are homotopic.
Since $g$ is simplicial, it takes the $n$-skeleton $K^n$ of $K$ to the $n$-skeleton $L^n$ of $L$, for
each $n$. Since $K$ is a subdivision of $L$, $L^n$ is contained in $K^n$, and hence $g(K^n) \subset K^n$
for all $n$. Thus $g$ induces a chain map of the cellular chain complex $\{H_n(K^n, K^{n-1})\}$
to itself. This can be used to compute $\tau(g)$ according to the formula

\[
\tau(g) = \sum_n (-1)^n \text{tr}(g_* : H_n(K^n, K^{n-1}) \to H_n(K^n, K^{n-1}))
\]

This is the analog of Theorem 2.44 for trace instead of rank, and is proved in precisely
the same way, based on the elementary algebraic fact that trace is additive for endomorphisms of short exact sequences: Given a commutative diagram as at the right with exact rows,



then $\text{tr } \beta = \text{tr } \alpha + \text{tr } \gamma$. This algebraic fact can be
proved by reducing to the easy case that $A$, $B$, and
$C$ are free by first factoring out the torsion in $B$, hence also the torsion in $A$, then
eliminating any remaining torsion in $C$ by replacing $A$ by a larger subgroup $A' \subset B$,
with $A$ having finite index in $A'$. The details of this argument are left to the reader.

Finally, note that $g_* : H_n(K^n, K^{n-1}) \to H_n(K^n, K^{n-1})$ has trace 0 since the matrix
for $g_*$ has zeros down the diagonal, in view of the fact that $g(\sigma) \cap \sigma = \emptyset$ for each
$n$-simplex $\sigma$. So $\tau(f) = \tau(g) = 0$.
\end{proof}
\begin{example}[Lefschetz number of surface rotation]
\label{ex:lefschetz-surface-rotation}
\dcref{ch2:2C.4}
\group{Singular Homology}
\uses{thm:lefschetz-fixed-point}
Let us verify the theorem in an example. Let $X$ be the closed orientable surface of genus 3 as shown in the figure below, with $f:X \to X$ the 180
degree rotation about a vertical axis passing through the central hole of
$X$. Since $f$ has no fixed points, we
should have $\tau(f) = 0$. The induced
map $f_* : H_0(X) \to H_0(X)$ is the identity, as always for a path-connected space, so this contributes 1 to $\tau(f)$. For $H_1(X)$
we saw in Example 2A.2 that the six loops $\alpha_i$ and $\beta_i$ represent a basis. The map
$f_*$ interchanges the homology classes of $\alpha_1$ and $\alpha_3$, and likewise for $\beta_1$ and $\beta_3$,
while $\beta_2$ is sent to itself and $\alpha_2$ is sent to $\alpha_2'$ which is homologous to $\alpha_2$ as we
saw in Example 2A.2. So $f_* : H_1(X) \to H_1(X)$ contributes $-2$ to $\tau(f)$. It remains
to check that $f_* : H_2(X) \to H_2(X)$ is the identity, which we do by the commutative diagram
at the right, where $x$ is a point of $X$ in the central torus and $y = f(x)$. We can see that the
\end{example}





%% --- END page-0190.tex ---
%% --- BEGIN page-0191.tex ---
left-hand vertical map is an isomorphism by considering the long exact sequence of
the triple $(X, X - \{x\}, X^1)$ where $X^1$ is the 1-skeleton of $X$ in its usual CW struc-
ture and $x$ is chosen in $X - X^1$, so that $X - \{x\}$ deformation retracts onto $X^1$ and
$H_n(X - \{x\}, X^1) = 0$ for all $n$. The same reasoning shows the right-hand vertical
map is an isomorphism. There is a similar commutative diagram with $f$ replaced by
a homeomorphism $g$ that is homotopic to the identity and equals $f$ in a neighbor-
hood of $x$, with $g$ the identity outside a disk in $X$ containing $x$ and $y$. Since $g$ is
homotopic to the identity, it induces the identity across the top row of the diagram,
and since $g$ equals $f$ near $x$, it induces the same map as $f$ in the bottom row of the
diagram, by excision. It follows that the map $f_*$ in the upper row is the identity.

This example generalizes to surfaces of any odd genus by adding symmetric pairs
of tori at the left and right. Examples for even genus are described in one of the
exercises.

Fixed point theory is a well-developed side branch of algebraic topology, but we
touch upon it only occasionally in this book. For a nice introduction see [Brown 1971].

\section{Simplicial Approximations to CW Complexes}

The simplicial approximation theorem allows arbitrary continuous maps to be
replaced by homotopic simplicial maps in many situations, and one might wonder
about the analogous question for spaces: Which spaces are homotopy equivalent to
simplicial complexes? We will show this is true for the most common class of spaces
in algebraic topology, CW complexes. In the Appendix the question is answered for a
few other classes of spaces as well.

\begin{theorem}[CW complex homotopy equivalent to simplicial complex]
\label{thm:cw-homotopy-equivalent-simplicial}
\dcref{ch2:2C.5}
\group{CW Complexes}
\uses{def:cell-complex}
Every CW complex $X$ is homotopy equivalent to a simplicial complex,
which can be chosen to be of the same dimension as $X$, finite if $X$ is finite, and
countable if $X$ is countable.
\end{theorem}

We will build a simplicial complex $Y \simeq X$ inductively as an increasing union of
subcomplexes $Y_n$ homotopy equivalent to the skeleta $X^n$. For the inductive step,
assuming we have already constructed $Y_n \simeq X^n$, let $e^{n+1}$ be an $(n+1)$-cell of $X$
attached by a map $\varphi : S^n \to X^n$. The map $S^n \to Y_n$ corresponding to $\varphi$ under the
homotopy equivalence $Y_n \simeq X^n$ is homotopic to a simplicial map $f : S^n \to Y_n$ by the
simplicial approximation theorem, and it is not hard to see that the spaces $X^n \cup_\varphi e^{n+1}$
and $Y_n \cup_f e^{n+1}$ are homotopy equivalent, where the subscripts denote attaching $e^{n+1}$
via $\varphi$ and $f$, respectively; see Proposition 0.18 for a proof. We can view $Y_n \cup_f e^{n+1}$
as the mapping cone $C_f$, obtained from the mapping cylinder of $f$ by collapsing the
domain end to a point. If we knew that the mapping cone of a simplicial map was a
simplicial complex, then by performing the same construction for all the $(n+1)$-cells
of $X$ we would have completed the induction step. Unfortunately, and somewhat
surprisingly, mapping cones and mapping cylinders are rather awkward objects in the

%% --- END page-0191.tex ---
%% --- BEGIN page-0192.tex ---
simplicial category. To avoid this awkwardness we will instead construct simplicial analogs of mapping cones and cylinders that have all the essential features of actual mapping cones and cylinders.

Let us first construct the simplicial analog of a mapping cylinder. For a simplicial map $f : K \to L$ this will be a simplicial complex $M(f)$ containing both $L$ and the barycentric subdivision $K'$ of $K$ as subcomplexes, and such that there is a deformation retraction $r_1$ of $M(f)$ onto $L$ with $r_1|_{K'} = f$. The figure shows the case that $f$ is a simplicial surjection $\Delta^2 \to \Delta^1$. The construction proceeds one simplex of $K$ at a time, by induction on dimension. To begin, the ordinary mapping cylinder of $f : K^0 \to L$ suffices for $M(f|K^0)$. Assume inductively that we have already constructed $M(f|K^{n-1})$. Let $\sigma$ be an $n$-simplex of $K$ and let $\tau = f(\sigma)$, a simplex of $L$ of dimension $n$ or less. By the inductive hypothesis we have already constructed $M(f : \partial \sigma \to \tau)$ with the desired properties, and we let $M(f : \sigma \to \tau)$ be the cone on $M(f : \partial \sigma \to \tau)$, as shown in the figure. The space $M(f : \partial \sigma \to \tau)$ is contractible since by induction it deformation retracts onto $\tau$ which is contractible. The cone $M(f : \sigma \to \tau)$ is of course contractible, so the inclusion of $M(f : \partial \sigma \to \tau)$ into $M(f : \sigma \to \tau)$ is a homotopy equivalence. This implies that $M(f : \sigma \to \tau)$ deformation retracts onto $M(f : \partial \sigma \to \tau)$ by Corollary 2.20, or one can give a direct argument using the fact that $M(f : \partial \sigma \to \tau)$ is contractible. By attaching $M(f : \sigma \to \tau)$ to $M(f|K^{n-1})$ along $M(f : \partial \sigma \to \tau) \subset M(f|K^{n-1})$ for all $n$-simplices $\sigma$ of $K$ we obtain $M(f|K^n)$ with a deformation retraction onto $M(f|K^{n-1})$. Taking the union over all $n$ yields $M(f)$ with a deformation retraction $r_1$ onto $L$, the infinite concatenation of the previous deformation retractions, with the deformation retraction of $M(f|K^n)$ onto $M(f|K^{n-1})$ performed in the $t$-interval $[1/2^{n+1}, 1/2^n]$. The map $r_1|K$ may not equal $f$, but it is homotopic to $f$ via the linear homotopy $tf + (1-t)r_1$, which is defined since $r_1(\sigma) \subset f(\sigma)$ for all simplices $\sigma$ of $K$. By applying the homotopy extension property to the homotopy of $r_1$ that equals $tf + (1-t)r_1$ on $K$ and the identity map on $L$, we can improve our deformation retraction of $M(f)$ onto $L$ so that its restriction to $K$ at time 1 is $f$.



From the simplicial analog $M(f)$ of a mapping cylinder we construct the simplicial `mapping cone' $C(f)$ by attaching the ordinary cone on $K'$ to the subcomplex $K' \subset M(f)$.
\begin{proof}[2C.5]
 We will construct for each $n$ a CW complex $Z_n$ containing $X^n$ as a deformation retract and also containing as a deformation retract a subcomplex $Y_n$ that is a simplicial complex. Beginning with $Y_0 = Z_0 = X^0$, suppose inductively that we have already constructed $Y_n$ and $Z_n$. Let the cells $e^{n+1}_\alpha$ of $X$ be attached by maps $\varphi_\alpha : S^n \to X^n$. Using the simplicial approximation theorem, there is a homotopy from $\varphi_\alpha$ to a simplicial map $f_\alpha : S^n \to Y_n$. The CW complex $W_n = Z_n \cup_\alpha M(f_\alpha)$ contains a

%% --- END page-0192.tex ---
%% --- BEGIN page-0193.tex ---
simplicial subcomplex $S_\alpha^n$ homeomorphic to $S^n$ at one end of $M(f_\alpha)$, and the homeomorphism $S^n \approx S_\alpha^n$ is homotopic in $W_n$ to the map $f_n\alpha$, hence also to $\varphi_\alpha$. Let $Z_{n+1}$ be obtained from $Z_n$ by attaching $D_\alpha^{n+1} \times I$'s via these homotopies between the $\varphi_\alpha$'s and the inclusions $S_\alpha^n \hookrightarrow W_n$. Thus $Z_{n+1}$ contains $X^{n+1}$ at one end, and at the other end we have a simplicial complex $Y_{n+1} = Y_n \cup_\alpha C(f_\alpha)$, where $C(f_\alpha)$ is obtained from $M(f_\alpha)$ by attaching a cone on the subcomplex $S_\alpha^n$. Since $D^{n+1} \times I$ deformation retracts onto $\partial D^{n+1} \times I \cup D^{n+1} \times \{1\}$, we see that $Z_{n+1}$ deformation retracts onto $Z_n \cup Y_{n+1}$, which in turn deformation retracts onto $Y_n \cup Y_{n+1} = Y_{n+1}$ by induction. Likewise, $Z_{n+1}$ deformation retracts onto $X^{n+1} \cup W_n$ which deformation retracts onto $X^{n+1} \cup Z_n$ and hence onto $X^{n+1} \cup X^n = X^{n+1}$ by induction.

Let $Y = \bigcup_n Y_n$ and $Z = \bigcup_n Z_n$. The deformation retractions of $Z_n$ onto $X^n$ give deformation retractions of $X \cup Z_n$ onto $X$, and the infinite concatenation of the latter deformation retractions is a deformation retraction of $Z$ onto $X$. Similarly, $Z$ deformation retracts onto $Y$.
\end{proof}
\textbf{Exercises}

%ORIGINAL-NUMBER: 1
\begin{exercise}
What is the minimum number of edges in simplicial complex structures $K$ and $L$ on $S^1$ such that there is a simplicial map $K \to L$ of degree $n$?
\end{exercise}

%ORIGINAL-NUMBER: 2
\begin{exercise}
Use the Lefschetz fixed point theorem to show that a map $S^n \to S^n$ has a fixed point unless its degree is equal to the degree of the antipodal map $x \mapsto -x$.
\end{exercise}

%ORIGINAL-NUMBER: 3
\begin{exercise}
Verify that the formula $f(z_1, \ldots, z_{2k}) = (\overline{z_2}, -\overline{z_1}, \overline{z_4}, -\overline{z_3}, \ldots, \overline{z_{2k}}, -\overline{z_{2k-1}})$ defines a map $f: \mathbb{C}^k \to \mathbb{C}^k$ inducing a quotient map $\mathbb{CP}^{2k-1} \to \mathbb{CP}^{2k-1}$ without fixed points.
\end{exercise}

%ORIGINAL-NUMBER: 4
\begin{exercise}
If $X$ is a finite simplicial complex and $f: X \to X$ is a simplicial homeomorphism, show that the Lefschetz number $\tau(f)$ equals the Euler characteristic of the set of fixed points of $f$. In particular, $\tau(f)$ is the number of fixed points if the fixed points are isolated. [Hint: Barycentrically subdivide $X$ to make the fixed point set a subcomplex.]
\end{exercise}

%ORIGINAL-NUMBER: 5
\begin{exercise}
Let $M$ be a closed orientable surface embedded in $\mathbb{R}^3$ in such a way that reflection across a plane $P$ defines a homeomorphism $r: M \to M$ fixing $M \cap P$, a collection of circles. Is it possible to homotope $r$ to have no fixed points?
\end{exercise}

%ORIGINAL-NUMBER: 6
\begin{exercise}
Do an even-genus analog of Example 2C.4 by replacing the central torus by a sphere, letting $f$ be a homeomorphism that restricts to the antipodal map on this sphere.
\end{exercise}

%ORIGINAL-NUMBER: 7
\begin{exercise}
Verify that the Lefschetz fixed point theorem holds also when $\tau(f)$ is defined using homology with coefficients in a field $F$.
\end{exercise}

%ORIGINAL-NUMBER: 8
\begin{exercise}
Let $X$ be homotopy equivalent to a finite simplicial complex and let $Y$ be homotopy equivalent to a finite or countably infinite simplicial complex. Using the simplicial approximation theorem, show that there are at most countably many homotopy classes of maps $X \to Y$.
\end{exercise}

%ORIGINAL-NUMBER: 9
\begin{exercise}
Show that there are only countably many homotopy types of finite CW complexes.
\end{exercise}

%% --- END page-0193.tex ---
